% Options for packages loaded elsewhere
\PassOptionsToPackage{unicode}{hyperref}
\PassOptionsToPackage{hyphens}{url}
\PassOptionsToPackage{dvipsnames,svgnames,x11names}{xcolor}
%
\documentclass[
  letterpaper,
  DIV=11,
  numbers=noendperiod]{scrreprt}

\usepackage{amsmath,amssymb}
\usepackage{iftex}
\ifPDFTeX
  \usepackage[T1]{fontenc}
  \usepackage[utf8]{inputenc}
  \usepackage{textcomp} % provide euro and other symbols
\else % if luatex or xetex
  \usepackage{unicode-math}
  \defaultfontfeatures{Scale=MatchLowercase}
  \defaultfontfeatures[\rmfamily]{Ligatures=TeX,Scale=1}
\fi
\usepackage{lmodern}
\ifPDFTeX\else  
    % xetex/luatex font selection
\fi
% Use upquote if available, for straight quotes in verbatim environments
\IfFileExists{upquote.sty}{\usepackage{upquote}}{}
\IfFileExists{microtype.sty}{% use microtype if available
  \usepackage[]{microtype}
  \UseMicrotypeSet[protrusion]{basicmath} % disable protrusion for tt fonts
}{}
\makeatletter
\@ifundefined{KOMAClassName}{% if non-KOMA class
  \IfFileExists{parskip.sty}{%
    \usepackage{parskip}
  }{% else
    \setlength{\parindent}{0pt}
    \setlength{\parskip}{6pt plus 2pt minus 1pt}}
}{% if KOMA class
  \KOMAoptions{parskip=half}}
\makeatother
\usepackage{xcolor}
\setlength{\emergencystretch}{3em} % prevent overfull lines
\setcounter{secnumdepth}{5}
% Make \paragraph and \subparagraph free-standing
\makeatletter
\ifx\paragraph\undefined\else
  \let\oldparagraph\paragraph
  \renewcommand{\paragraph}{
    \@ifstar
      \xxxParagraphStar
      \xxxParagraphNoStar
  }
  \newcommand{\xxxParagraphStar}[1]{\oldparagraph*{#1}\mbox{}}
  \newcommand{\xxxParagraphNoStar}[1]{\oldparagraph{#1}\mbox{}}
\fi
\ifx\subparagraph\undefined\else
  \let\oldsubparagraph\subparagraph
  \renewcommand{\subparagraph}{
    \@ifstar
      \xxxSubParagraphStar
      \xxxSubParagraphNoStar
  }
  \newcommand{\xxxSubParagraphStar}[1]{\oldsubparagraph*{#1}\mbox{}}
  \newcommand{\xxxSubParagraphNoStar}[1]{\oldsubparagraph{#1}\mbox{}}
\fi
\makeatother

\usepackage{color}
\usepackage{fancyvrb}
\newcommand{\VerbBar}{|}
\newcommand{\VERB}{\Verb[commandchars=\\\{\}]}
\DefineVerbatimEnvironment{Highlighting}{Verbatim}{commandchars=\\\{\}}
% Add ',fontsize=\small' for more characters per line
\usepackage{framed}
\definecolor{shadecolor}{RGB}{241,243,245}
\newenvironment{Shaded}{\begin{snugshade}}{\end{snugshade}}
\newcommand{\AlertTok}[1]{\textcolor[rgb]{0.68,0.00,0.00}{#1}}
\newcommand{\AnnotationTok}[1]{\textcolor[rgb]{0.37,0.37,0.37}{#1}}
\newcommand{\AttributeTok}[1]{\textcolor[rgb]{0.40,0.45,0.13}{#1}}
\newcommand{\BaseNTok}[1]{\textcolor[rgb]{0.68,0.00,0.00}{#1}}
\newcommand{\BuiltInTok}[1]{\textcolor[rgb]{0.00,0.23,0.31}{#1}}
\newcommand{\CharTok}[1]{\textcolor[rgb]{0.13,0.47,0.30}{#1}}
\newcommand{\CommentTok}[1]{\textcolor[rgb]{0.37,0.37,0.37}{#1}}
\newcommand{\CommentVarTok}[1]{\textcolor[rgb]{0.37,0.37,0.37}{\textit{#1}}}
\newcommand{\ConstantTok}[1]{\textcolor[rgb]{0.56,0.35,0.01}{#1}}
\newcommand{\ControlFlowTok}[1]{\textcolor[rgb]{0.00,0.23,0.31}{\textbf{#1}}}
\newcommand{\DataTypeTok}[1]{\textcolor[rgb]{0.68,0.00,0.00}{#1}}
\newcommand{\DecValTok}[1]{\textcolor[rgb]{0.68,0.00,0.00}{#1}}
\newcommand{\DocumentationTok}[1]{\textcolor[rgb]{0.37,0.37,0.37}{\textit{#1}}}
\newcommand{\ErrorTok}[1]{\textcolor[rgb]{0.68,0.00,0.00}{#1}}
\newcommand{\ExtensionTok}[1]{\textcolor[rgb]{0.00,0.23,0.31}{#1}}
\newcommand{\FloatTok}[1]{\textcolor[rgb]{0.68,0.00,0.00}{#1}}
\newcommand{\FunctionTok}[1]{\textcolor[rgb]{0.28,0.35,0.67}{#1}}
\newcommand{\ImportTok}[1]{\textcolor[rgb]{0.00,0.46,0.62}{#1}}
\newcommand{\InformationTok}[1]{\textcolor[rgb]{0.37,0.37,0.37}{#1}}
\newcommand{\KeywordTok}[1]{\textcolor[rgb]{0.00,0.23,0.31}{\textbf{#1}}}
\newcommand{\NormalTok}[1]{\textcolor[rgb]{0.00,0.23,0.31}{#1}}
\newcommand{\OperatorTok}[1]{\textcolor[rgb]{0.37,0.37,0.37}{#1}}
\newcommand{\OtherTok}[1]{\textcolor[rgb]{0.00,0.23,0.31}{#1}}
\newcommand{\PreprocessorTok}[1]{\textcolor[rgb]{0.68,0.00,0.00}{#1}}
\newcommand{\RegionMarkerTok}[1]{\textcolor[rgb]{0.00,0.23,0.31}{#1}}
\newcommand{\SpecialCharTok}[1]{\textcolor[rgb]{0.37,0.37,0.37}{#1}}
\newcommand{\SpecialStringTok}[1]{\textcolor[rgb]{0.13,0.47,0.30}{#1}}
\newcommand{\StringTok}[1]{\textcolor[rgb]{0.13,0.47,0.30}{#1}}
\newcommand{\VariableTok}[1]{\textcolor[rgb]{0.07,0.07,0.07}{#1}}
\newcommand{\VerbatimStringTok}[1]{\textcolor[rgb]{0.13,0.47,0.30}{#1}}
\newcommand{\WarningTok}[1]{\textcolor[rgb]{0.37,0.37,0.37}{\textit{#1}}}

\providecommand{\tightlist}{%
  \setlength{\itemsep}{0pt}\setlength{\parskip}{0pt}}\usepackage{longtable,booktabs,array}
\usepackage{calc} % for calculating minipage widths
% Correct order of tables after \paragraph or \subparagraph
\usepackage{etoolbox}
\makeatletter
\patchcmd\longtable{\par}{\if@noskipsec\mbox{}\fi\par}{}{}
\makeatother
% Allow footnotes in longtable head/foot
\IfFileExists{footnotehyper.sty}{\usepackage{footnotehyper}}{\usepackage{footnote}}
\makesavenoteenv{longtable}
\usepackage{graphicx}
\makeatletter
\newsavebox\pandoc@box
\newcommand*\pandocbounded[1]{% scales image to fit in text height/width
  \sbox\pandoc@box{#1}%
  \Gscale@div\@tempa{\textheight}{\dimexpr\ht\pandoc@box+\dp\pandoc@box\relax}%
  \Gscale@div\@tempb{\linewidth}{\wd\pandoc@box}%
  \ifdim\@tempb\p@<\@tempa\p@\let\@tempa\@tempb\fi% select the smaller of both
  \ifdim\@tempa\p@<\p@\scalebox{\@tempa}{\usebox\pandoc@box}%
  \else\usebox{\pandoc@box}%
  \fi%
}
% Set default figure placement to htbp
\def\fps@figure{htbp}
\makeatother
% definitions for citeproc citations
\NewDocumentCommand\citeproctext{}{}
\NewDocumentCommand\citeproc{mm}{%
  \begingroup\def\citeproctext{#2}\cite{#1}\endgroup}
\makeatletter
 % allow citations to break across lines
 \let\@cite@ofmt\@firstofone
 % avoid brackets around text for \cite:
 \def\@biblabel#1{}
 \def\@cite#1#2{{#1\if@tempswa , #2\fi}}
\makeatother
\newlength{\cslhangindent}
\setlength{\cslhangindent}{1.5em}
\newlength{\csllabelwidth}
\setlength{\csllabelwidth}{3em}
\newenvironment{CSLReferences}[2] % #1 hanging-indent, #2 entry-spacing
 {\begin{list}{}{%
  \setlength{\itemindent}{0pt}
  \setlength{\leftmargin}{0pt}
  \setlength{\parsep}{0pt}
  % turn on hanging indent if param 1 is 1
  \ifodd #1
   \setlength{\leftmargin}{\cslhangindent}
   \setlength{\itemindent}{-1\cslhangindent}
  \fi
  % set entry spacing
  \setlength{\itemsep}{#2\baselineskip}}}
 {\end{list}}
\usepackage{calc}
\newcommand{\CSLBlock}[1]{\hfill\break\parbox[t]{\linewidth}{\strut\ignorespaces#1\strut}}
\newcommand{\CSLLeftMargin}[1]{\parbox[t]{\csllabelwidth}{\strut#1\strut}}
\newcommand{\CSLRightInline}[1]{\parbox[t]{\linewidth - \csllabelwidth}{\strut#1\strut}}
\newcommand{\CSLIndent}[1]{\hspace{\cslhangindent}#1}

\KOMAoption{captions}{tableheading}
\makeatletter
\@ifpackageloaded{tcolorbox}{}{\usepackage[skins,breakable]{tcolorbox}}
\@ifpackageloaded{fontawesome5}{}{\usepackage{fontawesome5}}
\definecolor{quarto-callout-color}{HTML}{909090}
\definecolor{quarto-callout-note-color}{HTML}{0758E5}
\definecolor{quarto-callout-important-color}{HTML}{CC1914}
\definecolor{quarto-callout-warning-color}{HTML}{EB9113}
\definecolor{quarto-callout-tip-color}{HTML}{00A047}
\definecolor{quarto-callout-caution-color}{HTML}{FC5300}
\definecolor{quarto-callout-color-frame}{HTML}{acacac}
\definecolor{quarto-callout-note-color-frame}{HTML}{4582ec}
\definecolor{quarto-callout-important-color-frame}{HTML}{d9534f}
\definecolor{quarto-callout-warning-color-frame}{HTML}{f0ad4e}
\definecolor{quarto-callout-tip-color-frame}{HTML}{02b875}
\definecolor{quarto-callout-caution-color-frame}{HTML}{fd7e14}
\makeatother
\makeatletter
\@ifpackageloaded{bookmark}{}{\usepackage{bookmark}}
\makeatother
\makeatletter
\@ifpackageloaded{caption}{}{\usepackage{caption}}
\AtBeginDocument{%
\ifdefined\contentsname
  \renewcommand*\contentsname{Table of contents}
\else
  \newcommand\contentsname{Table of contents}
\fi
\ifdefined\listfigurename
  \renewcommand*\listfigurename{List of Figures}
\else
  \newcommand\listfigurename{List of Figures}
\fi
\ifdefined\listtablename
  \renewcommand*\listtablename{List of Tables}
\else
  \newcommand\listtablename{List of Tables}
\fi
\ifdefined\figurename
  \renewcommand*\figurename{Figure}
\else
  \newcommand\figurename{Figure}
\fi
\ifdefined\tablename
  \renewcommand*\tablename{Table}
\else
  \newcommand\tablename{Table}
\fi
}
\@ifpackageloaded{float}{}{\usepackage{float}}
\floatstyle{ruled}
\@ifundefined{c@chapter}{\newfloat{codelisting}{h}{lop}}{\newfloat{codelisting}{h}{lop}[chapter]}
\floatname{codelisting}{Listing}
\newcommand*\listoflistings{\listof{codelisting}{List of Listings}}
\makeatother
\makeatletter
\makeatother
\makeatletter
\@ifpackageloaded{caption}{}{\usepackage{caption}}
\@ifpackageloaded{subcaption}{}{\usepackage{subcaption}}
\makeatother

\usepackage{bookmark}

\IfFileExists{xurl.sty}{\usepackage{xurl}}{} % add URL line breaks if available
\urlstyle{same} % disable monospaced font for URLs
\hypersetup{
  pdftitle={Charlie's R for Psychological Science Repository},
  pdfauthor={Charlie Pittaway},
  colorlinks=true,
  linkcolor={blue},
  filecolor={Maroon},
  citecolor={Blue},
  urlcolor={Blue},
  pdfcreator={LaTeX via pandoc}}


\title{Charlie's R for Psychological Science Repository}
\author{Charlie Pittaway}
\date{20 July 2026}

\begin{document}
\maketitle

\renewcommand*\contentsname{Table of contents}
{
\hypersetup{linkcolor=}
\setcounter{tocdepth}{2}
\tableofcontents
}

\bookmarksetup{startatroot}

\chapter*{Preface}\label{preface}
\addcontentsline{toc}{chapter}{Preface}

\markboth{Preface}{Preface}

This is an ongoing collection of functions/workflows/code snippets to
run analyses in R, with a psychological science twist.

I learned how to do many of these things from other people's generously
developed guides, which I have cited and linked throughout. Some
particularly helpful guides are \href{guides.qmd}{linked here}.

The initial goal of this project was to provide a
\href{data-cleaning-overview.qmd}{workflow for data cleaning} for
psychology researchers, as I was unable to find something similar when I
was a PhD student. It has expanded to be a living repository of all of
the R code I have collected over the years, with working syntax taken
directly from my own projects.

If you find new packages or better processes, please let me know:

\begin{itemize}
\tightlist
\item
  📧
  \href{mailto:charlie.pittaway@flinders.edu.au}{\nolinkurl{charlie.pittaway@flinders.edu.au}}
\end{itemize}

\part{Data Wrangling \& Basics}

\chapter{Getting Started}\label{getting-started}

So, you've finished collecting data in Qualtrics. Now what?

\section{Setting up your R project}\label{setting-up-your-r-project}

Before you import anything, it's worth spending five minutes setting up
a proper R Project. This keeps your files organised and makes it much
easier to upload a clean, self-contained folder to OSF later, since
everything the project needs lives in one place with relative (not
absolute) file paths.\footnote{An \emph{absolute path} is the full file
  location starting from the root of your computer's file system,
  e.g.~\texttt{"C:/Users/yourname/Documents/study1\_analysis/data/my\_data.sav"}
  on Windows. This works fine on your machine, but breaks the moment the
  folder is moved, renamed, or opened by someone else. A \emph{relative
  path}, like \texttt{"data/my\_data.sav"}, is defined relative to
  wherever the \texttt{.Rproj} file lives, so it keeps working no matter
  who opens the folder or where it's saved.}

\begin{enumerate}
\def\labelenumi{\arabic{enumi}.}
\tightlist
\item
  In RStudio: \textbf{File → New Project → New Directory → New Project}.
\item
  Give it a sensible name that represents the study or package of
  studies
  (e.g.~\texttt{Longitudinal\ analysis\ of\ temporal\ orientation}) and
  choose where to save it.
\item
  Inside that project folder, create a few subfolders to keep things
  tidy. My standard folders are: \textbf{data} (where the raw and
  cleaned data files live), \textbf{scripts} (where data cleaning and
  analysis code lives), and \textbf{outputs} (figures, tables, and
  exported results).
\end{enumerate}

\begin{tcolorbox}[enhanced jigsaw, toprule=.15mm, coltitle=black, colframe=quarto-callout-important-color-frame, colbacktitle=quarto-callout-important-color!10!white, opacityback=0, left=2mm, breakable, rightrule=.15mm, bottomrule=.15mm, leftrule=.75mm, titlerule=0mm, arc=.35mm, opacitybacktitle=0.6, toptitle=1mm, bottomtitle=1mm, title=\textcolor{quarto-callout-important-color}{\faExclamation}\hspace{0.5em}{Important}, colback=white]

Save your raw Qualtrics export into \texttt{data/} as soon as you
download it, and don't edit it directly --- treat it as read-only. Any
cleaning should happen in a script, not in the file itself, so you (or
anyone else) can always trace exactly what was done to the data and
reproduce it.

\end{tcolorbox}

\section{Installing R packages}\label{installing-r-packages}

R has a completely bonkers amount of versatility. There is a custom
built package for most things. To install a package, select
\textbf{Tools} from the toolbar along the top of the window, then
``Install packages'' and then search for the package you are looking
for. You can also install using code, as below. Once you have installed
a package, you must load it before using it.

\begin{Shaded}
\begin{Highlighting}[]
\FunctionTok{install.packages}\NormalTok{(}\StringTok{"package\_name"}\NormalTok{) }\CommentTok{\#Load package}
\FunctionTok{library}\NormalTok{(package\_name) }\CommentTok{\# Install package}
\end{Highlighting}
\end{Shaded}

\section{Importing data into R}\label{importing-data-into-r}

You're probably familiar with downloading your data from Qualtrics in
\texttt{.sav} format to import into SPSS, or \texttt{.csv} format to
open in Excel. Both of these file types
(\href{https://rfaqs.com/r-language-basics/files-in-r/}{and many more})
can be used in R.

My recommendation is to save as an \texttt{.sav} file for the utility of
having the item wording retained, which you lose in \texttt{.csv}
format. To do this, you'll need to install the \textbf{haven} package
(Wickham et al., 2025).

\begin{Shaded}
\begin{Highlighting}[]
\FunctionTok{install.packages}\NormalTok{(}\StringTok{"haven"}\NormalTok{)}
\FunctionTok{library}\NormalTok{(haven)}

\CommentTok{\# Import your data into R, using the relative file path}
\NormalTok{dat }\OtherTok{\textless{}{-}} \FunctionTok{read\_sav}\NormalTok{(}\StringTok{"data/my\_qualtrics\_export.sav"}\NormalTok{)}
\end{Highlighting}
\end{Shaded}

If you don't care about retaining the item stem and \texttt{.csv} is
sufficient, you'll need the \textbf{readr} package instead (Wickham et
al., 2026).

\begin{Shaded}
\begin{Highlighting}[]
\FunctionTok{install.packages}\NormalTok{(}\StringTok{"readr"}\NormalTok{)}
\FunctionTok{library}\NormalTok{(readr)}

\CommentTok{\# Import your data into R, using the relative file path}
\NormalTok{dat }\OtherTok{\textless{}{-}} \FunctionTok{read\_csv}\NormalTok{(}\StringTok{"data/my\_qualtrics\_export.csv"}\NormalTok{)}
\end{Highlighting}
\end{Shaded}

\begin{tcolorbox}[enhanced jigsaw, toprule=.15mm, coltitle=black, colframe=quarto-callout-warning-color-frame, colbacktitle=quarto-callout-warning-color!10!white, opacityback=0, left=2mm, breakable, rightrule=.15mm, bottomrule=.15mm, leftrule=.75mm, titlerule=0mm, arc=.35mm, opacitybacktitle=0.6, toptitle=1mm, bottomtitle=1mm, title=\textcolor{quarto-callout-warning-color}{\faExclamationTriangle}\hspace{0.5em}{Common error}, colback=white]

If \texttt{read\_sav()} throws an error about the file not being found,
double check you're working inside your \texttt{.Rproj} (look for the
project name in the top-right corner of RStudio) and that the file path
matches your actual folder structure, e.g.~\texttt{"data/filename.sav"}
not just \texttt{"filename.sav"}.

\end{tcolorbox}

Now that you have your data saved and your R project structure set up,
you are ready to move on to \href{data-cleaning.qmd}{data
cleaning}\ldots{}

\section{References}\label{references}

\chapter{Selecting/Deselecting
Variables}\label{selectingdeselecting-variables}

\section{Loading full dataset from
file}\label{loading-full-dataset-from-file}

Using \textbf{haven} for .sav files (Wickham et al., 2025):

\begin{Shaded}
\begin{Highlighting}[]
\FunctionTok{library}\NormalTok{(haven)}
\NormalTok{data }\OtherTok{\textless{}{-}} \FunctionTok{read\_sav}\NormalTok{(}\StringTok{"data/my\_qualtrics\_export.sav"}\NormalTok{)}
\end{Highlighting}
\end{Shaded}

Using \textbf{readr} from \textbf{tidyverse} for .csv files (Wickham et
al., 2019):

\begin{Shaded}
\begin{Highlighting}[]
\FunctionTok{library}\NormalTok{(tidyverse)}
\NormalTok{data }\OtherTok{\textless{}{-}}\NormalTok{ readr}\SpecialCharTok{::}\FunctionTok{read\_csv}\NormalTok{(}\StringTok{"data/my\_qualtrics\_export.csv"}\NormalTok{)}
\end{Highlighting}
\end{Shaded}

Using base R for R files.

\begin{Shaded}
\begin{Highlighting}[]
\NormalTok{data }\OtherTok{\textless{}{-}} \FunctionTok{readRDS}\NormalTok{(}\StringTok{"data/my\_data.RDS"}\NormalTok{)}
\end{Highlighting}
\end{Shaded}

\section{Loading reduced dataset}\label{loading-reduced-dataset}

Perhaps your data contains many variables that are not relevant to the
current project, you can load your data with specific variables only.

You can use the select() function from \textbf{tidyverse} to specify
which variables you would like.

\begin{Shaded}
\begin{Highlighting}[]
\NormalTok{data }\OtherTok{\textless{}{-}} \FunctionTok{read\_csv}\NormalTok{(}\StringTok{"data/my\_data.csv"}\NormalTok{) }\SpecialCharTok{\%\textgreater{}\%}
\NormalTok{  dplyr}\SpecialCharTok{::}\FunctionTok{select}\NormalTok{(var1, var2, var3)}
\end{Highlighting}
\end{Shaded}

\begin{tcolorbox}[enhanced jigsaw, toprule=.15mm, coltitle=black, colframe=quarto-callout-warning-color-frame, colbacktitle=quarto-callout-warning-color!10!white, opacityback=0, left=2mm, breakable, rightrule=.15mm, bottomrule=.15mm, leftrule=.75mm, titlerule=0mm, arc=.35mm, opacitybacktitle=0.6, toptitle=1mm, bottomtitle=1mm, title=\textcolor{quarto-callout-warning-color}{\faExclamationTriangle}\hspace{0.5em}{Common Error}, colback=white]

I often find that R produces errors when I do not include
\texttt{dplyr::} before the \texttt{select()} function. I'm not really
sure why this is -- perhaps there is a conflict between two packages I
often have loaded at the same time -- but whatever the reason, my
recommendation is to always make sure you're specifying that you're
calling the \texttt{select()} function from the \textbf{dplyr} package
in this way: \texttt{dplyr::select()}

\end{tcolorbox}

\section{Deleting specific variables}\label{deleting-specific-variables}

Perhaps you want most of the variables, but there are a select few that
you would like to drop. You can use the tidyverse package to drop
specific variables

\subsection{When loading the data}\label{when-loading-the-data}

\begin{Shaded}
\begin{Highlighting}[]
\NormalTok{data }\OtherTok{\textless{}{-}} \FunctionTok{read\_csv}\NormalTok{(}\StringTok{"data/my\_data.csv"}\NormalTok{) }\SpecialCharTok{\%\textgreater{}\%}
\NormalTok{  dplyr}\SpecialCharTok{::}\FunctionTok{select}\NormalTok{(}\SpecialCharTok{{-}}\NormalTok{var1, }\SpecialCharTok{{-}}\NormalTok{var2, }\SpecialCharTok{{-}}\NormalTok{var3)}
\end{Highlighting}
\end{Shaded}

\subsection{In an already loaded
dataframe}\label{in-an-already-loaded-dataframe}

\begin{Shaded}
\begin{Highlighting}[]
\NormalTok{data }\OtherTok{\textless{}{-}}\NormalTok{ data }\SpecialCharTok{\%\textgreater{}\%}
\NormalTok{  dplyr}\SpecialCharTok{::}\FunctionTok{select}\NormalTok{(}\SpecialCharTok{{-}}\NormalTok{var1)}
\end{Highlighting}
\end{Shaded}

\section{References}\label{references-1}

\chapter{Renaming Variables}\label{renaming-variables}

Rename variables using the \textbf{dplyr} package from
\textbf{tidyverse} (Wickham et al., 2019). This is useful for tidying up
cryptic export names (e.g., Q1\_1 from Qualtrics) into something more
meaningful.

\begin{Shaded}
\begin{Highlighting}[]
\FunctionTok{library}\NormalTok{(tidyverse)}
\NormalTok{data }\OtherTok{\textless{}{-}}\NormalTok{ data }\SpecialCharTok{\%\textgreater{}\%} 
  \FunctionTok{rename}\NormalTok{(}\AttributeTok{A =}\NormalTok{ q, }\AttributeTok{B =}\NormalTok{ e)}
\end{Highlighting}
\end{Shaded}

\textbf{What this is doing:}

\begin{itemize}
\item
  \texttt{rename()} changes column names without touching any of the
  data inside them.
\item
  The syntax is always \textbf{new name = old name} --- so
  \texttt{A\ =\ q} renames the column currently called \texttt{q} to
  \texttt{A}, and \texttt{B\ =\ e} renames \texttt{e} to \texttt{B}.
\end{itemize}

You can rename as many columns as you like in one call by adding more
\texttt{new\_name\ =\ old\_name} pairs, separated by commas:

\begin{Shaded}
\begin{Highlighting}[]
\NormalTok{data }\OtherTok{\textless{}{-}}\NormalTok{ data }\SpecialCharTok{\%\textgreater{}\%}
  \FunctionTok{rename}\NormalTok{(}\AttributeTok{A =}\NormalTok{ q,}
         \AttributeTok{B =}\NormalTok{ e,}
         \AttributeTok{C =}\NormalTok{ f,}
         \AttributeTok{D =}\NormalTok{ z)}
\end{Highlighting}
\end{Shaded}

\section{References}\label{references-2}

\chapter{Recoding Variables}\label{recoding-variables}

A collection of common recoding situations you'll run into when
preparing survey data for analysis.

\textbf{Packages required:}

\begin{itemize}
\tightlist
\item
  \textbf{tidyverse}: All examples on this page use \textbf{dplyr} (part
  of the tidyverse)
\end{itemize}

\begin{center}\rule{0.5\linewidth}{0.5pt}\end{center}

\section{\texorpdfstring{\textbf{Situation 1: Shifting a scale's
starting
point}}{Situation 1: Shifting a scale's starting point}}\label{situation-1-shifting-a-scales-starting-point}

In this first example, I have three items coded as 1--4, but I need them
to be coded 0--3.

\begin{Shaded}
\begin{Highlighting}[]
\CommentTok{\# Recode variables from 1{-}4 to 0{-}3}
\NormalTok{data }\OtherTok{\textless{}{-}}\NormalTok{ data }\SpecialCharTok{\%\textgreater{}\%}
  \FunctionTok{mutate}\NormalTok{(}\FunctionTok{across}\NormalTok{(}\FunctionTok{c}\NormalTok{(var1, var2, var3), }\SpecialCharTok{\textasciitilde{}}\NormalTok{ .x }\SpecialCharTok{{-}} \DecValTok{1}\NormalTok{, }\AttributeTok{.names =} \StringTok{"\{.col\}\_r"}\NormalTok{))}
\end{Highlighting}
\end{Shaded}

\textbf{What this is doing:}

\begin{itemize}
\item
  \texttt{across(c(var1,\ var2,\ var3),...)} applies the same
  transformation to all three variables at once
\item
  \texttt{\textasciitilde{}\ .x\ -\ 1} is the transformation itself:
  \texttt{.x} stands in for ``whatever variable is currently being
  processed,'' and this subtracts 1 from every value in it.
\item
  \texttt{.names\ =\ "\{.col\}\_r"} tells R to create new columns named
  var1\_r, var2\_r, and var3\_r, rather than overwriting the original
  variables.
\end{itemize}

\begin{center}\rule{0.5\linewidth}{0.5pt}\end{center}

\section{\texorpdfstring{\textbf{Situation 2: Reverse
coding}}{Situation 2: Reverse coding}}\label{situation-2-reverse-coding}

In this second example, I have items coded 1--7 that need to be
reverse-coded --- i.e.~a response of 1 should become 7, a response of 7
should become 1, and so on.

The formula for this is (n + 1) − original\_value, where n is the
highest point on the scale. For a 1--7 scale, that's 8 − .x, because:

\begin{itemize}
\tightlist
\item
  Original 1 becomes 8 − 1 = 7
\item
  Original 2 becomes 8 − 2 = 6
\item
  Original 3 becomes 8 − 3 = 5
\item
  Original 4 becomes 8 − 4 = 4 (midpoint stays the same)
\item
  Original 5 becomes 8 − 5 = 3
\item
  Original 6 becomes 8 − 6 = 2
\item
  Original 7 becomes 8 − 7 = 1
\end{itemize}

\textbf{For multiple items at once:}

\begin{Shaded}
\begin{Highlighting}[]
\NormalTok{data }\OtherTok{\textless{}{-}}\NormalTok{ data }\SpecialCharTok{\%\textgreater{}\%}
  \FunctionTok{mutate}\NormalTok{(}\FunctionTok{across}\NormalTok{(var1, var2, var3), }\SpecialCharTok{\textasciitilde{}} \DecValTok{8} \SpecialCharTok{{-}}\NormalTok{ .x, }\AttributeTok{.names =} \StringTok{"\{.col\}\_r"}\NormalTok{)}\ErrorTok{)}
\end{Highlighting}
\end{Shaded}

\textbf{For a single item:}

\begin{Shaded}
\begin{Highlighting}[]
\NormalTok{data }\OtherTok{\textless{}{-}}\NormalTok{ data }\SpecialCharTok{\%\textgreater{}\%}
  \FunctionTok{mutate}\NormalTok{(}\AttributeTok{var1\_r =} \DecValTok{8} \SpecialCharTok{{-}}\NormalTok{ var1)}
\end{Highlighting}
\end{Shaded}

\textbf{What this is doing:}

This follows the same logic as the first example.

\begin{itemize}
\item
  \texttt{across(c(var1,\ var2,\ var3),...)} applies the same
  transformation to all three variables at once
\item
  \texttt{\textasciitilde{}\ .x\ -\ 1} is the transformation itself:
  \texttt{.x} stands in for ``whatever variable is currently being
  processed,'' and this subtracts 1 from every value in it.
\item
  \texttt{.names\ =\ "\{.col\}\_r"} tells R to create new columns named
  var1\_r, var2\_r, and var3\_r, rather than overwriting the original
  variables.
\end{itemize}

The single item just skips \texttt{across()} since there's only one
variable to transform.

\textbf{Adjusting the formula for other scale lengths:}

To reverse-code for any scale, the formula is simply
\texttt{(highest\ scale\ point\ +\ 1)\ -\ original\ value}.

\begin{center}\rule{0.5\linewidth}{0.5pt}\end{center}

\section{\texorpdfstring{\textbf{Situation 3: Recoding
factors/levels}}{Situation 3: Recoding factors/levels}}\label{situation-3-recoding-factorslevels}

\subsection{Character variables to
factors}\label{character-variables-to-factors}

Categorical variables are often stored as \textbf{character} (text)
strings when they're first imported --- e.g.~\texttt{"control"} and
\texttt{"treatment"}, or \texttt{"low"}, \texttt{"medium"},
\texttt{"high"}. For most analyses, you'll want these as a
\textbf{factor} instead, since factors carry information about the
category \emph{levels} and their \emph{order}, which R uses when running
models or drawing legends.

\textbf{Check the unique values first:}

Before recoding, it's worth checking exactly what values currently exist
in the column --- this catches typos, inconsistent capitalisation
(\texttt{"Control"} vs \texttt{"control"}), or unexpected values before
they cause problems downstream.

\begin{Shaded}
\begin{Highlighting}[]
\FunctionTok{unique}\NormalTok{(data}\SpecialCharTok{$}\NormalTok{condition)}
\end{Highlighting}
\end{Shaded}

\textbf{What this is doing:} \texttt{unique()} returns every distinct
value present in \texttt{data\$condition}, with duplicates removed ---
so you get a quick list of exactly what you're working with.

\begin{tcolorbox}[enhanced jigsaw, toprule=.15mm, coltitle=black, colframe=quarto-callout-tip-color-frame, colbacktitle=quarto-callout-tip-color!10!white, opacityback=0, left=2mm, breakable, rightrule=.15mm, bottomrule=.15mm, leftrule=.75mm, titlerule=0mm, arc=.35mm, opacitybacktitle=0.6, toptitle=1mm, bottomtitle=1mm, title=\textcolor{quarto-callout-tip-color}{\faLightbulb}\hspace{0.5em}{Tip}, colback=white]

If you have a lot of variables to check, \texttt{table(data\$condition)}
is often more useful than \texttt{unique()} --- it shows you the same
distinct values, but also how many times each one appears, which makes
typos and rare stray values (e.g.~a single row with ``contorl'') much
easier to spot.

\end{tcolorbox}

\subsubsection{Converting to a factor --- the default (alphabetical)
behavior}\label{converting-to-a-factor-the-default-alphabetical-behavior}

If you convert a character variable straight to a factor without
specifying anything else, R will order the levels alphabetically by
default.

\begin{Shaded}
\begin{Highlighting}[]
\NormalTok{data }\OtherTok{\textless{}{-}}\NormalTok{ data }\SpecialCharTok{\%\textgreater{}\%}
  \FunctionTok{mutate}\NormalTok{(}\AttributeTok{condition =} \FunctionTok{factor}\NormalTok{(condition))}
\end{Highlighting}
\end{Shaded}

\textbf{What this is doing:}

\begin{itemize}
\item
  \texttt{factor()} converts the character column into a factor, but
  since no level order was specified, R falls back to sorting the
  distinct values alphabetically.
\item
  So \texttt{"control"} and \texttt{"treatment"} would end up ordered as
  control = 1, treatment = 2. This becomes a problem when you have
  something like \texttt{"low"}, \texttt{"medium"}, \texttt{"high"},
  which would end up being ordered high = 1, low = 2, medium = 3, which
  is probably not what you want.
\item
  This matters because the first level of a factor becomes the
  reference/baseline category in regression models by default. If R has
  chosen this alphabetically rather than by your intended design, your
  mode output and coefficient interpretation could be wrong.
\end{itemize}

\subsubsection{Converting to a factor --- with your own level
order}\label{converting-to-a-factor-with-your-own-level-order}

To control the order (and therefore which level is treated as the
reference/baseline), specify it explicitly using \texttt{levels\ =}:

\begin{Shaded}
\begin{Highlighting}[]
\NormalTok{data }\OtherTok{\textless{}{-}}\NormalTok{ data }\SpecialCharTok{\%\textgreater{}\%}
  \FunctionTok{mutate}\NormalTok{(}\AttributeTok{condition =} \FunctionTok{factor}\NormalTok{(condition, }\AttributeTok{levels =} \FunctionTok{c}\NormalTok{(}\StringTok{"control"}\NormalTok{, }\StringTok{"treatment"}\NormalTok{)))}
\end{Highlighting}
\end{Shaded}

\textbf{What this is doing:} the \texttt{levels\ =\ c(...)} argument
tells R exactly which order to use, overriding the alphabetical default.
Whatever value you list \textbf{first} becomes the reference level in
models.

For a naturally ordered scale (like low/medium/high), the same approach
applies:

\begin{Shaded}
\begin{Highlighting}[]
\NormalTok{data }\OtherTok{\textless{}{-}}\NormalTok{ data }\SpecialCharTok{\%\textgreater{}\%}
  \FunctionTok{mutate}\NormalTok{(}\AttributeTok{condition =} \FunctionTok{factor}\NormalTok{(condition, }\AttributeTok{levels =} \FunctionTok{c}\NormalTok{(}\StringTok{"low"}\NormalTok{, }\StringTok{"medium"}\NormalTok{, }\StringTok{"high"}\NormalTok{)))}
\end{Highlighting}
\end{Shaded}

If the variable has a genuine order (not just categories, but a rank ---
like low/medium/high), you can also mark it as an \textbf{ordered
factor} using \texttt{ordered\ =\ TRUE}. This tells R the levels are not
just labels, but have a true ranked order.

\begin{Shaded}
\begin{Highlighting}[]
\NormalTok{data }\OtherTok{\textless{}{-}}\NormalTok{ data }\SpecialCharTok{\%\textgreater{}\%}
  \FunctionTok{mutate}\NormalTok{(}\AttributeTok{condition =} \FunctionTok{factor}\NormalTok{(condition, }\AttributeTok{levels =} \FunctionTok{c}\NormalTok{(}\StringTok{"low"}\NormalTok{, }\StringTok{"medium"}\NormalTok{, }\StringTok{"high"}\NormalTok{), }\AttributeTok{ordered =} \ConstantTok{TRUE}\NormalTok{))}
\end{Highlighting}
\end{Shaded}

\section{Situation 4: Dummy coding}\label{situation-4-dummy-coding}

Perhaps you want to recode your descriptive or factor variable to a
numeric variable, e.g., 0 and 1.

\begin{Shaded}
\begin{Highlighting}[]
\NormalTok{data }\OtherTok{\textless{}{-}}\NormalTok{ data }\SpecialCharTok{\%\textgreater{}\%}
  \FunctionTok{mutate}\NormalTok{(}\AttributeTok{new\_variable =} \FunctionTok{case\_when}\NormalTok{(}
\NormalTok{    old\_variable }\SpecialCharTok{==} \StringTok{"control"} \SpecialCharTok{\textasciitilde{}} \DecValTok{0}\NormalTok{,}
    \ConstantTok{TRUE} \SpecialCharTok{\textasciitilde{}} \DecValTok{1}
\NormalTok{  ))}
\end{Highlighting}
\end{Shaded}

\textbf{What this is doing}:

\begin{itemize}
\item
  \texttt{case\_when()} checks each row against a series of conditions,
  in order, and assigns a value based on the first one that matches.
\item
  \texttt{old\_variable\ ==\ "control"\ \textasciitilde{}\ 0} is the
  first condition: wherever \texttt{old\_variable} equals
  \texttt{"control"}, assign \texttt{0} to the new column.
\item
  \texttt{TRUE\ \textasciitilde{}\ 1} acts as the catch-all/default: for
  any row that didn't match the condition above, assign \texttt{1}.
  Since \texttt{TRUE} is always true, this line always matches
  whatever's left over --- so in a two-group variable like
  control/treatment, everything that isn't ``control''
  (i.e.~```treatment''`) gets coded as \texttt{1}.
\item
  The result, \texttt{new\_variable}, is a plain numeric column (not a
  factor) --- useful when a function or model specifically expects 0/1
  numeric input rather than a labeled factor.
\end{itemize}

\subsection{Using the FastDummies
package}\label{using-the-fastdummies-package}

\textbf{fastDummies} (Kaplan, 2026) is a package which you can quickly
dummy-code multiple variables at once. You could create dummy variables
individually with \texttt{case\_when()} (as above), but that means
writing one line per level, per variable --- a tedious and error-prone
process once you have several categorical variables with multiple levels
each. The \texttt{dummy\_cols()} function in \textbf{fastDummies} does
this for as many variables as you like, at any number of levels, in a
single function call.

\begin{Shaded}
\begin{Highlighting}[]
\FunctionTok{library}\NormalTok{(fastDummies)}

\NormalTok{data }\OtherTok{\textless{}{-}} \FunctionTok{dummy\_cols}\NormalTok{(data, }\AttributeTok{select\_columns =} \FunctionTok{c}\NormalTok{(}\StringTok{"gender"}\NormalTok{, }\StringTok{"state"}\NormalTok{, }\StringTok{"age\_bracket"}\NormalTok{, }\StringTok{"ethnicity"}\NormalTok{))}
\end{Highlighting}
\end{Shaded}

\textbf{What this is doing:}

\begin{itemize}
\item
  \texttt{dummy\_cols()} takes a dataframe and one or more categorical
  columns, and creates a new binary (0/1) column for every
  leve\textbf{l} found in each of those columns.
\item
  \texttt{select\_columns\ =\ c("gender",\ "state",\ "age\_bracket",\ "ethnicity")}
  tells it exactly which columns to dummy-code at once.
\item
  By default, \texttt{dummy\_cols()} keeps the original columns and just
  adds the new dummy columns alongside them --- so \texttt{data} ends up
  with both \texttt{gender}(original) and new columns like
  \texttt{gender\_male}, \texttt{gender\_female} (or \texttt{gender\_1},
  \texttt{gender\_2}; one per level found in \texttt{gender}), plus the
  same pattern repeated for the other four variables.
\end{itemize}

There's also a function to drop the first level of a variable, if
necessary:

\begin{Shaded}
\begin{Highlighting}[]
\NormalTok{data }\OtherTok{\textless{}{-}} \FunctionTok{dummy\_cols}\NormalTok{(data, }\AttributeTok{select\_columns =} \FunctionTok{c}\NormalTok{(}\StringTok{"gender"}\NormalTok{, }\StringTok{"state"}\NormalTok{, }\StringTok{"age\_bracket"}\NormalTok{, }\StringTok{"ethnicity"}\NormalTok{, }\AttributeTok{remove\_first\_dummy =} \ConstantTok{TRUE}\NormalTok{))}
\end{Highlighting}
\end{Shaded}

\begin{itemize}
\tightlist
\item
  Turn this feature on if you're using dummy-coded variables in
  regression models, as including a dummy column for \textbf{every}
  level of a variable, alongside an intercept, creates perfect
  multicollinearity (known as the
  ``\href{https://www.statology.org/dummy-variable-trap/}{dummy variable
  trap}''). Dropping one level resolves this: the dropped level becomes
  your implicit reference/baseline category, and every other dummy
  column's coefficient is interpreted relative to it.
\item
  Levels are dropped alphabetically by default.
\end{itemize}

\section{Situation 5: Effect coding}\label{situation-5-effect-coding}

Effect coding (also called contrast coding) uses −1 and +1 instead of
0/1 to represent a two-level categorical variable. This changes how the
coefficient in a regression model is interpreted: with effect coding,
the coefficient represents the deviation from the grand mean, rather
than the difference from a reference group (as with standard 0/1 dummy
coding; (Cohen et al., 2013).

The same process is applied to create effect coded variables as above:

\begin{Shaded}
\begin{Highlighting}[]
\NormalTok{data }\OtherTok{\textless{}{-}}\NormalTok{ data }\SpecialCharTok{\%\textgreater{}\%}
  \FunctionTok{mutate}\NormalTok{(}\AttributeTok{condition\_effect =} \FunctionTok{case\_when}\NormalTok{(}
\NormalTok{    condition }\SpecialCharTok{==} \StringTok{"control"} \SpecialCharTok{\textasciitilde{}} \SpecialCharTok{{-}}\DecValTok{1}\NormalTok{,}
\NormalTok{    condition }\SpecialCharTok{==} \StringTok{"treatment"} \SpecialCharTok{\textasciitilde{}} \DecValTok{1}
\NormalTok{  ))}
\end{Highlighting}
\end{Shaded}

The result, \texttt{condition\_effect}, can be dropped straight into a
regression model in place of the original grouping variable, giving you
effect-coded coefficients.

\textbf{When to use this}:

\begin{itemize}
\item
  When you're running ANOVA-style designs through the \texttt{lm()}
  function. When you have a factorial design and want to run it as a
  regression rather than through \texttt{aov()}, effect coding makes the
  intercept represent the grand mean across all conditions, matching how
  ANOVA output is typically interpreted (Kugler et al., 2018).

  \begin{itemize}
  \item
    With standard dummy coding (0/1), the intercept instead represents
    the mean of whichever group happens to be coded 0, which usually
    isn't what you want to report.
  \item
    Dummy coding and effect coding will yield the same omnibus F
    statistic, but can produce different estimates, test statistics, and
    p-values for main and interaction effects.
  \end{itemize}
\item
  When you're interpreting main effects in the presence of interactions.
  In models with interaction terms, effect-coded main effects can be
  interpreted as averaged across the levels of the other factor(s).

  \begin{itemize}
  \tightlist
  \item
    With dummy coding, the coefficient represents the effect only within
    the reference group, which changes the interpretation substantially
    once an interaction is in the model (Cohen et al., 2013).
  \end{itemize}
\item
  When you're comparing coefficient magnitudes across multiple
  categorical predictors, because effect-coded predictors are more
  directly comparable to one another than dummy coded coefficients would
  be.
\item
  However, if your research question is specifically about comparing
  every other group back to one meaningful reference group (e.g., every
  treatment conditions vs a single control condition), then standard 0/1
  dummy coding is usually more interpretable.
\end{itemize}

\subsection{\texorpdfstring{\textbf{Situation 5: Converting NAs to
0}}{Situation 5: Converting NAs to 0}}\label{situation-5-converting-nas-to-0}

\subsubsection{For one variable}\label{for-one-variable}

Using base R:

\begin{Shaded}
\begin{Highlighting}[]
\NormalTok{data}\SpecialCharTok{$}\NormalTok{variable[}\FunctionTok{is.na}\NormalTok{(data}\SpecialCharTok{$}\NormalTok{variable)] }\OtherTok{\textless{}{-}} \DecValTok{0}
\end{Highlighting}
\end{Shaded}

\textbf{What this is doing:}

\begin{itemize}
\tightlist
\item
  \texttt{is.na(data\$variable)} checks every value in the column and
  returns \texttt{TRUE} wherever it's missing, \texttt{FALSE} otherwise.
\item
  \texttt{data\$variable{[}...{]}} uses that TRUE/FALSE result to select
  only the missing entries.
\item
  \texttt{\textless{}-\ 0} assigns 0 to just those selected (missing)
  entries, leaving every other value in the column untouched.
\end{itemize}

\subsubsection{For multiple, specific
variables}\label{for-multiple-specific-variables}

Using tidyverse syntax:

\begin{Shaded}
\begin{Highlighting}[]
\NormalTok{data }\OtherTok{\textless{}{-}}\NormalTok{ data }\SpecialCharTok{\%\textgreater{}\%}
  \FunctionTok{mutate}\NormalTok{(}\FunctionTok{across}\NormalTok{(}\FunctionTok{c}\NormalTok{(var1, var2, var3), }\SpecialCharTok{\textasciitilde{}} \FunctionTok{replace\_na}\NormalTok{(.x, }\DecValTok{0}\NormalTok{)))}
\end{Highlighting}
\end{Shaded}

\textbf{What this is doing:}

\begin{itemize}
\item
  \texttt{across(c(var1,\ var2,\ var3),\ ...)} applies the same
  transformation to each of the three named variables in turn, rather
  than writing a separate line for each one.
\item
  \texttt{\textasciitilde{}\ replace\_na(.x,\ 0)} is the transformation:
  \texttt{.x} stands in for whichever variable is currently being
  processed, and \texttt{replace\_na(.x,\ 0)} replaces any \texttt{NA}
  in it with \texttt{0}.
\item
  This overwrites \texttt{var1}, \texttt{var2}, and \texttt{var3} in
  place. If you'd rather keep the originals and create new columns
  instead, add \texttt{.names\ =\ "\{.col\}\_r"} as in earlier examples.
\end{itemize}

\subsubsection{For the whole dataframe}\label{for-the-whole-dataframe}

Using tidyverse syntax:

\begin{Shaded}
\begin{Highlighting}[]
\NormalTok{data }\OtherTok{\textless{}{-}}\NormalTok{ data }\SpecialCharTok{\%\textgreater{}\%}
  \FunctionTok{mutate}\NormalTok{(}\FunctionTok{across}\NormalTok{(}\FunctionTok{everything}\NormalTok{(), }\SpecialCharTok{\textasciitilde{}} \FunctionTok{replace\_na}\NormalTok{(.x, }\DecValTok{0}\NormalTok{)))}
\end{Highlighting}
\end{Shaded}

\textbf{What this is doing:}

\begin{itemize}
\item
  \texttt{across(everything(),\ ...)} applies the transformation to
  \textbf{every} column in the dataframe, rather than a specified
  subset.
\item
  \texttt{\textasciitilde{}\ replace\_na(.x,\ 0)} replaces any
  \texttt{NA} with \texttt{0}, the same as above, just run across all
  columns instead of a chosen few.
\end{itemize}

\begin{tcolorbox}[enhanced jigsaw, toprule=.15mm, coltitle=black, colframe=quarto-callout-important-color-frame, colbacktitle=quarto-callout-important-color!10!white, opacityback=0, left=2mm, breakable, rightrule=.15mm, bottomrule=.15mm, leftrule=.75mm, titlerule=0mm, arc=.35mm, opacitybacktitle=0.6, toptitle=1mm, bottomtitle=1mm, title=\textcolor{quarto-callout-important-color}{\faExclamation}\hspace{0.5em}{Important}, colback=white]

Applying this to the whole dataframe will also affect any character or
factor columns that happen to contain \texttt{NA}, not just numeric
ones. \texttt{replace\_na(.x,\ 0)} on a character column will throw an
error (since \texttt{0} isn't a valid character value), so this
whole-dataframe version really only works cleanly if every column is
numeric. If your dataframe has a mix of column types, it's safer to
either subset to just the numeric columns first
(\texttt{across(where(is.numeric),\ ...)}) or stick with the
named-variable version above.

\end{tcolorbox}

\section{Situation 6: Replacing a specific value with
another}\label{situation-6-replacing-a-specific-value-with-another}

\subsection{Special missing to NA}\label{special-missing-to-na}

Sometimes a placeholder value has been used during data collection
(e.g.~`-99` for ``not applicable'' or ``prefer not to say'') and needs
to be converted to a proper NA before analysis.

\subsubsection{For a single variable}\label{for-a-single-variable}

Using base R:

\begin{Shaded}
\begin{Highlighting}[]
\NormalTok{data}\SpecialCharTok{$}\NormalTok{variable[data}\SpecialCharTok{$}\NormalTok{variable }\SpecialCharTok{==} \SpecialCharTok{{-}}\DecValTok{99}\NormalTok{] }\OtherTok{\textless{}{-}} \ConstantTok{NA}
\end{Highlighting}
\end{Shaded}

\subsubsection{For every column in the dataframe at
once}\label{for-every-column-in-the-dataframe-at-once}

Using tidyverse syntax:

\begin{Shaded}
\begin{Highlighting}[]
\NormalTok{data }\OtherTok{\textless{}{-}}\NormalTok{ data }\SpecialCharTok{\%\textgreater{}\%}
  \FunctionTok{mutate}\NormalTok{(}\FunctionTok{across}\NormalTok{(}\FunctionTok{everything}\NormalTok{(), }\SpecialCharTok{\textasciitilde{}} \FunctionTok{na\_if}\NormalTok{(.x, }\SpecialCharTok{{-}}\DecValTok{99}\NormalTok{)))}
\end{Highlighting}
\end{Shaded}

Or, if you'd rather use base R:

\begin{Shaded}
\begin{Highlighting}[]
\NormalTok{data[data }\SpecialCharTok{==} \SpecialCharTok{{-}}\DecValTok{99}\NormalTok{] }\OtherTok{\textless{}{-}} \ConstantTok{NA}
\end{Highlighting}
\end{Shaded}

\subsection{SPSS special missing}\label{spss-special-missing}

If special missing values have been used in SPSS (e.g., -99 or 999
tagged as missing, but different from NA), there is a function in the
\textbf{haven} package to quickly convert them all to NA. I have not
used this myself, but you can find out about this function
\href{https://haven.tidyverse.org/reference/zap_missing.html}{here}
(Wickham et al., 2025).

\section{References}\label{references-3}

\chapter{Row/ID Numbers}\label{rowid-numbers}

If you want to add a sequential numbered variable for the number of rows
in the dataframe, such as an ID code, use the following code using base
R:

\begin{Shaded}
\begin{Highlighting}[]
\NormalTok{data}\SpecialCharTok{$}\NormalTok{id }\SpecialCharTok{\textless{}}\DecValTok{1}\SpecialCharTok{:}\FunctionTok{nrow}\NormalTok{(data)}
\end{Highlighting}
\end{Shaded}

This will add a variable to your dataframe called `id', which numbers
your participants sequentially.

\chapter{Saving Data}\label{saving-data}

\section{Saving to .csv}\label{saving-to-.csv}

\subsection{Using base R}\label{using-base-r}

\begin{Shaded}
\begin{Highlighting}[]
\FunctionTok{write.csv}\NormalTok{(data, }\StringTok{"data/data.csv"}\NormalTok{, }\AttributeTok{row.names =} \ConstantTok{FALSE}\NormalTok{)}
\end{Highlighting}
\end{Shaded}

\subsection{A better option, using
tidyverse}\label{a-better-option-using-tidyverse}

\begin{Shaded}
\begin{Highlighting}[]
\FunctionTok{library}\NormalTok{(tidyverse)}
\FunctionTok{write\_csv}\NormalTok{(data, }\StringTok{"data/data.csv"}\NormalTok{)}
\end{Highlighting}
\end{Shaded}

\begin{itemize}
\tightlist
\item
  Automatically does not save row names.
\end{itemize}

\section{Saving to .sav}\label{saving-to-.sav}

\begin{Shaded}
\begin{Highlighting}[]
\FunctionTok{library}\NormalTok{(haven)}
\FunctionTok{write\_sav}\NormalTok{(data, }\StringTok{"outdata.sav"}\NormalTok{)}
\end{Highlighting}
\end{Shaded}

\section{Saving as an R data file}\label{saving-as-an-r-data-file}

\begin{Shaded}
\begin{Highlighting}[]
\FunctionTok{saveRDS}\NormalTok{(df, }\StringTok{"data/data.rds"}\NormalTok{)}
\end{Highlighting}
\end{Shaded}

\textbf{The benefits of saving an R data file:}

\begin{itemize}
\tightlist
\item
  It preserves the complete R object, including data types (factors,
  dates, logicals, integers, lists, etc.), attributes (labels, metadata,
  row names, factor levels), complex objects (like models, lists,
  tibbles, spatial objects, plots, etc.), no additional parsing is
  required because data are saved exactly as they are, and they're
  usually quicker to read than .cvs or .sav files.
\end{itemize}

\part{Data Cleaning}

\chapter{Data Cleaning: Step-by-Step}\label{data-cleaning-step-by-step}

Researchers have different processes for preparing their data for
analysis, often which appears to be by convention. On this page, I will
set out my process and reasoning---take what works, leave what doesn't.
Notably, I conduct my reliability analysis and scale creation during my
data cleaning process, as this makes the most sense for me.

\section{Step 1: Participant
exclusions}\label{step-1-participant-exclusions}

The first step I like to do is exclude participants based on my
pre-registered criteria. One example is filtering out likely bot
responses using Qualtrics' built-in reCAPTCHA score.

\textbf{Q\_RecaptchaScore values \textless{} 0.5:} this is one of
Qualtrics' built-in measures. Scores range from 0 to 1; scores less than
0.5 are likely a bot, whereas scores greater than or equal to 0.5 are
likely to be human.

\begin{Shaded}
\begin{Highlighting}[]
\FunctionTok{library}\NormalTok{(tidyverse)}
\NormalTok{data }\OtherTok{\textless{}{-}}\NormalTok{ data }\SpecialCharTok{\%\textgreater{}\%}
  \FunctionTok{filter}\NormalTok{(Q\_RecaptchaScore }\SpecialCharTok{\textgreater{}=} \FloatTok{0.5}\NormalTok{)}
\end{Highlighting}
\end{Shaded}

\textbf{What this is doing:}

\begin{itemize}
\tightlist
\item
  \texttt{filter()} keeps only the rows that meet the condition you give
  it, and drops every row that doesn't --- here, that condition is
  \texttt{Q\_RecaptchaScore\ \textgreater{}=\ 0.5}.
\item
  \texttt{Q\_RecaptchaScore\ \textgreater{}=\ 0.5} checks each row's
  score and evaluates to \texttt{TRUE} for anyone at or above the 0.5
  threshold, \texttt{FALSE} for anyone below it. Only the \texttt{TRUE}
  rows are kept in the resulting \texttt{data}.
\item
  The result overwrites \texttt{data}, so from this point onward in your
  script, every subsequent step is working with the filtered sample
  rather than the original raw import. If you wanted to save this into a
  new dataframe, you would run the code
  \texttt{data\_filtered\ \textless{}-\ data\ \%\textgreater{}\%\ filter(Q\_RecaptchaScore\ \textgreater{}=\ 0.5)}
\end{itemize}

You can also exclude participants using the following code template:

\begin{Shaded}
\begin{Highlighting}[]
\NormalTok{data }\OtherTok{\textless{}{-}}\NormalTok{ data }\SpecialCharTok{\%\textgreater{}\%}
  \FunctionTok{filter}\NormalTok{(participant\_type }\SpecialCharTok{!=} \StringTok{"exclude"}\NormalTok{)}
\end{Highlighting}
\end{Shaded}

\textbf{What this is doing:}

\begin{itemize}
\tightlist
\item
  In this example, I have a character variable called
  \texttt{participant\_type} in which I have already flagged
  participants I wanted to exclude. The operators \texttt{!=} indicate
  ``does not equal,'' so we are \emph{not} retaining participants coded
  ``exclude.''
\end{itemize}

\section{Step 2: Drop, rename, and check variable
types}\label{step-2-drop-rename-and-check-variable-types}

Next, I like to clean up my datafile and
\href{load-specific-variables.qmd}{drop any unnecessary variables},
\href{renaming-variables.qmd}{rename them}, and
\href{check-variable-types.qmd}{ensure that no variables which should be
numeric are being read as character variables}.

\section{Step 3: Basic item-level
descriptives}\label{step-3-basic-item-level-descriptives}

I then like to make sure there are no glaringly obvious issues with the
data; for example, I haven't accidentally mapped two scale response
options to the same number. In that case, I would go back into Qualtrics
and recode that response and then redownload my data.

I use the \texttt{describe()} function from the \textbf{psych} package
(William Revelle, 2026) to quickly investigate: do the min/max values
make sense? Are special missing data codes present (e.g., -99)? I will
\href{recoding.qmd}{recode if necessary}.

\section{Step 4: Reliability and scale
creation}\label{step-4-reliability-and-scale-creation}

Scale by scale, I conduct reliability analysis and create my scales. I
use functions from the \textbf{psych} package (William Revelle, 2026)
for scales with 3 or more items, and spearman\_brown() from the
\textbf{splithalfr} package (Pronk, 2025) for scales with only 2 items.

\subsection{Cronbach's alpha}\label{cronbachs-alpha}

\begin{Shaded}
\begin{Highlighting}[]
\FunctionTok{library}\NormalTok{(psych)}
\FunctionTok{alpha}\NormalTok{(data[, }\FunctionTok{c}\NormalTok{(}\StringTok{"var1"}\NormalTok{, }\StringTok{"var2"}\NormalTok{, }\StringTok{"var3"}\NormalTok{)], }\AttributeTok{na.rm =}\NormalTok{ T)}
\end{Highlighting}
\end{Shaded}

\textbf{What this is doing:}

\begin{itemize}
\item
  \texttt{dat{[},\ c("var1",\ "var2",\ "var3"){]}} selects just the
  columns making up your scale.
\item
  \texttt{alpha()} computes Cronbach's alpha across those items, along
  with the fuller diagnostic output: item-total correlations, ``alpha if
  item dropped,'' and both raw and standardised alpha.
\item
  \texttt{na.rm\ =\ TRUE} tells \texttt{alpha()} to exclude missing
  values when calculating, rather than letting a single \texttt{NA} in
  any item cause the whole calculation to fail. Internally, this uses
  pairwise deletion for the correlation matrix \texttt{alpha()} relies
  on --- meaning each pair of items is correlated using whatever rows
  have data for \emph{both} of them, not just the rows that are complete
  across \emph{every} item in the scale.
\end{itemize}

\begin{tcolorbox}[enhanced jigsaw, toprule=.15mm, coltitle=black, colframe=quarto-callout-note-color-frame, colbacktitle=quarto-callout-note-color!10!white, opacityback=0, left=2mm, breakable, rightrule=.15mm, bottomrule=.15mm, leftrule=.75mm, titlerule=0mm, arc=.35mm, opacitybacktitle=0.6, toptitle=1mm, bottomtitle=1mm, title=\textcolor{quarto-callout-note-color}{\faInfo}\hspace{0.5em}{Reading the output}, colback=white]

The value you'll usually want to report is \textbf{\texttt{raw\_alpha}},
which is the standard Cronbach's alpha. The \texttt{std.alpha} value
(standardised alpha) is relevant if your items are on notably different
scales. The ``alpha if item dropped'' section shows what
\texttt{raw\_alpha} would be if that specific row's item were removed
from the scale. This is the most practically useful part of the output
for scale refinement; e.g., if the reliability is below the acceptable
threshold, perhaps this is due to one pesky item which should be dropped
from the scale.

\end{tcolorbox}

\subsection{McDonald's omega}\label{mcdonalds-omega}

Another measure of internal consistency is \emph{McDonald's omega}.
Omega estimates the reliability of a scale using a factor-analytic model
rather than a simple summed score (as with alpha).

Cronbach's alpha assumes that all the items measure the underlying
construct equally well, but in practice this assumption is rarely
met---some items just measure a construct better than others! When this
assumption is violated, Cronbach's alpha tends to underestimate the
scale reliability.

McDonald's omega, on the other hand, run a factor analysis on the items
first, then calculates reliability based on actual factor loadings so
that items which load more strongly onto the factor are weighted more
heavily than weaker items.

\begin{Shaded}
\begin{Highlighting}[]
\FunctionTok{library}\NormalTok{(psych)}
\FunctionTok{omega}\NormalTok{(data[, }\FunctionTok{c}\NormalTok{(}\StringTok{"var1"}\NormalTok{, }\StringTok{"var2"}\NormalTok{, }\StringTok{"var3"}\NormalTok{)], }\AttributeTok{na.rm =}\NormalTok{ T)}
\end{Highlighting}
\end{Shaded}

The output includes several different reliability coefficients, but the
two most commonly reported are:

\begin{itemize}
\item
  \textbf{\texttt{omega\_h}} (omega hierarchical): this is the
  proportion of variance in total scores attributable to a single
  general factor underlying all items. This is the closest omega
  equivalent to ``how much of this scale is measuring one coherent
  thing.''

  \begin{itemize}
  \tightlist
  \item
    This is typically what you want to report if using omega.
  \end{itemize}
\item
  \textbf{\texttt{omega\_t}} (omega total): this is the reliability of
  the total score treating all common variance (general factor plus any
  group/sub-factors) as true score. This is generally closer to what
  \texttt{alpha()} is estimating, but calculated via the factor model
  rather than raw covariances.
\end{itemize}

To my understanding, omega looks like it will be the way forward for new
scales/psychometric development, or at least that it will be reported
alongside alpha. However, it's not really useful for short scales, in
which case alpha (if 3 items) or the Spearman Brown coefficient (for 2
item scales) are more appropriate.

\subsection{Spearman Brown
Coefficient}\label{spearman-brown-coefficient}

Eisinga et al. (2013) compared Pearson r, Cronbach's alpha, and the
Spearman-Brown coefficient for assessing the reliability of two-item
scales. They concluded that Pearson r is not appropriate (in fact, were
quite sassy about it: they said ``one may rather call that the
reliability of a one-item test,'' p.~641), and nor is Cronbach's alpha.
Instead, they recommended using the Spearman-Brown coefficient, which is
generally higher than alpha and less biased on average.

You can calculate the Spearman-Brown coefficient using the
\textbf{splithalfr} package (Pronk, 2025).

\begin{Shaded}
\begin{Highlighting}[]
\FunctionTok{library}\NormalTok{(splithalfr)}
\FunctionTok{spearman\_brown}\NormalTok{(data}\SpecialCharTok{$}\NormalTok{var1, data}\SpecialCharTok{$}\NormalTok{var2,}
               \AttributeTok{use =} \StringTok{"pairwise.complete.obs"}\NormalTok{)}
\end{Highlighting}
\end{Shaded}

\textbf{What this is doing:}

\begin{itemize}
\tightlist
\item
  \texttt{use\ =\ "pairwise.complete.obs"} is necessary if you have any
  missing data; without this, if missing data is present, the result
  will be NA.
\end{itemize}

\subsection{Scale creation}\label{scale-creation}

If reliability is sufficient, I create my scales using the
\textbf{tidyverse} package (Wickham et al., 2019).

\begin{Shaded}
\begin{Highlighting}[]
\FunctionTok{library}\NormalTok{(tidyverse)}

\NormalTok{data }\OtherTok{\textless{}{-}}\NormalTok{ data }\SpecialCharTok{\%\textgreater{}\%}
  \FunctionTok{mutate}\NormalTok{(}\AttributeTok{scale =} \FunctionTok{rowMeans}\NormalTok{(}\FunctionTok{across}\NormalTok{(}\FunctionTok{c}\NormalTok{(var1, var2, var3)), }\AttributeTok{na.rm =} \ConstantTok{TRUE}\NormalTok{),}
         \AttributeTok{scale =} \FunctionTok{ifelse}\NormalTok{(}\FunctionTok{is.nan}\NormalTok{(scale), }\ConstantTok{NA}\NormalTok{, scale))}
\end{Highlighting}
\end{Shaded}

\textbf{What's happening here:}

\begin{itemize}
\item
  The first line creates \texttt{scale} and assigns it as a new column
  in \texttt{dat}.

  \begin{itemize}
  \item
    \texttt{across(c(var1,\ var2,\ var3))} selects the three items that
    make up the scale.
  \item
    \texttt{rowMeans(...,\ na.rm\ =\ TRUE)} calculates the average of
    those three columns per row (i.e.~per participant), rather than a
    single average across the whole column. \texttt{na.rm\ =\ TRUE}
    tells it to ignore any missing values when averaging, so a
    participant who answered 2 out of 3 items still gets a scale score
    based on those 2, rather than automatically becoming \texttt{NA}.
  \end{itemize}
\item
  The second line,
  \texttt{scale\ =\ ifelse(is.nan(scale),\ NA,\ scale)}, is a cleanup
  step: \texttt{rowMeans()} with \texttt{na.rm\ =\ TRUE} has a quirk
  where a row with \textbf{no} valid responses at all (all three items
  missing) doesn't return \texttt{NA} like you'd expect --- it returns
  \texttt{NaN} (``Not a Number'', R's result for an undefined
  calculation like averaging zero values). This line checks for that
  specific case with \texttt{is.nan()} and converts it to a proper
  \texttt{NA}, so downstream functions (which generally know how to
  handle \texttt{NA}, but may not handle \texttt{NaN} the same way)
  treat it correctly as missing.
\end{itemize}

\subsection{Alternatively\ldots{} reliability and scale creation in one
step}\label{alternatively-reliability-and-scale-creation-in-one-step}

I have recently discovered that the \textbf{psych} package has a
function called scoreItems() which allows you to calculate reliability
and create multiple scales in one step. First, you define the variables
that make up each scale and prefix any reverse-scored items with a minus
sign. Next, you score all scales at once. Finally, you can print the
reliability and descriptives for every scale at the same time.

\begin{Shaded}
\begin{Highlighting}[]
\FunctionTok{library}\NormalTok{(psych)}

\CommentTok{\# Define each scale\textquotesingle{}s items, prefixing reverse{-}scored items with a minus sign}
\NormalTok{scale\_keys }\OtherTok{\textless{}{-}} \FunctionTok{list}\NormalTok{(}
  \AttributeTok{scale\_1 =} \FunctionTok{c}\NormalTok{(}\StringTok{"var1"}\NormalTok{, }\StringTok{"var2"}\NormalTok{, }\StringTok{"{-}var3"}\NormalTok{, }\StringTok{"var4"}\NormalTok{), }\CommentTok{\# In this scale, var3 is reverse{-}coded}
  \AttributeTok{scale\_2 =} \FunctionTok{c}\NormalTok{(}\StringTok{"var5"}\NormalTok{, }\StringTok{"{-}var6"}\NormalTok{, }\StringTok{"var7"}\NormalTok{), }\CommentTok{\# var6 is reverse{-}coded}
  \AttributeTok{scale\_3 =} \FunctionTok{c}\NormalTok{(}\StringTok{"var8"}\NormalTok{, }\StringTok{"var9"}\NormalTok{, }\StringTok{"var10"}\NormalTok{, }\StringTok{"{-}var11"}\NormalTok{, }\StringTok{"var12"}\NormalTok{) }\CommentTok{\# var11 is reverse{-}coded}
\NormalTok{)}

\CommentTok{\# Score all scales at once}
\NormalTok{scores }\OtherTok{\textless{}{-}} \FunctionTok{scoreItems}\NormalTok{(scale\_keys, data)}

\CommentTok{\# Check reliability and descriptives for every scale together}
\FunctionTok{summary}\NormalTok{(scores)}

\CommentTok{\# Attach the scale scores back onto your main dataframe}
\NormalTok{data }\OtherTok{\textless{}{-}} \FunctionTok{cbind}\NormalTok{(data, scores}\SpecialCharTok{$}\NormalTok{scores)}
\end{Highlighting}
\end{Shaded}

\textbf{What this is doing:}

\begin{itemize}
\item
  \texttt{scale\_keys} is a named list, where each element defines one
  scale. The name you give each element (\texttt{scale\_1},
  \texttt{scale\_2}, \texttt{scale\_3}) becomes that scale's column name
  in the output.
\item
  Within each scale's item vector, a \texttt{-} prefix
  (e.g.~\texttt{"-var3"}) tells \texttt{scoreItems()} to automatically
  reverse-code that specific item before combining it with the rest ---
  no need to manually create \texttt{\_r} columns beforehand.
\item
  \texttt{scoreItems(scale\_keys,\ data)} scores every scale defined in
  \texttt{scale\_keys} in one call.
\item
  \texttt{summary(scores)} prints reliability (alpha) and basic
  descriptives for all three scales at once.
\item
  \texttt{scores\$scores} is a matrix with one column per scale, each
  row matching the corresponding row in \texttt{data}.
  \texttt{cbind(dat,\ scores\$scores)} binds these new columns onto your
  existing dataframe, so all three scale scores become permanent columns
  in \texttt{data}, ready to use in later analysis.
\end{itemize}

\begin{tcolorbox}[enhanced jigsaw, toprule=.15mm, coltitle=black, colframe=quarto-callout-note-color-frame, colbacktitle=quarto-callout-note-color!10!white, opacityback=0, left=2mm, breakable, rightrule=.15mm, bottomrule=.15mm, leftrule=.75mm, titlerule=0mm, arc=.35mm, opacitybacktitle=0.6, toptitle=1mm, bottomtitle=1mm, title=\textcolor{quarto-callout-note-color}{\faInfo}\hspace{0.5em}{Note}, colback=white]

It goes without saying, but if using this approach, don't reverse-code
your items manually before hand, otherwise they will be double
reverse-coded and just end up forward-coded again.

\end{tcolorbox}

\section{Step 6: Univariate Outliers}\label{step-6-univariate-outliers}

Some researchers like to transform univariate outliers, and some do not.
Both options introduce different forms of bias (transforming that
distort the underlying distribution and relationships between variables,
while leaving outliers untouched can influence means, correlations, and
model estimates). If you are the kind of researcher who would like to
transform your univariate outliers, this is the approach that worked for
me in the past.

\begin{tcolorbox}[enhanced jigsaw, toprule=.15mm, coltitle=black, colframe=quarto-callout-note-color-frame, colbacktitle=quarto-callout-note-color!10!white, opacityback=0, left=2mm, breakable, rightrule=.15mm, bottomrule=.15mm, leftrule=.75mm, titlerule=0mm, arc=.35mm, opacitybacktitle=0.6, toptitle=1mm, bottomtitle=1mm, title=\textcolor{quarto-callout-note-color}{\faInfo}\hspace{0.5em}{Note}, colback=white]

Sometimes you might want to check for univariate outliers on individual
items (e.g., for factor analysis, where the assumption is that outliers
are not present); but generally, I would conduct this analysis on the
scales rather than the items.

\end{tcolorbox}

\subsection{Identify scales with
outliers}\label{identify-scales-with-outliers}

First, I developed a function to identify scales on which outliers were
present, with outliers defined as those more than 3 standard deviations
below or above the mean.

\begin{Shaded}
\begin{Highlighting}[]
\NormalTok{detect\_outliers }\OtherTok{\textless{}{-}} \ControlFlowTok{function}\NormalTok{(df) \{}
\NormalTok{  outlier\_flags }\SpecialCharTok{\textless{}}\FunctionTok{sapply}\NormalTok{(df, }\ControlFlowTok{function}\NormalTok{(x) \{}
    \ControlFlowTok{if}\NormalTok{ (}\FunctionTok{is.numeric}\NormalTok{(x)) \{}
      \FunctionTok{any}\NormalTok{(x }\SpecialCharTok{\textgreater{}}\NormalTok{ (}\FunctionTok{mean}\NormalTok{(x, }\AttributeTok{na.rm =} \ConstantTok{TRUE}\NormalTok{) }\SpecialCharTok{+} \DecValTok{3} \SpecialCharTok{*} \FunctionTok{sd}\NormalTok{(x, }\AttributeTok{na.rm =} \ConstantTok{TRUE}\NormalTok{)) }\SpecialCharTok{|}
\NormalTok{            x }\SpecialCharTok{\textless{}}\NormalTok{ (}\FunctionTok{mean}\NormalTok{(x, }\AttributeTok{na.rm =} \ConstantTok{TRUE}\NormalTok{) }\SpecialCharTok{{-}} \DecValTok{3} \SpecialCharTok{*} \FunctionTok{sd}\NormalTok{(x, }\AttributeTok{na.rm =} \ConstantTok{TRUE}\NormalTok{)),}
          \AttributeTok{na.rm =} \ConstantTok{TRUE}\NormalTok{)  }\CommentTok{\# \textless{}{-}{-} added here}
\NormalTok{    \} }\ControlFlowTok{else}\NormalTok{ \{}
      \ConstantTok{FALSE}
\NormalTok{    \}}
\NormalTok{  \})}
  
\NormalTok{  outlier\_results }\OtherTok{\textless{}{-}} \FunctionTok{tibble}\NormalTok{(}
    \AttributeTok{Variable =} \FunctionTok{names}\NormalTok{(outlier\_flags),}
    \AttributeTok{Has\_Outliers =}\NormalTok{ outlier\_flags}
\NormalTok{  )}
  
  \FunctionTok{return}\NormalTok{(outlier\_results)}
\NormalTok{\}}

\CommentTok{\# Run univariate outlier analysis}
\FunctionTok{detect\_outlers}\NormalTok{(data[, }\FunctionTok{c}\NormalTok{(var1, var2, var3, var4)])}
\end{Highlighting}
\end{Shaded}

\textbf{What this is doing:}

\begin{itemize}
\item
  \texttt{detect\_outliers\ \textless{}-\ function(df)\ \{\ ...\ \}}
  defines a custom function --- rather than checking each variable one
  at a time, this packages the whole outlier-detection process so it can
  be reused on any set of columns, in this workbook or future projects,
  just by calling \texttt{detect\_outliers()} on a new dataframe.
\item
  \texttt{sapply(df,\ function(x)\ \{\ ...\ \})} runs the inner function
  on every column in \texttt{df} in turn, with \texttt{x} standing in
  for whichever column is currently being processed.
\item
  \texttt{if\ (is.numeric(x))\ \{\ ...\ \}\ else\ \{\ FALSE\ \}} ---
  this ensures that the function is only run on numeric variables. Any
  character, factor, or other non-numeric column just gets flagged
  \texttt{FALSE} (no outliers checked) automatically.
\item
  Inside the numeric branch, the outlier rule is:
  \texttt{x\ \textgreater{}\ (mean\ +\ 3*sd)\ \textbar{}\ x\ \textless{}\ (mean\ -\ 3*sd)}.
  This is the standard ``3 SDs from the mean'' definition of a
  univariate outlier --- anything further than 3 standard deviations
  above or below the mean (the \texttt{\textbar{}} means ``or'') is
  flagged. \texttt{mean(x,\ na.rm\ =\ TRUE)} and
  \texttt{sd(x,\ na.rm\ =\ TRUE)} both ignore missing values when
  calculating, so a few \texttt{NA}s in a column don't distort the
  threshold.
\item
  \texttt{any(...,\ na.rm\ =\ TRUE)} collapses that row-by-row
  \texttt{TRUE}/\texttt{FALSE} result down to a single answer per
  column: \texttt{TRUE} if \emph{any} value in that column meets the
  outlier criteria, \texttt{FALSE} if none do. The
  \texttt{na.rm\ =\ TRUE} here (on \texttt{any()}, not the inner mean/sd
  calls) makes sure that a missing value in the comparison itself
  doesn't cause the whole \texttt{any()} check to return \texttt{NA}
  instead of a real \texttt{TRUE}/\texttt{FALSE}.
\item
  \texttt{tibble(Variable\ =\ names(outlier\_flags),\ Has\_Outliers\ =\ outlier\_flags)}
  takes the results --- a named vector, one \texttt{TRUE}/\texttt{FALSE}
  per column --- and turns it into a tidy two-column table: which
  variable, and whether it has any outliers.
\item
  \texttt{df{[},\ c("var1",\ "var2",\ "var3",\ "var4"){]}} selects just
  the specific columns you want to check (your scale variables, say),
  rather than running the function on your entire dataframe.
\end{itemize}

The result is a small table listing every variable you checked and a
\texttt{TRUE}/\texttt{FALSE} for whether it contains at least one
outlier by this definition.

\subsection{Winsorise outliers}\label{winsorise-outliers}

If outliers are present, one technique is to Winsorise the items so that
responses that are greater than three standard deviations above or below
the mean are replaced with \textbf{exactly} 3 standard deviations above
or below the mean. Again, I created a function to do this across
multiple variables at once. The code below creates a new variable
identifiable by a suffix `win'.

\begin{Shaded}
\begin{Highlighting}[]
\NormalTok{winsorise\_outliers }\OtherTok{\textless{}{-}} \ControlFlowTok{function}\NormalTok{(data, var\_name, }\AttributeTok{na.rm =} \ConstantTok{TRUE}\NormalTok{) \{}
  \CommentTok{\# Get mean}
\NormalTok{  var\_m }\OtherTok{\textless{}{-}} \FunctionTok{mean}\NormalTok{(data[[var\_name]], }\AttributeTok{na.rm =}\NormalTok{ na.rm)}
  \CommentTok{\# Get sd}
\NormalTok{  var\_sd }\OtherTok{\textless{}{-}} \FunctionTok{sd}\NormalTok{(data[[var\_name]], }\AttributeTok{na.rm =}\NormalTok{ na.rm)}
  \CommentTok{\# Get lower and upper thresholds}
\NormalTok{  var\_lower }\OtherTok{\textless{}{-}}\NormalTok{ var\_m }\SpecialCharTok{{-}} \DecValTok{3} \SpecialCharTok{*}\NormalTok{ var\_sd}
\NormalTok{  var\_upper }\OtherTok{\textless{}{-}}\NormalTok{ var\_m }\SpecialCharTok{+} \DecValTok{3} \SpecialCharTok{*}\NormalTok{ var\_sd}
  
  \CommentTok{\# Create new variable name with \_win suffix}
\NormalTok{  new\_var\_name }\OtherTok{\textless{}{-}} \FunctionTok{paste0}\NormalTok{(var\_name, }\StringTok{"\_win"}\NormalTok{)}
  
  \CommentTok{\# Create new variable with suffix \_win, replacing values beyond the thresholds with the threshold}
\NormalTok{  data[[new\_var\_name]] }\OtherTok{\textless{}{-}} \FunctionTok{case\_when}\NormalTok{(}
\NormalTok{    data[[var\_name]] }\SpecialCharTok{\textgreater{}}\NormalTok{ var\_upper }\SpecialCharTok{\textasciitilde{}}\NormalTok{ var\_upper,}
\NormalTok{    data[[var\_name]] }\SpecialCharTok{\textless{}}\NormalTok{ var\_lower }\SpecialCharTok{\textasciitilde{}}\NormalTok{ var\_lower,}
    \ConstantTok{TRUE} \SpecialCharTok{\textasciitilde{}}\NormalTok{ data[[var\_name]]}
\NormalTok{  )}
  
  \FunctionTok{return}\NormalTok{(data)}
\NormalTok{\}}

\NormalTok{df }\OtherTok{\textless{}{-}} \FunctionTok{winsorise\_outliers}\NormalTok{(df, }\StringTok{"variable name"}\NormalTok{) }\CommentTok{\# For one variable}
\end{Highlighting}
\end{Shaded}

To winsorise several outliers at once, create a list of variables and
then loop over.

\begin{Shaded}
\begin{Highlighting}[]
\CommentTok{\# Create list of variables to loop over}
\NormalTok{vars\_to\_winsorise }\OtherTok{\textless{}{-}} \FunctionTok{c}\NormalTok{(}\StringTok{"var1"}\NormalTok{, }\StringTok{"var2"}\NormalTok{, }\StringTok{"var3"}\NormalTok{)}

\CommentTok{\# Run the loop}
\ControlFlowTok{for}\NormalTok{(var }\ControlFlowTok{in}\NormalTok{ vars\_to\_winsorise) \{}
\NormalTok{  df }\OtherTok{\textless{}{-}} \FunctionTok{winsorise\_outliers}\NormalTok{(df, var)}
\NormalTok{\}}
\end{Highlighting}
\end{Shaded}

\section{Step 7: Normality}\label{step-7-normality}

Next, check for normality. If you have Winsorised scales with outliers,
make sure to check the normality of the \emph{Winsorised} scales rather
than the raw scales.

\begin{tcolorbox}[enhanced jigsaw, toprule=.15mm, coltitle=black, colframe=quarto-callout-note-color-frame, colbacktitle=quarto-callout-note-color!10!white, opacityback=0, left=2mm, breakable, rightrule=.15mm, bottomrule=.15mm, leftrule=.75mm, titlerule=0mm, arc=.35mm, opacitybacktitle=0.6, toptitle=1mm, bottomtitle=1mm, title=\textcolor{quarto-callout-note-color}{\faInfo}\hspace{0.5em}{Why outliers before normality?}, colback=white]

Outliers inflate skewness and kurtosis statistics. Most normality checks
are sensitive to extreme values. A handful of outliers can make an
otherwise normal-ish distribution look meaningfully skewed. If you check
normality first, you might conclude a variable is non-normal and needs
transforming, when really it's a few outliers doing the damage.

\end{tcolorbox}

\subsection{Skew and kurtosis}\label{skew-and-kurtosis}

\textbf{Skew} measures asymmetry --- whether one tail of the
distribution is longer than the other.

\begin{itemize}
\item
  A skew of \textbf{0} indicates a perfectly symmetric distribution.
\item
  \textbf{Positive skew} (a long tail stretching to the right) means
  most scores cluster at the lower end, with a few high outlying scores
  pulling the tail out.
\item
  \textbf{Negative skew} (a long tail stretching to the left) means most
  scores cluster at the higher end, with a few low outlying scores ---
  common in ceiling-effect items.
\end{itemize}

\textbf{Kurtosis} measures ``tailedness'' --- how heavy or light a
distribution's tails are relative to a normal distribution.

\begin{itemize}
\item
  A kurtosis of \textbf{0} (using the \emph{excess} kurtosis convention,
  which is what \texttt{psych::describe()} reports) indicates tails
  similar to a normal distribution.
\item
  \textbf{Positive kurtosis} (leptokurtic) means heavier tails and a
  sharper peak than normal --- more extreme values than you'd expect
  under normality.
\item
  \textbf{Negative kurtosis} (platykurtic) means lighter tails and a
  flatter peak --- fewer extreme values than expected, a more ``spread
  out,'' uniform-looking distribution.
\end{itemize}

Use the \textbf{psych} package (William Revelle, 2026) to check the skew
and kurtosis.

\begin{Shaded}
\begin{Highlighting}[]
\FunctionTok{library}\NormalTok{(psych) }\FunctionTok{describe}\NormalTok{(dat[, }\FunctionTok{c}\NormalTok{(var1, var2, var3)]) }\CommentTok{\#Describe specific variables describe(dat) \# Describe every variable in the dataset}
\end{Highlighting}
\end{Shaded}

\textbf{What this is doing:}

\begin{itemize}
\tightlist
\item
  \texttt{describe()} produces a table of descriptive statistics for
  each variable --- including \texttt{mean}, \texttt{sd}, \texttt{min},
  \texttt{max}, and, importantly here, \textbf{\texttt{skew}} and
  \textbf{\texttt{kurtosis}} columns.
\end{itemize}

\textbf{Interpreting the output:}

There's no single universally agreed cutoff, but a commonly used rule of
thumb is that skew or kurtosis values beyond ±2 suggest a meaningful
departure from normality.

\subsection{Formal test of normality}\label{formal-test-of-normality}

The Shapiro-Wilk test gives you a formal statistical test of whether one
variable's distribution differs significantly from normal.

\begin{Shaded}
\begin{Highlighting}[]
\FunctionTok{shapiro.test}\NormalTok{(dat}\SpecialCharTok{$}\NormalTok{var1)}
\end{Highlighting}
\end{Shaded}

The output gives you a test statistic (\texttt{W}) and a p-value. If
significant, you reject the null hypothesis (i.e.~your data
significantly deviates from normal).

\begin{tcolorbox}[enhanced jigsaw, toprule=.15mm, coltitle=black, colframe=quarto-callout-caution-color-frame, colbacktitle=quarto-callout-caution-color!10!white, opacityback=0, left=2mm, breakable, rightrule=.15mm, bottomrule=.15mm, leftrule=.75mm, titlerule=0mm, arc=.35mm, opacitybacktitle=0.6, toptitle=1mm, bottomtitle=1mm, title=\textcolor{quarto-callout-caution-color}{\faFire}\hspace{0.5em}{Sensitive to sample size}, colback=white]

Shapiro-Wilk is notoriously sensitive to sample size {[}Yazici \&
Yolacan (2007){]}(Razali \& Wah, 2011). With a large sample (roughly N
\textgreater{} 300--500), even tiny, practically trivial deviations from
normality will return a significant p-value.

\end{tcolorbox}

\subsection{Visual checks}\label{visual-checks}

\subsubsection{Histogram/Density plot}\label{histogramdensity-plot}

A histogram gives you an intuitive sense of a variable's shape, and is
often more informative than a single test statistic, since you can
actually \emph{see} where any skew or unusual clustering is coming from.

You can do this really quickly using base R:

\begin{Shaded}
\begin{Highlighting}[]
\FunctionTok{hist}\NormalTok{(dat}\SpecialCharTok{$}\NormalTok{var)}
\end{Highlighting}
\end{Shaded}

Or using \textbf{ggplot2} from \textbf{tidyverse} (Wickham et al.,
2019). The code below shows a normal curve on the graph -- this means it
is actually a density plot (rather than a histogram).

\begin{Shaded}
\begin{Highlighting}[]
\FunctionTok{ggplot}\NormalTok{(dat, }\FunctionTok{aes}\NormalTok{(var1)) }\SpecialCharTok{+}
  \FunctionTok{theme}\NormalTok{(}\AttributeTok{legend.position =} \StringTok{"none"}\NormalTok{) }\SpecialCharTok{+}
  \FunctionTok{geom\_histogram}\NormalTok{(}\FunctionTok{aes}\NormalTok{(}\AttributeTok{y =} \FunctionTok{after\_stat}\NormalTok{(density)), }\AttributeTok{colour =} \StringTok{"black"}\NormalTok{, }\AttributeTok{fill =} \StringTok{"skyblue"}\NormalTok{, }\AttributeTok{binwidth =} \DecValTok{1}\NormalTok{) }\SpecialCharTok{+}
  \FunctionTok{labs}\NormalTok{(}\AttributeTok{x =} \StringTok{"Item Responses"}\NormalTok{, }\AttributeTok{y =} \StringTok{"Density"}\NormalTok{) }\SpecialCharTok{+}
  \FunctionTok{stat\_function}\NormalTok{(}\AttributeTok{fun =}\NormalTok{ dnorm, }\AttributeTok{args =} \FunctionTok{list}\NormalTok{(}\AttributeTok{mean =} \FunctionTok{mean}\NormalTok{(dat}\SpecialCharTok{$}\NormalTok{var1, }\AttributeTok{na.rm =} \ConstantTok{TRUE}\NormalTok{), }\AttributeTok{sd =} \FunctionTok{sd}\NormalTok{(dat}\SpecialCharTok{$}\NormalTok{var1, }\AttributeTok{na.rm =} \ConstantTok{TRUE}\NormalTok{)), }\AttributeTok{colour =} \StringTok{"black"}\NormalTok{, }\AttributeTok{linewidth =} \DecValTok{1}\NormalTok{) }\SpecialCharTok{+}
  \FunctionTok{theme\_classic}\NormalTok{() }\SpecialCharTok{+}
  \FunctionTok{scale\_x\_continuous}\NormalTok{(}\AttributeTok{breaks =} \FunctionTok{seq}\NormalTok{(}\DecValTok{1}\NormalTok{, }\DecValTok{7}\NormalTok{, }\DecValTok{1}\NormalTok{)) }\SpecialCharTok{+} \CommentTok{\# First number = min scale, second number = max scale, last number leave as 1}
  \FunctionTok{theme}\NormalTok{(}\AttributeTok{text =} \FunctionTok{element\_text}\NormalTok{(}\AttributeTok{size =} \DecValTok{25}\NormalTok{))}
\end{Highlighting}
\end{Shaded}

\textbf{What this is doing:}

\begin{itemize}
\tightlist
\item
  \texttt{ggplot(dat,\ aes(var1))} sets up the plot, telling ggplot
  which dataframe and variable to use.
\end{itemize}

\begin{itemize}
\tightlist
\item
  \texttt{theme(legend.position\ =\ "none")} removes the plot legend.
  Since this plot only has one histogram and one line (no grouping
  variable), there's nothing a legend would usefully explain.
\end{itemize}

\begin{itemize}
\item
  \texttt{geom\_histogram(aes(y\ =\ after\_stat(density)),\ ...)} draws
  the histogram.

  \begin{itemize}
  \item
    Normally a histogram's y-axis shows \textbf{raw counts} (how many
    participants fell in each bin), but that can't be directly compared
    to a normal curve, which is a probability density, not a count. In
    this code, \texttt{aes(y\ =\ after\_stat(density))}.
    \texttt{after\_stat(density)} tells ggplot to rescale the histogram
    bars so the total area under them sums to 1 --- the same scale a
    density curve uses --- so that a normal curve can be overlaid.
  \item
    \texttt{colour\ =\ "black"} outlines each bar,
    \texttt{fill\ =\ "skyblue"} colors them in, and
    \texttt{binwidth\ =\ 1} sets each bar to span 1 unit (sensible for a
    Likert-type item).
  \end{itemize}
\end{itemize}

\begin{itemize}
\tightlist
\item
  \texttt{labs(x\ =\ "Item\ Responses",\ y\ =\ "Density")} sets axis
  labels.
\end{itemize}

\begin{itemize}
\item
  \texttt{stat\_function(fun\ =\ dnorm,\ args\ =\ list(mean\ =\ ...,\ sd\ =\ ...),\ ...)}
  draws the actual normal curve.
\item
  \texttt{theme\_classic()} applies a clean, minimal ggplot theme
  (removes the default grey background and gridlines).
\end{itemize}

\begin{itemize}
\tightlist
\item
  \texttt{scale\_x\_continuous(breaks\ =\ seq(1,\ 7,\ 1))} controls
  where the tick marks appear on the x-axis. \texttt{seq(1,\ 7,\ 1)}
  generates the sequence 1, 2, 3, 4, 5, 6, 7 --- matching a 1--7
  response scale, so every scale point gets its own labeled tick mark
  rather than R choosing arbitrary breaks. If your scale was scored
  differently, you'd adjust the first two numbers to match your actual
  scale's minimum and maximum (e.g.~\texttt{seq(1,\ 5,\ 1)} for a 1--5
  scale).
\end{itemize}

\subsubsection{Q-Q plot}\label{q-q-plot}

A Q-Q (``quantile-quantile'') plot compares the data's distribution
directly against a theoretical normal distribution --- points falling
along the diagonal reference line indicate normality; systematic
deviation from that line indicates non-normality, and the \emph{shape}
of the deviation tells you what kind.

\begin{Shaded}
\begin{Highlighting}[]
\FunctionTok{ggplot}\NormalTok{(dat, }\FunctionTok{aes}\NormalTok{(}\AttributeTok{sample =}\NormalTok{ var1)) }\SpecialCharTok{+}
  \FunctionTok{stat\_qq}\NormalTok{() }\SpecialCharTok{+}
  \FunctionTok{stat\_qq\_line}\NormalTok{(}\AttributeTok{color =} \StringTok{"red"}\NormalTok{) }\SpecialCharTok{+}
  \FunctionTok{labs}\NormalTok{(}\AttributeTok{title =} \StringTok{"Q{-}Q Plot of Variable 1"}\NormalTok{)}
\end{Highlighting}
\end{Shaded}

\textbf{What this is doing:}

\begin{itemize}
\tightlist
\item
  \texttt{stat\_qq()} plots the data's quantiles against the quantiles
  you'd expect if the data were perfectly normally distributed.
\end{itemize}

\begin{itemize}
\tightlist
\item
  \texttt{stat\_qq\_line()} draws the reference diagonal line
  representing ``perfect normality'' --- the closer your points hug this
  line, the closer your data is to normal.
\end{itemize}

\begin{itemize}
\tightlist
\item
  Points curving away from the line at the ends (``tails'') typically
  indicate skew or heavy/light tails --- the same issue skew and
  kurtosis values are picking up numerically, but here you can see
  exactly \emph{where} in the distribution (upper end, lower end, or
  both) the deviation is occurring.
\end{itemize}

\subsection{So what\ldots?}\label{so-what}

None of these three checks --- skew/kurtosis, Shapiro-Wilk, and visual
plots --- is sufficient on its own. A sensible workflow is: use
skew/kurtosis as a fast first-pass screen across many variables at once
(as in the previous section), then for any variable that looks
borderline or concerning, follow up with a Shapiro-Wilk test
\textbf{and} a histogram/Q-Q plot together, since the visual check will
often explain \emph{why} the formal test came back significant
(e.g.~``ah, it's these three outlying values,'' rather than a genuinely
non-normal underlying shape).

\subsection{Transforming skewed data}\label{transforming-skewed-data}

If you have determined that your data \emph{is} skewed and you wish to
transform it, here is some helpful code.

Which transformation to use depends on two things: \textbf{how severe}
the skew is, and \textbf{which direction} it goes (positive or
negative). For more detail on choosing between these approaches, see
\href{https://rcompanion.org/handbook/I_12.html}{rcompanion's handbook
on transforming skewed data} and
\href{https://www.datanovia.com/en/lessons/transform-data-to-normal-distribution-in-r/}{Datanovia's
guide to normal distribution transformations}.

As a general starting point: try a \textbf{log transformation} first if
your DV increases more rapidly as your IV increases. Try a
\textbf{square transformation} instead if your DV decreases more rapidly
as your IV increases. The three options below (square root, log,
inverse) are ordered from mildest to strongest --- start with the
mildest transformation that resolves the skew, rather than jumping
straight to the strongest one available.

\subsubsection{Square root transformation for moderately skewed
data}\label{square-root-transformation-for-moderately-skewed-data}

\begin{Shaded}
\begin{Highlighting}[]
\CommentTok{\# For positive skew}
\NormalTok{dat }\OtherTok{\textless{}{-}}\NormalTok{ dat }\SpecialCharTok{\%\textgreater{}\%}
  \FunctionTok{mutate}\NormalTok{(}\AttributeTok{variable\_sqrt =} \FunctionTok{sqrt}\NormalTok{(variable))}

\CommentTok{\# For negative skew}
\NormalTok{dat }\OtherTok{\textless{}{-}}\NormalTok{ dat }\SpecialCharTok{\%\textgreater{}\%}
  \FunctionTok{mutate}\NormalTok{(}\AttributeTok{variable\_sqrt =} \FunctionTok{sqrt}\NormalTok{(}\FunctionTok{max}\NormalTok{(variable }\SpecialCharTok{+} \DecValTok{1}\NormalTok{) }\SpecialCharTok{{-}}\NormalTok{ variable))}
\end{Highlighting}
\end{Shaded}

\textbf{What this is doing:}

\begin{itemize}
\item
  For \textbf{positive skew}, \texttt{sqrt(variable)} takes the square
  root of every value. Since square roots compress large values more
  than small ones, this pulls in the long right tail characteristic of
  positive skew, bringing the distribution closer to symmetric.
\item
  For \textbf{negative skew}, the transformation is first
  \textbf{reflected} before taking the square root:
  \texttt{max(variable\ +\ 1)\ -\ variable} flips the distribution's
  direction (turning a left-skewed tail into a right-skewed one), and
  adding \texttt{1} to the max ensures no value in the result becomes
  exactly 0, since a square root transformation can't meaningfully
  handle 0 or negative values sitting right at the boundary.
  \texttt{sqrt(...)} is then applied to that reflected version.
\end{itemize}

\subsubsection{Log transformation for highly skewed
data}\label{log-transformation-for-highly-skewed-data}

\begin{Shaded}
\begin{Highlighting}[]
\CommentTok{\# For positive skew}
\NormalTok{dat }\OtherTok{\textless{}{-}}\NormalTok{ dat }\SpecialCharTok{\%\textgreater{}\%}
  \FunctionTok{mutate}\NormalTok{(}\AttributeTok{variable\_log =} \FunctionTok{log10}\NormalTok{(variable))}

\CommentTok{\# For negative skew}
\NormalTok{dat }\OtherTok{\textless{}{-}}\NormalTok{ dat }\SpecialCharTok{\%\textgreater{}\%}
  \FunctionTok{mutate}\NormalTok{(}\AttributeTok{variable\_log =} \FunctionTok{log10}\NormalTok{(}\FunctionTok{max}\NormalTok{(variable }\SpecialCharTok{+} \DecValTok{1}\NormalTok{) }\SpecialCharTok{{-}}\NormalTok{ variable))}
\end{Highlighting}
\end{Shaded}

\textbf{What this is doing:}

\begin{itemize}
\item
  \texttt{log10(variable)} applies a base-10 logarithm to every value.
  Logs compress large values far more aggressively than a square root
  does, which is why this is the next step up when a square root
  transformation isn't sufficient to resolve stronger skew.
\item
  Same reflection logic as above applies for negative skew:
  \texttt{max(variable\ +\ 1)\ -\ variable} flips the distribution's
  direction before the log is applied, and the \texttt{+\ 1} again
  avoids taking the log of 0 (which is mathematically undefined and
  would return \texttt{-Inf}).
\end{itemize}

\begin{tcolorbox}[enhanced jigsaw, toprule=.15mm, coltitle=black, colframe=quarto-callout-warning-color-frame, colbacktitle=quarto-callout-warning-color!10!white, opacityback=0, left=2mm, breakable, rightrule=.15mm, bottomrule=.15mm, leftrule=.75mm, titlerule=0mm, arc=.35mm, opacitybacktitle=0.6, toptitle=1mm, bottomtitle=1mm, title=\textcolor{quarto-callout-warning-color}{\faExclamationTriangle}\hspace{0.5em}{Warning}, colback=white]

\texttt{log10()} (and \texttt{log()}) will return \texttt{NaN} for any
negative value, and \texttt{-Inf} for exactly 0. If your variable can
genuinely contain 0 or negative values (not just as an artifact of the
reflection step above), you'll need to shift the whole variable to be
positive first --- e.g.~\texttt{log10(variable\ +\ 1)} --- before
transforming, or the function will silently produce unusable values for
those rows.

\end{tcolorbox}

\subsubsection{Inverse transformation for severely skewed
data}\label{inverse-transformation-for-severely-skewed-data}

\begin{Shaded}
\begin{Highlighting}[]
\CommentTok{\# For positive skew}
\NormalTok{dat }\OtherTok{\textless{}{-}}\NormalTok{ dat }\SpecialCharTok{\%\textgreater{}\%}
  \FunctionTok{mutate}\NormalTok{(}\AttributeTok{variable\_inverse =} \DecValTok{1} \SpecialCharTok{/}\NormalTok{ variable)}

\CommentTok{\# For negative skew}
\NormalTok{dat }\OtherTok{\textless{}{-}}\NormalTok{ dat }\SpecialCharTok{\%\textgreater{}\%}
  \FunctionTok{mutate}\NormalTok{(}\AttributeTok{variable\_inverse =} \DecValTok{1} \SpecialCharTok{/}\NormalTok{ (}\FunctionTok{max}\NormalTok{(variable }\SpecialCharTok{+} \DecValTok{1}\NormalTok{) }\SpecialCharTok{{-}}\NormalTok{ variable))}
\end{Highlighting}
\end{Shaded}

\textbf{What this is doing:}

\begin{itemize}
\item
  \texttt{1\ /\ variable} takes the reciprocal of every value --- the
  strongest of the three transformations here, and the most aggressive
  at compressing extreme values. This is generally reserved for cases
  where square root and log transformations have both been tried and
  haven't sufficiently resolved the skew.
\item
  Same reflection approach applies for negative skew, flipping the
  distribution's direction with
  \texttt{max(variable\ +\ 1)\ -\ variable} before taking the
  reciprocal.
\end{itemize}

\begin{tcolorbox}[enhanced jigsaw, toprule=.15mm, coltitle=black, colframe=quarto-callout-warning-color-frame, colbacktitle=quarto-callout-warning-color!10!white, opacityback=0, left=2mm, breakable, rightrule=.15mm, bottomrule=.15mm, leftrule=.75mm, titlerule=0mm, arc=.35mm, opacitybacktitle=0.6, toptitle=1mm, bottomtitle=1mm, title=\textcolor{quarto-callout-warning-color}{\faExclamationTriangle}\hspace{0.5em}{Warning}, colback=white]

I would recommend against the inverse transformation, to be honest. It
reverses the order of your values, which makes the transformed variable
much harder to interpret.

\end{tcolorbox}

\subsubsection{After transforming}\label{after-transforming}

Whichever transformation you use, re-run the skew/kurtosis check,
Shapiro-Wilk test, and histogram/Q-Q plot on the new transformed
variable to confirm whether the transformation actually resolved the
normality issue.

\subsection{Should you transform?}\label{should-you-transform}

Some data is genuinely, meaningfully non-normal by nature. Social change
and collective action data is a common example in this space (e.g.,
measures like willingness to protest, donation behavior, or
petition-signing frequently show floor effects with a long right tail of
highly engaged responders). This is often a real, substantive feature of
the underlying phenomenon, not a data quality issue (Ferr, 2025).

Rather than defaulting straight to transformation, a few approaches are
designed to handle non-normal data directly:

\begin{itemize}
\tightlist
\item
  Non-parametric tests (e.g.~Mann-Whitney U instead of an
  independent-samples t-test, Spearman's rho instead of Pearson's r)
  don't assume normality in the first place, and handle skewed or
  ordinal data well.
\end{itemize}

\begin{itemize}
\tightlist
\item
  Bootstrapping/resampling methods estimate variability empirically from
  your actual data, rather than relying on distributional assumptions
  --- a flexible option when the shape of your data doesn't map onto
  standard textbook methods.
\end{itemize}

\begin{itemize}
\tightlist
\item
  Robust regression estimators (e.g.~M-estimators, or the
  Satorra-Bentler scaled correction commonly used in structural equation
  modeling) can produce more accurate standard errors and significance
  tests under non-normality, without requiring you to transform the
  variable itself first.
\end{itemize}

\section{Step 8: Multivariate outlier
analysis}\label{step-8-multivariate-outlier-analysis}

This identifies participants who are unusual in their~\emph{pattern}~of
responses and flags them. This coded uses all base r functions. Run this
on the raw scores, not the Winsorised/transformed scores.

\begin{Shaded}
\begin{Highlighting}[]
\CommentTok{\# Multivariate outlier detection}
\NormalTok{dat}\SpecialCharTok{$}\NormalTok{mahal\_d }\OtherTok{\textless{}{-}} \FunctionTok{mahalanobis}\NormalTok{(}
\NormalTok{  dat[, }\FunctionTok{c}\NormalTok{(}\StringTok{"var1"}\NormalTok{, }\StringTok{"var2"}\NormalTok{, }\StringTok{"var3"}\NormalTok{, }\StringTok{"var4"}\NormalTok{)], }
  \FunctionTok{colMeans}\NormalTok{(dat[, }\FunctionTok{c}\NormalTok{(}\StringTok{"var1"}\NormalTok{, }\StringTok{"var2"}\NormalTok{, }\StringTok{"var3"}\NormalTok{, }\StringTok{"var4"}\NormalTok{)], }\AttributeTok{na.rm=}\NormalTok{T), }
  \FunctionTok{cov}\NormalTok{(dat[, }\FunctionTok{c}\NormalTok{(}\StringTok{"var1"}\NormalTok{, }\StringTok{"var2"}\NormalTok{, }\StringTok{"var3"}\NormalTok{, }\StringTok{"var4"}\NormalTok{)], }\AttributeTok{use =} \StringTok{"pairwise.complete.obs"}\NormalTok{)}
\NormalTok{)}

\NormalTok{dat}\SpecialCharTok{$}\NormalTok{mahal\_p }\OtherTok{\textless{}{-}} \FunctionTok{pchisq}\NormalTok{(dat}\SpecialCharTok{$}\NormalTok{mahal\_d, }\AttributeTok{df =}\NormalTok{ numberofvariables, }\AttributeTok{lower.tail =} \ConstantTok{FALSE}\NormalTok{) }\CommentTok{\#Note that you need to insert your number of variables included in the mahalanobis distance calculation above.}

\CommentTok{\# Flag multivariate outliers}
\NormalTok{dat }\OtherTok{\textless{}{-}}\NormalTok{ dat }\SpecialCharTok{\%\textgreater{}\%}
  \FunctionTok{mutate}\NormalTok{(}\AttributeTok{mv\_outlier =} \FunctionTok{ifelse}\NormalTok{(mahal\_p }\SpecialCharTok{\textless{}} \FloatTok{0.001}\NormalTok{, }\DecValTok{1}\NormalTok{, }\DecValTok{0}\NormalTok{))}
\end{Highlighting}
\end{Shaded}

\textbf{What this is doing:}

\begin{itemize}
\item
  \texttt{mahalanobis(...)} calculates the \textbf{Mahalanobis distance}
  for every participant --- a single number representing how far that
  participant's overall pattern of responses across
  \texttt{var1}--\texttt{var4} is from the ``average'' pattern,
  accounting for how those variables typically relate to each other.
  This is different from checking each variable separately: a
  participant could have entirely unremarkable individual scores, but
  still be a multivariate outlier if that particular \emph{combination}
  of scores is unusual given how the variables normally correlate.
\item
  The function needs three pieces, matching its three arguments here:

  \begin{itemize}
  \item
    \texttt{dat{[},\ c("var1",\ "var2",\ "var3",\ "var4"){]}} --- the
    set of variables you're checking.
  \item
    \texttt{colMeans(...,\ na.rm\ =\ TRUE)} --- the average score for
    each variable, which acts as the ``center'' that every participant's
    pattern is compared against.
  \item
    \texttt{cov(...,\ use\ =\ "pairwise.complete.obs")} --- the
    covariance matrix, capturing how the variables relate to and vary
    with each other. \texttt{use\ =\ "pairwise.complete.obs"} handles
    missing data.
  \end{itemize}
\item
  The result, stored in \texttt{dat\$mahal\_d}, is one Mahalanobis
  distance value per participant --- larger values mean a more unusual
  overall response pattern.
\item
  \texttt{pchisq(dat\$mahal\_d,\ df\ =\ numberofvariables,\ lower.tail\ =\ FALSE)}
  converts each distance into a p-value. \texttt{lower.tail\ =\ FALSE}
  gives you the probability of a distance \emph{at least this extreme}
  occurring by chance --- so a very small p-value means that
  participant's response pattern would be highly unlikely to occur under
  normal, expected variation.
\item
  \texttt{mutate(mv\_outlier\ =\ ifelse(mahal\_p\ \textless{}\ 0.001,\ 1,\ 0))}
  applies a conventional threshold --- \texttt{p\ \textless{}\ .001} is
  a commonly used cutoff for flagging multivariate outliers --- and
  creates a simple 0/1 flag column, making it easy to filter, count, or
  review flagged cases afterward.
\end{itemize}

Typically, you would run your analysis on the full sample and then run
sensitivity analyses excluding people with \texttt{mv\_outlier\ ==\ 1}
to see if the pattern of results holds without unusual participants.

\section{Step 9: Missing value
analysis}\label{step-9-missing-value-analysis}

Next, run a missing value analysis to determine whether your data are
missing completely at random. This uses the package \textbf{naniar}, and
has a hold bunch of useful functions that are documented
\href{https://naniar.njtierney.com/}{here}.

For the most part, I only check the proportion of missing data on each
variable, and use the Little's MCAR test, unless I need to do some
exploring.

First, run this code to determine the proportion of each variable which
is missing data. Variables missing more than 30\% are usually cause for
concern.

\begin{Shaded}
\begin{Highlighting}[]
\FunctionTok{install.packages}\NormalTok{(}\StringTok{"naniar"}\NormalTok{)}
\FunctionTok{library}\NormalTok{(naniar)}

\FunctionTok{miss\_var\_summary}\NormalTok{(data[, }\FunctionTok{c}\NormalTok{(}\StringTok{"var1"}\NormalTok{, }\StringTok{"var2"}\NormalTok{, }\StringTok{"var3"}\NormalTok{, }\StringTok{"var4"}\NormalTok{, }\StringTok{"var5"}\NormalTok{)]) }
\end{Highlighting}
\end{Shaded}

\textbf{What this is doing:}

\begin{itemize}
\tightlist
\item
  This provides a summary for each variable of the number, percentage,
  and cumulative sum of missings, ordered by the most missing.
\end{itemize}

If there is no missing data on any of your key variables, there is no
need to run Little's MCAR test (Little, 1988). If there is missing data,
run the following code:

\begin{Shaded}
\begin{Highlighting}[]
\FunctionTok{mcar\_test}\NormalTok{(data[, }\FunctionTok{c}\NormalTok{(}\StringTok{"var1"}\NormalTok{, }\StringTok{"var2"}\NormalTok{, }\StringTok{"var3"}\NormalTok{)])}
\end{Highlighting}
\end{Shaded}

\textbf{What this is doing:}

\begin{itemize}
\tightlist
\item
  Use Little's (1988) test statistic to assess if data is missing
  completely at random (MCAR). The null hypothesis in this test is that
  the data is MCAR, and the test statistic is a chi-squared value.
\end{itemize}

\section{Next Steps\ldots{}}\label{next-steps}

Once you have cleaned your data, created your scales, you might want to
generate \href{sample-descriptives.qmd}{sample demographics} and
\href{correlation-table.qmd}{correlations} tables.

\part{Reproducible Tables}

\chapter{Overview}\label{overview}

You could run your analyses and then manually build your tables, copying
over every relevant value by hand as you write up your results section.
Or\ldots{} you could use packages and functions to do that for you.

The benefit of this is threefold. First, it saves you time and cognitive
effort, since you're no longer trying to remember exactly which values a
given psychology journal expects to see reported, in what order, and to
how many decimal places. Second, it removes a huge number of
opportunities for human error. Rather than manually copying dozens (or
hundreds) of values --- many of which may need updating repeatedly
throughout drafting, revision, and re-analysis --- you can handle this
directly in R, so your tables update automatically whenever your
underlying analysis changes. Finally, it allows your tables to be
reproducible, so that if other researchers download your R code from
your OSF repository, they, too, can create the same tables you have
created.

In this section, I have included the code to create reproducible sample
demographics tables and correlations tables. The code to produce tables
for specific analyses can be found on the respective analysis pages.

\chapter{Sample Descriptives}\label{sample-descriptives}

The first artifact in your manuscript will probably be a table of sample
descriptives.

\section{Using the tableone package}\label{using-the-tableone-package}

This example shows an end-to-end workflow: recoding numeric-coded
demographic variables into labeled factors, building a ``Table 1''-style
summary table with \textbf{tableone}, converting it into a formatted
\textbf{flextable} object, and exporting the result directly to a Word
document.

The \textbf{tableone} package is great because it's whole purpose is
literally this, so no other functionality is needed. However, it does
not allow exporting directly to a Word document, so I have provided a
function to do this that some other fantastic, altruistic genius created
and made freely available.

\textbf{Packages needed:}

\begin{itemize}
\item
  dplyr (or tidyverse) (Wickham et al., 2019)
\item
  tableone (Yoshida \& Bartel, 2022)
\item
  flextable (Gohel \& Skintzos, 2026)
\item
  officer (Gohel et al., 2026)
\end{itemize}

\subsection{Step 1: Recode numeric codes into labelled
factors}\label{step-1-recode-numeric-codes-into-labelled-factors}

\begin{Shaded}
\begin{Highlighting}[]
\FunctionTok{library}\NormalTok{(dplyr)  data }\OtherTok{\textless{}{-}}\NormalTok{ data }\SpecialCharTok{\%\textgreater{}\%}   \FunctionTok{mutate}\NormalTok{(}\AttributeTok{gender\_r =} \FunctionTok{factor}\NormalTok{(gender,                           }\AttributeTok{levels =} \FunctionTok{c}\NormalTok{(}\DecValTok{1}\NormalTok{, }\DecValTok{2}\NormalTok{, }\DecValTok{3}\NormalTok{, }\DecValTok{4}\NormalTok{),                           }\AttributeTok{labels =} \FunctionTok{c}\NormalTok{(}\StringTok{"Male"}\NormalTok{, }\StringTok{"Female"}\NormalTok{, }\StringTok{"Non{-}binary"}\NormalTok{, }\StringTok{"Other"}\NormalTok{)),                    }\AttributeTok{state\_r =} \FunctionTok{factor}\NormalTok{(state,                           }\AttributeTok{levels =} \FunctionTok{c}\NormalTok{(}\DecValTok{1}\NormalTok{, }\DecValTok{2}\NormalTok{, }\DecValTok{3}\NormalTok{, }\DecValTok{4}\NormalTok{, }\DecValTok{5}\NormalTok{, }\DecValTok{6}\NormalTok{, }\DecValTok{7}\NormalTok{, }\DecValTok{8}\NormalTok{),                           }\AttributeTok{labels =} \FunctionTok{c}\NormalTok{(}\StringTok{"NSW"}\NormalTok{, }\StringTok{"VIC"}\NormalTok{, }\StringTok{"QLD"}\NormalTok{, }\StringTok{"SA"}\NormalTok{, }\StringTok{"WA"}\NormalTok{, }\StringTok{"TAS"}\NormalTok{, }\StringTok{"NT"}\NormalTok{, }\StringTok{"ACT"}\NormalTok{)),                    }\AttributeTok{education\_r =} \FunctionTok{factor}\NormalTok{(education,                          }\AttributeTok{levels =} \FunctionTok{c}\NormalTok{(}\DecValTok{1}\NormalTok{, }\DecValTok{2}\NormalTok{, }\DecValTok{3}\NormalTok{, }\DecValTok{4}\NormalTok{, }\DecValTok{5}\NormalTok{),                          }\AttributeTok{labels =} \FunctionTok{c}\NormalTok{(}\StringTok{"Less than high school"}\NormalTok{, }\StringTok{"High school"}\NormalTok{, }\StringTok{"TAFE"}\NormalTok{,                                      }\StringTok{"Undergraduate degree"}\NormalTok{, }\StringTok{"Postgraduate degree"}\NormalTok{)))}
\end{Highlighting}
\end{Shaded}

\textbf{What this is doing:}

\begin{itemize}
\tightlist
\item
  This builds on the \href{recoding.qmd}{factor-conversion approach},
  but adds a \texttt{labels\ =} argument alongside \texttt{levels\ =}.
  \texttt{levels} tells R the \emph{order} and \emph{original numeric
  values} to match against; \texttt{labels} tells it what \emph{text} to
  display in place of each number --- so \texttt{1} becomes
  \texttt{"Male"}, \texttt{2} becomes \texttt{"Female"}, and so on.
\end{itemize}

\subsection{Step 2: Build the summary table with
createTableOne()}\label{step-2-build-the-summary-table-with-createtableone}

\begin{Shaded}
\begin{Highlighting}[]
\FunctionTok{library}\NormalTok{(tableone)  vars }\OtherTok{\textless{}{-}} \FunctionTok{c}\NormalTok{(}\StringTok{"age"}\NormalTok{, }\StringTok{"gender\_r"}\NormalTok{, }\StringTok{"education\_r"}\NormalTok{)  table\_1 }\OtherTok{\textless{}{-}} \FunctionTok{CreateTableOne}\NormalTok{(}\AttributeTok{data =}\NormalTok{ data,                             }\AttributeTok{vars =}\NormalTok{ vars,                             }\AttributeTok{strata =} \StringTok{"state\_r"}\NormalTok{,                            }\AttributeTok{addOverall =} \ConstantTok{TRUE}\NormalTok{,                            }\AttributeTok{test =} \ConstantTok{FALSE}\NormalTok{) table\_1}
\end{Highlighting}
\end{Shaded}

\textbf{What this is doing:}

\begin{itemize}
\item
  \texttt{vars\ \textless{}-\ c("age",\ "gender\_r",\ "education\_r")}
  defines which variables should appear as rows in the table..
\item
  \texttt{CreateTableOne()} builds the summary table itself. It
  automatically detects each variable's type and picks an appropriate
  summary: means and SDs for continuous variables like \texttt{age},
  counts and percentages for factors like \texttt{gender\_r}and
  \texttt{education\_r}--- you don't need to specify this manually.
\item
  \texttt{strata\ =\ "state\_r"} splits the table into separate columns,
  one per state, so you can see how demographics differ by state at a
  glance, rather than only getting one pooled summary. This is useful if
  you have a representative survey; however, you could also list this as
  another count variable like gender and education. Similarly, you could
  add in an ethnicity variable, if relevant.
\item
  \texttt{addOverall\ =\ TRUE} adds an additional column showing the
  summary across \emph{all} participants combined, alongside the
  state-by-state breakdown. This is useful if you have stratified the
  table, as we have done here.
\item
  \texttt{test\ =\ FALSE} turns off tableone's default behavior of
  automatically running significance tests (e.g.~ANOVA, chi-square)
  comparing groups and adding a p-value column. Turn this off if the
  table is meant to be purely descriptive (as Table 1s usually are).
  Turn this on with \texttt{TRUE} if you specifically want group
  comparisons included.
\item
  Printing \texttt{table\_1} on its own displays the table in your R
  console.
\end{itemize}

\subsection{Step 3: Converting the tableone object into a
flextable}\label{step-3-converting-the-tableone-object-into-a-flextable}

This is a function helpfully created by Graham Cooper (Cooper, 2024) and
made available on GitHub to quickly convert a \textbf{tableone} object
to something that you can open in a Word document using the
\textbf{flextable} and \textbf{officer} packages. Make sure you have
those two packages installed and loaded to be able to run the function.

\begin{Shaded}
\begin{Highlighting}[]
\NormalTok{tableone2flextable }\OtherTok{\textless{}{-}} \ControlFlowTok{function}\NormalTok{(tableone)\{   tableone }\OtherTok{\textless{}{-}} \FunctionTok{print}\NormalTok{(tableone, }\AttributeTok{quote =} \ConstantTok{FALSE}\NormalTok{, }\AttributeTok{noSpaces =} \ConstantTok{TRUE}\NormalTok{)  }\CommentTok{\# convert to matrix      rows \textless{}{-} nrow(tableone)   cols \textless{}{-} ncol(tableone)   colnames \textless{}{-} colnames(tableone)      listoflists \textless{}{-} list()      for (i in 1:cols)\{     start \textless{}{-} (i * rows + 1) {-} rows     end \textless{}{-} i * rows     listoflists[[i]] \textless{}{-} tableone[start:end]   \}      dataframe \textless{}{-} as.data.frame(listoflists, col.names = colnames)   dataframe \textless{}{-} cbind(Variable = rownames(tableone), dataframe)     rownames(dataframe) \textless{}{-} NULL                                        flex \textless{}{-} flextable::flextable(dataframe)      return(flex) \}}
\end{Highlighting}
\end{Shaded}

\textbf{What this is doing:}

\begin{itemize}
\item
  \texttt{print(tableone,\ quote\ =\ FALSE,\ noSpaces\ =\ TRUE)}
  converts the tableone object into a plain \textbf{character matrix},
  stripping out quote marks and extra padding spaces that tableone
  normally adds for console display.
\item
  \texttt{nrow()}, \texttt{ncol()}, and \texttt{colnames()} capture the
  matrix's dimensions and column headers, which are needed to correctly
  rebuild the table as a dataframe in the next steps.
\item
  The \texttt{for} loop walks through the matrix column by column
  (\texttt{1:cols}), extracting each column's values from the matrix
  (which is stored internally as one long vector) using
  \texttt{start}/\texttt{end} index positions, and stores each column as
  one element of \texttt{listoflists}.
\item
  \texttt{as.data.frame(listoflists,\ col.names\ =\ colnames)} combines
  those extracted columns into an actual dataframe, restoring the
  original column headers.
\item
  \texttt{cbind(Variable\ =\ rownames(tableone),\ dataframe)} adds the
  row labels (variable names like ``age,'' ``gender\_r'') back in as
  their own proper column, since they'd otherwise just be dataframe row
  names rather than visible data.
\item
  \texttt{rownames(dataframe)\ \textless{}-\ NULL} resets the
  dataframe's row names back to plain sequential numbers, since the row
  names were just used above and aren't needed anymore as actual row
  labels.
\item
  \texttt{flextable::flextable(dataframe)} converts the finished
  dataframe into a \textbf{flextable} object --- the format needed for
  the polished formatting and Word export in the next step. Using
  \texttt{flextable::} here (rather than \texttt{library(flextable)}
  beforehand) explicitly specifies which package the function comes
  from, without needing to load the whole package into your session.
\end{itemize}

\subsection{Step 4: Format and export to
Word}\label{step-4-format-and-export-to-word}

\begin{Shaded}
\begin{Highlighting}[]
\NormalTok{flextable2word }\OtherTok{\textless{}{-}}\ControlFlowTok{function}\NormalTok{(flextab, filepath)\{            }\CommentTok{\# Add filepath so you can specify folder and name of file   require(tidyverse)   require(flextable)   ft \textless{}{-}autofit(theme\_apa(flextab)) \%\textgreater{}\%     align(j = 1, align = "left", part = "all") \%\textgreater{}\%        \# left{-}align first column     line\_spacing(space = 1, part = "all")                  \# remove double spacing   officer::read\_docx() \%\textgreater{}\%     body\_add\_flextable(ft) \%\textgreater{}\%     body\_end\_section\_landscape() \%\textgreater{}\%     print(target = filepath)   browseURL(filepath) \}  tab\_1 \textless{}{-}tableone2flextable(table\_1)    flextable2word(tab\_1, "output/Table 1 Sample demographics.docx")  tab\_1}
\end{Highlighting}
\end{Shaded}

\textbf{What this is doing:}

\begin{itemize}
\item
  \texttt{flextable2word} is a second custom function that takes a
  flextable object (from Step 3) and a \texttt{filepath} argument,
  letting you specify exactly where and under what filename the finished
  Word document should be saved.
\item
  \texttt{require(tidyverse)} and \texttt{require(flextable)} ensure
  both packages are loaded whenever this function runs, even if you
  forgot to load them earlier in your script.
\item
  \texttt{theme\_apa(flextab)} applies flextable's built-in
  \textbf{APA-style formatting} preset.
\item
  \texttt{autofit(...)} automatically adjusts column widths to fit their
  content, rather than leaving default widths that might cut off longer
  text.
\item
  \texttt{align(j\ =\ 1,\ align\ =\ "left",\ part\ =\ "all")}
  left-aligns the first column (\texttt{j\ =\ 1}, i.e.~your
  \texttt{Variable} column) across the whole table
  (\texttt{part\ =\ "all"} covers header and body).
\item
  \texttt{line\_spacing(space\ =\ 1,\ part\ =\ "all")} sets single line
  spacing throughout the table, overriding Word's common default of
  double spacing, which looks cluttered in a table context.
\item
  \texttt{officer::read\_docx()} starts a new, blank Word document
  object using the \textbf{officer} package (which flextable relies on
  for Word export) --- \texttt{officer::} here again explicitly
  specifies the package without a separate \texttt{library()} call.
\item
  \texttt{body\_add\_flextable(ft)} inserts your formatted table into
  that blank document.
\item
  \texttt{body\_end\_section\_landscape()} sets the page orientation to
  landscape for this section --- useful for wider tables (like this one,
  with multiple state columns) that wouldn't fit well on a standard
  portrait page.
\item
  \texttt{print(target\ =\ filepath)} writes the finished Word document
  out to disk, at the location and filename you specify in
  \texttt{filepath}.
\item
  \texttt{browseURL(filepath)} automatically opens the newly created
  Word file once it's saved, so you can immediately check the result
  without manually navigating to find it.
\end{itemize}

\subsection{Step 5: Run it}\label{step-5-run-it}

\begin{Shaded}
\begin{Highlighting}[]
\NormalTok{tab\_1 }\OtherTok{\textless{}{-}} \FunctionTok{tableone2flextable}\NormalTok{(table\_1)    }\FunctionTok{flextable2word}\NormalTok{(tab\_1, }\StringTok{"output/Table 1 Sample demographics.docx"}\NormalTok{)  tab\_1}
\end{Highlighting}
\end{Shaded}

\textbf{What this is doing:}

\begin{itemize}
\item
  \texttt{tableone2flextable(table\_1)} runs the Step 3 function on the
  table you built in Step 2, producing the flextable version.
\item
  \texttt{flextable2word(tab\_1,\ "output/Table\ 1\ Sample\ demographics.docx")}
  runs the Step 4 function, exporting the finished, formatted table
  straight to a Word file in your \texttt{output/} folder.
\item
  Running \texttt{tab\_1} on its own at the end displays the flextable
  directly in your R/RStudio viewer, letting you do a final visual
  check.
\end{itemize}

\chapter{Table of Correlations}\label{table-of-correlations}

The second table you will probably create in your manuscript is the
table of intercorrelations. I use the \textbf{apaTables} package to do
this.

\begin{itemize}
\tightlist
\item
  apaTables (Stanley, 2026)
\end{itemize}

\begin{Shaded}
\begin{Highlighting}[]
\FunctionTok{library}\NormalTok{(apaTables)}
\NormalTok{corr\_table }\OtherTok{\textless{}{-}} \FunctionTok{apa.cor.table}\NormalTok{(data[, }\FunctionTok{c}\NormalTok{(}\StringTok{"var1"}\NormalTok{, }\StringTok{"var2"}\NormalTok{, }\StringTok{"var3"}\NormalTok{, }\StringTok{"var4"}\NormalTok{)], }
                        \AttributeTok{table.number =} \DecValTok{2}\NormalTok{, }
                        \AttributeTok{show.sig.stars =} \ConstantTok{TRUE}\NormalTok{, }
                        \AttributeTok{landscape =} \ConstantTok{TRUE}\NormalTok{,}
                        \AttributeTok{filename =} \StringTok{"output/Table 2 Correlations pooled sample.doc"}\NormalTok{)}
\end{Highlighting}
\end{Shaded}

\textbf{What this is doing:}

\begin{itemize}
\item
  \texttt{corr\_table\ \textless{}-} saves the created table to an
  object in your R environment.
\item
  \texttt{data{[},\ c("var1",\ "var2",\ "var3",\ "var4"){]}} selects the
  variables to include, otherwise it will produce a correlations table
  for every variable in the dataframe.
\item
  \texttt{table.number\ =\ 2} allows you to set the number of the table
  in the produced artifact.
\item
  \texttt{landscape\ =\ TRUE} is pretty self-explanatory, it creates the
  table in a landscape oriented Word document. This might not be
  necessary if you only have a few variables; in this case, set this to
  \texttt{FALSE}.
\item
  \texttt{filename\ =\ "output/Table\ 2\ Correlations\ pooled\ sample.doc"}
  saves the table to a word document in your `outputs' folder in your R
  project.
\end{itemize}

You can then call it in your R environment using the code
\texttt{corr\_table}.

\part{Analyses}

\chapter{T-Tests}\label{t-tests}

I use the \textbf{rempsyc} package to run t-tests, which is a
``wrapper'' on top of other functions designed to produce APA-formatted
t-test tables directly. Rather than running a t-test, then manually
formatting the output into a table, \textbf{rempsyc} does both in two
function calls.

\textbf{Packages required:}

\begin{itemize}
\tightlist
\item
  rempsych (Thériault, 2023)
\end{itemize}

\section{Running a single t-test}\label{running-a-single-t-test}

\begin{Shaded}
\begin{Highlighting}[]
\FunctionTok{library}\NormalTok{(rempsyc)}

\NormalTok{t\_test\_results }\OtherTok{\textless{}{-}} \FunctionTok{nice\_t\_test}\NormalTok{(}
  \AttributeTok{data =}\NormalTok{ dat,}
  \AttributeTok{response =} \StringTok{"dependent\_variable"}\NormalTok{,}
  \AttributeTok{group =} \StringTok{"grouping\_variable"}
\NormalTok{)}

\NormalTok{t\_test\_table }\OtherTok{\textless{}{-}} \FunctionTok{nice\_table}\NormalTok{(t\_test\_results)}
\NormalTok{t\_test\_table}
\end{Highlighting}
\end{Shaded}

\textbf{What this is doing:}

\begin{itemize}
\item
  \texttt{nice\_t\_test()} runs the actual t-test. Unlike base R's
  \texttt{t.test()}, it returns results already formatted for reporting,
  including the test statistic, degrees of freedom, p-value, and an
  effect size (Cohen's d), all in one output.
\item
  \texttt{response\ =\ "dependent\_variable"} and
  \texttt{group\ =\ "grouping\_variable"} are given as \textbf{character
  strings} (in quotes) rather than bare column names.
\item
  \texttt{nice\_table()} takes the results from \texttt{nice\_t\_test()}
  and converts them into a polished, APA-styled table, with leading
  zeros dropped.
\end{itemize}

\section{Running multiple t-tests at
once}\label{running-multiple-t-tests-at-once}

The function accepts either a single variable or a vector of several
variables (e.g.~\texttt{response\ =\ c("var1",\ "var2",\ "var3")}) to
run multiple t-tests at once, all formatted consistently in the same
output.

\begin{Shaded}
\begin{Highlighting}[]
\NormalTok{t\_test\_results }\OtherTok{\textless{}{-}} \FunctionTok{nice\_t\_test}\NormalTok{(}
  \AttributeTok{data =}\NormalTok{ dat,}
  \AttributeTok{response =} \FunctionTok{c}\NormalTok{(}\StringTok{"var1"}\NormalTok{, }\StringTok{"var2"}\NormalTok{, }\StringTok{"var3"}\NormalTok{, }\StringTok{"var4"}\NormalTok{),}
  \AttributeTok{group =} \StringTok{"grouping\_variable"}\NormalTok{,}
  \AttributeTok{correction =} \StringTok{"bonferroni"}
\NormalTok{)}

\NormalTok{t\_test\_table }\OtherTok{\textless{}{-}} \FunctionTok{nice\_table}\NormalTok{(t\_test\_results)}
\NormalTok{t\_test\_table}
\end{Highlighting}
\end{Shaded}

\textbf{What this is doing:}

\begin{itemize}
\item
  \texttt{correction\ =\ "bonferroni"} applies a \textbf{Bonferroni
  correction} to the p-values across the four tests.
\item
  Bonferroni is, of course, the most conservative correction for
  multiple tests. You can select the \textbf{Holm} correction
  (\texttt{correction\ =\ "holm"}) or \textbf{false discovery rate}
  methods (\texttt{correction\ =\ "fdr"}) for less conservative
  corrections that are also supported by \texttt{nice\_t\_test()}.

  \begin{itemize}
  \tightlist
  \item
    Perneger (1998) discusses when Bonferroni adjustments might be the
    wrong choice, and Benjamini \& Hochberg (1995) explain the false
    discovery rate method.
  \end{itemize}
\end{itemize}

\section{Running t-tests using lavaan/SEM
frameworks}\label{running-t-tests-using-lavaansem-frameworks}

I can't imagine this being particularly useful to anyone but me, but I
once had to figure out how to fit a t-test but using full information
maximum likelihood (FIML). For background, a large representative survey
contained around 50\% missingness by design (to reduce questionnaire
completion time and cognitive load for participants). Running standard
t-tests would therefore require lisewise deletion of a devastatingly
large amount of data.

My supervisor asked me to devise a way to run t-tests using FIML
estimation, which allows partially missing cases to be retained and
contribute to the overall model estimation, reducing bias.\footnote{And
  can handle up to 60\% missing data (Enders \& Bandalos, 2001)!}
However, FIML is an estimation method used in structural equation
modelling and is not something that has been widely incorporated into
other statistical analyses. As such, I had to figure out how to run the
\emph{equivalent} of a t-test using the SEM framework.

In order to run a t-test equivalent with FIML in R, the following
packages are required:

\begin{itemize}
\tightlist
\item
  lavaan {[}Rosseel et al. (2026){]}(Rosseel, 2012)
\end{itemize}

\subsection{Building the models}\label{building-the-models}

Essentially what I wanted to do was compare two different models: one in
which the mean of the variable of interest is constrained to be equal
(equivalent to the null hypothesis) and one in which the mean is able to
vary (alternative hypothesis). I used the multigroups function in
\textbf{lavaan}, which allowed me to specify the grouping variable.
Variance in the dependent variables is freely estimated across groups to
approximate a Welch t-test when group sizes are unequal. The models are
then compared. A significant likelihood ratio test indicates that the
unconstrained model fit the data better than the constrained model;
therefore, the groups differed significantly.

The code below shows a full workflow in which I listed all the variables
I wanted to run t-tests on (in this case, many!), created an analysis
function, looped it over all the dependent variables, and exported the
results to a table with correction applied for multiple testing.

\subsection{Step 1: List dependent
variables}\label{step-1-list-dependent-variables}

\begin{Shaded}
\begin{Highlighting}[]
\StringTok{\textasciigrave{}}\AttributeTok{\%||\%}\StringTok{\textasciigrave{}} \OtherTok{\textless{}{-}} \ControlFlowTok{function}\NormalTok{(a, b) }\ControlFlowTok{if}\NormalTok{ (}\SpecialCharTok{!}\FunctionTok{is.null}\NormalTok{(a)) a }\ControlFlowTok{else}\NormalTok{ b}

\NormalTok{dv\_config\_controlled }\OtherTok{\textless{}{-}} \FunctionTok{list}\NormalTok{(}
  \FunctionTok{list}\NormalTok{(}\AttributeTok{dv =} \StringTok{"stress"}\NormalTok{,                  }\AttributeTok{scale\_factor =} \ConstantTok{NULL}\NormalTok{),}
  \FunctionTok{list}\NormalTok{(}\AttributeTok{dv =} \StringTok{"dangerous"}\NormalTok{,               }\AttributeTok{scale\_factor =} \ConstantTok{NULL}\NormalTok{),}
  \FunctionTok{list}\NormalTok{(}\AttributeTok{dv =} \StringTok{"emotional\_wb"}\NormalTok{,            }\AttributeTok{scale\_factor =} \ConstantTok{NULL}\NormalTok{),}
  \FunctionTok{list}\NormalTok{(}\AttributeTok{dv =} \StringTok{"psych\_wb"}\NormalTok{,                }\AttributeTok{scale\_factor =} \ConstantTok{NULL}\NormalTok{),}
  \FunctionTok{list}\NormalTok{(}\AttributeTok{dv =} \StringTok{"swb\_contribution"}\NormalTok{,        }\AttributeTok{scale\_factor =} \ConstantTok{NULL}\NormalTok{),}
  \FunctionTok{list}\NormalTok{(}\AttributeTok{dv =} \StringTok{"swb\_acceptance"}\NormalTok{,          }\AttributeTok{scale\_factor =} \ConstantTok{NULL}\NormalTok{),}
  \FunctionTok{list}\NormalTok{(}\AttributeTok{dv =} \StringTok{"swb\_integration"}\NormalTok{,         }\AttributeTok{scale\_factor =} \ConstantTok{NULL}\NormalTok{),}
  \FunctionTok{list}\NormalTok{(}\AttributeTok{dv =} \StringTok{"swb\_own\_actualisation"}\NormalTok{,   }\AttributeTok{scale\_factor =} \ConstantTok{NULL}\NormalTok{),}
  \FunctionTok{list}\NormalTok{(}\AttributeTok{dv =} \StringTok{"swb\_other\_actualisation"}\NormalTok{, }\AttributeTok{scale\_factor =} \ConstantTok{NULL}\NormalTok{),}
  \FunctionTok{list}\NormalTok{(}\AttributeTok{dv =} \StringTok{"swb\_coherence"}\NormalTok{,           }\AttributeTok{scale\_factor =} \ConstantTok{NULL}\NormalTok{),}
  \FunctionTok{list}\NormalTok{(}\AttributeTok{dv =} \StringTok{"anomie"}\NormalTok{,                  }\AttributeTok{scale\_factor =} \ConstantTok{NULL}\NormalTok{),}
  \FunctionTok{list}\NormalTok{(}\AttributeTok{dv =} \StringTok{"prej\_immigrant\_s"}\NormalTok{,        }\AttributeTok{scale\_factor =} \DecValTok{10}\NormalTok{),}
  \FunctionTok{list}\NormalTok{(}\AttributeTok{dv =} \StringTok{"prej\_muslim\_s"}\NormalTok{,           }\AttributeTok{scale\_factor =} \DecValTok{10}\NormalTok{),}
  \FunctionTok{list}\NormalTok{(}\AttributeTok{dv =} \StringTok{"prej\_jewish\_s"}\NormalTok{,           }\AttributeTok{scale\_factor =} \DecValTok{10}\NormalTok{),}
  \FunctionTok{list}\NormalTok{(}\AttributeTok{dv =} \StringTok{"collangst"}\NormalTok{,               }\AttributeTok{scale\_factor =} \ConstantTok{NULL}\NormalTok{),}
  \FunctionTok{list}\NormalTok{(}\AttributeTok{dv =} \StringTok{"fear\_immigrants"}\NormalTok{,         }\AttributeTok{scale\_factor =} \ConstantTok{NULL}\NormalTok{),}
  \FunctionTok{list}\NormalTok{(}\AttributeTok{dv =} \StringTok{"norms\_immigration\_s"}\NormalTok{,     }\AttributeTok{scale\_factor =} \DecValTok{10}\NormalTok{),}
  \FunctionTok{list}\NormalTok{(}\AttributeTok{dv =} \StringTok{"norms\_racism\_s"}\NormalTok{,          }\AttributeTok{scale\_factor =} \DecValTok{10}\NormalTok{),}
  \FunctionTok{list}\NormalTok{(}\AttributeTok{dv =} \StringTok{"norms\_demo\_s"}\NormalTok{,            }\AttributeTok{scale\_factor =} \DecValTok{10}\NormalTok{),}
  \FunctionTok{list}\NormalTok{(}\AttributeTok{dv =} \StringTok{"poleff"}\NormalTok{,                  }\AttributeTok{scale\_factor =} \ConstantTok{NULL}\NormalTok{),}
  \FunctionTok{list}\NormalTok{(}\AttributeTok{dv =} \StringTok{"sysjust"}\NormalTok{,                 }\AttributeTok{scale\_factor =} \ConstantTok{NULL}\NormalTok{),}
  \FunctionTok{list}\NormalTok{(}\AttributeTok{dv =} \StringTok{"demo\_indecisive"}\NormalTok{,         }\AttributeTok{scale\_factor =} \ConstantTok{NULL}\NormalTok{),}
  \FunctionTok{list}\NormalTok{(}\AttributeTok{dv =} \StringTok{"demo\_best\_option"}\NormalTok{,        }\AttributeTok{scale\_factor =} \ConstantTok{NULL}\NormalTok{),}
  \FunctionTok{list}\NormalTok{(}\AttributeTok{dv =} \StringTok{"undemo"}\NormalTok{,                  }\AttributeTok{scale\_factor =} \ConstantTok{NULL}\NormalTok{),}
  \FunctionTok{list}\NormalTok{(}\AttributeTok{dv =} \StringTok{"polviol"}\NormalTok{,                 }\AttributeTok{scale\_factor =} \ConstantTok{NULL}\NormalTok{),}
  \FunctionTok{list}\NormalTok{(}\AttributeTok{dv =} \StringTok{"anger\_govt"}\NormalTok{,              }\AttributeTok{scale\_factor =} \ConstantTok{NULL}\NormalTok{)}
\NormalTok{)}
\end{Highlighting}
\end{Shaded}

\textbf{What this is doing:}

\begin{itemize}
\item
  \texttt{\%\textbar{}\textbar{}\%\ \textless{}-\ function(a,\ b)\ if\ (!is.null(a))\ a\ else\ b}
  defines a custom operator, \texttt{\%\textbar{}\textbar{}\%}, borrowed
  from a common pattern in other languages. It returns \texttt{a} if
  \texttt{a} isn't \texttt{NULL}, and falls back to \texttt{b}
  otherwise. It's used later as a shorthand:
  \texttt{dv\_config\_controlled{[}{[}i{]}{]}\$scale\_factor\ \%\textbar{}\textbar{}\%\ 1}
  means ``use this DV's specified scale\_factor, or default to 1 if none
  was given'' --- this is exactly what lets you leave
  \texttt{scale\_factor\ =\ NULL} for any DV that doesn't need
  rescaling, without writing a separate \texttt{if}/\texttt{else} check
  for every single variable.
\item
  \texttt{dv\_config\_controlled} is a list of DVs to run the models
  over (equivalent to running multiple t-tests). In this particular
  case, most of my variables ranged from 1-7 but a select few were on
  0-100 scales. Estimating a model with variables on wildly different
  numeric scales can cause convergence problems in SEM software, since
  parameters for the larger-scale variables end up many times bigger
  than parameters for the smaller-scale ones. Dividing the 0--100
  variables by 10 brings them onto a comparable 0--10 range, closer to
  the other DVs, which helps the model estimate more reliably.
\item
  If every variable is already on a comparable scale, you can simply set
  \texttt{scale\_factor\ =\ NULL} for every entry (or omit it entirely
  and let \texttt{\%\textbar{}\textbar{}\%} default to 1 throughout).
\end{itemize}

\subsection{Step 2: Created scaled
variables}\label{step-2-created-scaled-variables}

\begin{Shaded}
\begin{Highlighting}[]
\CommentTok{\# ── Scale prejudice and norms variables from 0{-}100 to 0{-}10 ────────────────────}
\NormalTok{group\_comparisons }\OtherTok{\textless{}{-}}\NormalTok{ group\_comparisons }\SpecialCharTok{\%\textgreater{}\%}
  \FunctionTok{mutate}\NormalTok{(}
    \AttributeTok{prej\_jewish\_s        =}\NormalTok{ prej\_jewish }\SpecialCharTok{/} \DecValTok{10}\NormalTok{,}
    \AttributeTok{prej\_immigrant\_s     =}\NormalTok{ prej\_immigrant }\SpecialCharTok{/} \DecValTok{10}\NormalTok{,}
    \AttributeTok{prej\_muslim\_s        =}\NormalTok{ prej\_muslim }\SpecialCharTok{/} \DecValTok{10}\NormalTok{,}
    \AttributeTok{norms\_immigration\_s  =}\NormalTok{ norms\_immigration }\SpecialCharTok{/} \DecValTok{10}\NormalTok{,}
    \AttributeTok{norms\_racism\_s       =}\NormalTok{ norms\_racism }\SpecialCharTok{/} \DecValTok{10}\NormalTok{,}
    \AttributeTok{norms\_demo\_s         =}\NormalTok{ norms\_demo }\SpecialCharTok{/} \DecValTok{10}\NormalTok{)}
\end{Highlighting}
\end{Shaded}

\textbf{What this is doing:}

\begin{itemize}
\tightlist
\item
  Created the rescaled variables referred to in
  \texttt{dv\_config\_controlled}
\end{itemize}

\subsection{Step 3: Define auxiliary
variables}\label{step-3-define-auxiliary-variables}

\begin{Shaded}
\begin{Highlighting}[]
\NormalTok{COMMON\_AUXILIARY }\OtherTok{\textless{}{-}} \StringTok{"gender\_r"}
\end{Highlighting}
\end{Shaded}

\textbf{What this is doing:}

\begin{itemize}
\item
  FIML estimates model parameters using all available data, including
  partially missing cases.
\item
  An auxiliary variable is a variable that isn't part of your
  substantive hypothesis (it's not a DV or a theoretically important
  predictor), but is included in the model purely to improve the
  plausibility of that assumption and increase estimation precision. You
  want to include a variable that has very little missing data, thereby
  giving FIML more information to work with when handling the missing
  cases.
\end{itemize}

This code is written for \textbf{one} common auxiliary variable, applied
to every DV. You can adapt it either direction:

\begin{itemize}
\item
  \textbf{For zero auxiliary variables:} just set
  \texttt{COMMON\_AUXILIARY\ \textless{}-\ NULL}. The
  \texttt{use\_auxiliary} check inside the function (below) already
  handles this --- if \texttt{common\_auxiliary} is \texttt{NULL}, the
  auxiliary syntax is simply left blank, and the model runs without one.
\item
  \textbf{For multiple auxiliary variables:} you'd adjust the
  \texttt{aux\_syntax} construction to loop over a vector of auxiliary
  variables instead of a single one:
\end{itemize}

\begin{Shaded}
\begin{Highlighting}[]
\NormalTok{COMMON\_AUXILIARIES }\OtherTok{\textless{}{-}} \FunctionTok{c}\NormalTok{(}\StringTok{"gender\_r"}\NormalTok{, }\StringTok{"AGE"}\NormalTok{, }\StringTok{"EDUC"}\NormalTok{)}

\ControlFlowTok{if}\NormalTok{ (}\FunctionTok{length}\NormalTok{(common\_auxiliaries) }\SpecialCharTok{\textgreater{}} \DecValTok{0}\NormalTok{) \{}
\NormalTok{  aux\_lines }\OtherTok{\textless{}{-}} \FunctionTok{unlist}\NormalTok{(}\FunctionTok{lapply}\NormalTok{(common\_auxiliaries, }\ControlFlowTok{function}\NormalTok{(aux) \{}
    \FunctionTok{paste0}\NormalTok{(}
      \StringTok{"}\SpecialCharTok{\textbackslash{}n}\StringTok{  "}\NormalTok{, aux, }\StringTok{" \textasciitilde{} 1"}\NormalTok{,}
      \StringTok{"}\SpecialCharTok{\textbackslash{}n}\StringTok{  "}\NormalTok{, dv\_name, }\StringTok{" \textasciitilde{}\textasciitilde{} "}\NormalTok{, aux}
\NormalTok{    )}
\NormalTok{  \}))}
\NormalTok{  aux\_syntax }\OtherTok{\textless{}{-}} \FunctionTok{paste}\NormalTok{(aux\_lines, }\AttributeTok{collapse =} \StringTok{"}\SpecialCharTok{\textbackslash{}n}\StringTok{"}\NormalTok{)}
\NormalTok{\} }\ControlFlowTok{else}\NormalTok{ \{}
\NormalTok{  aux\_syntax }\OtherTok{\textless{}{-}} \StringTok{""}
\NormalTok{\}}
\end{Highlighting}
\end{Shaded}

\subsection{Step 4: Mean-centring
covariates}\label{step-4-mean-centring-covariates}

If you wanted to add controls/covariates to the model, you must
mean-centre continuous variables and ensure that dichotomous variables
are effect-coded.

\begin{Shaded}
\begin{Highlighting}[]
\NormalTok{group\_comparisons }\OtherTok{\textless{}{-}}\NormalTok{ group\_comparisons }\SpecialCharTok{\%\textgreater{}\%}
  \FunctionTok{mutate}\NormalTok{(}
    \AttributeTok{AGE\_c  =}\NormalTok{ AGE  }\SpecialCharTok{{-}} \FunctionTok{mean}\NormalTok{(AGE,  }\AttributeTok{na.rm =} \ConstantTok{TRUE}\NormalTok{),}
    \AttributeTok{EDUC\_c =}\NormalTok{ EDUC }\SpecialCharTok{{-}} \FunctionTok{mean}\NormalTok{(EDUC, }\AttributeTok{na.rm =} \ConstantTok{TRUE}\NormalTok{),}
    \AttributeTok{CONS\_c =}\NormalTok{ CONS }\SpecialCharTok{{-}} \FunctionTok{mean}\NormalTok{(CONS, }\AttributeTok{na.rm =} \ConstantTok{TRUE}\NormalTok{)}
\NormalTok{  )}

\NormalTok{group\_comparisons }\OtherTok{\textless{}{-}}\NormalTok{ group\_comparisons }\SpecialCharTok{\%\textgreater{}\%}
  \FunctionTok{mutate}\NormalTok{(}\AttributeTok{VIC =} \FunctionTok{case\_when}\NormalTok{(}
\NormalTok{    STATE\_2 }\SpecialCharTok{==} \DecValTok{0} \SpecialCharTok{\textasciitilde{}} \SpecialCharTok{{-}}\DecValTok{1}\NormalTok{,}
    \ConstantTok{TRUE} \SpecialCharTok{\textasciitilde{}} \DecValTok{1}
\NormalTok{  ))}
\end{Highlighting}
\end{Shaded}

\textbf{What this is doing:}

\begin{itemize}
\item
  \texttt{AGE\_c}, \texttt{EDUC\_c}, and \texttt{CONS\_c} are
  mean-centered versions of the covariates (subtracting each variable's
  own mean from every value), which makes the model's intercept
  interpretable as the predicted value at the \emph{average} level of
  these covariates, rather than at a covariate value of zero (which is
  often meaningless --- e.g.~an age of 0).
\item
  \texttt{VIC} is created using effect coding (covered earlier in this
  chapter) rather than standard 0/1 dummy coding, so this covariate's
  coefficient reflects deviation from the grand mean rather than from an
  arbitrary reference group.
\end{itemize}

\subsection{Step 5: The analysis
function}\label{step-5-the-analysis-function}

\begin{Shaded}
\begin{Highlighting}[]
\NormalTok{run\_group\_comparison\_controlled }\OtherTok{\textless{}{-}} \ControlFlowTok{function}\NormalTok{(dv\_name, }\AttributeTok{common\_auxiliary =} \ConstantTok{NULL}\NormalTok{, data,}
                                            \AttributeTok{scale\_factor =} \DecValTok{1}\NormalTok{) \{}
  
  \CommentTok{\# Build auxiliary syntax — a single common auxiliary variable used for every DV}
  \CommentTok{\# (skip it if the DV itself happens to be the common auxiliary)}
\NormalTok{  use\_auxiliary }\OtherTok{\textless{}{-}} \SpecialCharTok{!}\FunctionTok{is.null}\NormalTok{(common\_auxiliary) }\SpecialCharTok{\&\&}\NormalTok{ common\_auxiliary }\SpecialCharTok{!=}\NormalTok{ dv\_name}
  
  \ControlFlowTok{if}\NormalTok{ (use\_auxiliary) \{}
\NormalTok{    aux\_syntax }\OtherTok{\textless{}{-}} \FunctionTok{paste0}\NormalTok{(}
      \StringTok{"}\SpecialCharTok{\textbackslash{}n}\StringTok{  \# Auxiliary variable {-} estimate its mean}\SpecialCharTok{\textbackslash{}n}\StringTok{  "}\NormalTok{,}
\NormalTok{      common\_auxiliary, }\StringTok{" \textasciitilde{} 1"}\NormalTok{,}
      \StringTok{"}\SpecialCharTok{\textbackslash{}n\textbackslash{}n}\StringTok{  \# Covariance between DV and auxiliary}\SpecialCharTok{\textbackslash{}n}\StringTok{  "}\NormalTok{,}
\NormalTok{      dv\_name, }\StringTok{" \textasciitilde{}\textasciitilde{} "}\NormalTok{, common\_auxiliary}
\NormalTok{    )}
\NormalTok{  \} }\ControlFlowTok{else}\NormalTok{ \{}
\NormalTok{    aux\_syntax }\OtherTok{\textless{}{-}} \StringTok{""}
\NormalTok{  \}}
  
  \CommentTok{\# Unconstrained model}
\NormalTok{  model\_unconstrained }\OtherTok{\textless{}{-}} \FunctionTok{paste0}\NormalTok{(}\StringTok{"}
\StringTok{  "}\NormalTok{, dv\_name, }\StringTok{" \textasciitilde{} c(mu1, mu2)*1}
\StringTok{  "}\NormalTok{, dv\_name, }\StringTok{" \textasciitilde{} AGE\_c + EDUC\_c + CONS\_c + VIC}
\StringTok{  "}\NormalTok{, dv\_name, }\StringTok{" \textasciitilde{}\textasciitilde{} c(v1, v2)*"}\NormalTok{, dv\_name, aux\_syntax, }\StringTok{"}

\StringTok{  diff     := mu2 {-} mu1}
\StringTok{  sd\_welch := sqrt((v1 + v2) / 2)}
\StringTok{  d\_welch  := diff / sd\_welch}
\StringTok{"}\NormalTok{)}
  
  \CommentTok{\# Constrained model (means equal across groups)}
\NormalTok{  model\_constrained }\OtherTok{\textless{}{-}} \FunctionTok{paste0}\NormalTok{(}\StringTok{"}
\StringTok{  "}\NormalTok{, dv\_name, }\StringTok{" \textasciitilde{} c(mu\_equal, mu\_equal)*1}
\StringTok{  "}\NormalTok{, dv\_name, }\StringTok{" \textasciitilde{} AGE\_c + EDUC\_c + CONS\_c + VIC}
\StringTok{  "}\NormalTok{, dv\_name, }\StringTok{" \textasciitilde{}\textasciitilde{} c(v1, v2)*"}\NormalTok{, dv\_name, aux\_syntax, }\StringTok{"}
\StringTok{"}\NormalTok{)}
  
\NormalTok{  fit\_unconstrained }\OtherTok{\textless{}{-}} \FunctionTok{sem}\NormalTok{(}
\NormalTok{    model\_unconstrained,}
    \AttributeTok{data          =}\NormalTok{ data,}
    \AttributeTok{group         =} \StringTok{"pre\_post"}\NormalTok{,}
    \AttributeTok{meanstructure =} \ConstantTok{TRUE}\NormalTok{,}
    \AttributeTok{fixed.x       =} \ConstantTok{FALSE}\NormalTok{,}
    \AttributeTok{missing       =} \StringTok{"FIML"}
\NormalTok{  )}
  
\NormalTok{  fit\_constrained }\OtherTok{\textless{}{-}} \FunctionTok{sem}\NormalTok{(}
\NormalTok{    model\_constrained,}
    \AttributeTok{data          =}\NormalTok{ data,}
    \AttributeTok{group         =} \StringTok{"pre\_post"}\NormalTok{,}
    \AttributeTok{meanstructure =} \ConstantTok{TRUE}\NormalTok{,}
    \AttributeTok{fixed.x       =} \ConstantTok{FALSE}\NormalTok{,}
    \AttributeTok{missing       =} \StringTok{"FIML"}
\NormalTok{  )}
  
  \CommentTok{\# Verify group ordering (pre = 1, post = 2)}
\NormalTok{  group\_labels }\OtherTok{\textless{}{-}} \FunctionTok{lavInspect}\NormalTok{(fit\_unconstrained, }\StringTok{"group.label"}\NormalTok{)}
  \ControlFlowTok{if}\NormalTok{ (}\SpecialCharTok{!}\FunctionTok{identical}\NormalTok{(group\_labels, }\FunctionTok{c}\NormalTok{(}\StringTok{"pre"}\NormalTok{, }\StringTok{"post"}\NormalTok{))) \{}
    \FunctionTok{stop}\NormalTok{(}
      \StringTok{"Unexpected group ordering for \textquotesingle{}"}\NormalTok{, dv\_name, }\StringTok{"\textquotesingle{}: "}\NormalTok{,}
      \FunctionTok{paste}\NormalTok{(group\_labels, }\AttributeTok{collapse =} \StringTok{", "}\NormalTok{),}
      \StringTok{". Expected c(\textquotesingle{}pre\textquotesingle{}, \textquotesingle{}post\textquotesingle{})."}
\NormalTok{    )}
\NormalTok{  \}}
  
  \CommentTok{\# Sample sizes}
\NormalTok{  N\_total }\OtherTok{\textless{}{-}} \FunctionTok{lavInspect}\NormalTok{(fit\_unconstrained, }\StringTok{"nobs"}\NormalTok{)}
\NormalTok{  n1      }\OtherTok{\textless{}{-}}\NormalTok{ N\_total[}\DecValTok{1}\NormalTok{]}
\NormalTok{  n2      }\OtherTok{\textless{}{-}}\NormalTok{ N\_total[}\DecValTok{2}\NormalTok{]}
  
  \CommentTok{\# Likelihood ratio test}
\NormalTok{  lr\_test }\OtherTok{\textless{}{-}} \FunctionTok{anova}\NormalTok{(fit\_constrained, fit\_unconstrained)}
\NormalTok{  chi\_sq  }\OtherTok{\textless{}{-}}\NormalTok{ lr\_test}\SpecialCharTok{$}\StringTok{\textasciigrave{}}\AttributeTok{Chisq diff}\StringTok{\textasciigrave{}}\NormalTok{[}\DecValTok{2}\NormalTok{]}
\NormalTok{  df      }\OtherTok{\textless{}{-}}\NormalTok{ lr\_test}\SpecialCharTok{$}\StringTok{\textasciigrave{}}\AttributeTok{Df diff}\StringTok{\textasciigrave{}}\NormalTok{[}\DecValTok{2}\NormalTok{]}
\NormalTok{  p\_raw   }\OtherTok{\textless{}{-}}\NormalTok{ lr\_test}\SpecialCharTok{$}\StringTok{\textasciigrave{}}\AttributeTok{Pr(\textgreater{}Chisq)}\StringTok{\textasciigrave{}}\NormalTok{[}\DecValTok{2}\NormalTok{]}
  
  \CommentTok{\# Extract parameters}
\NormalTok{  params    }\OtherTok{\textless{}{-}} \FunctionTok{parameterEstimates}\NormalTok{(fit\_unconstrained, }\AttributeTok{ci =} \ConstantTok{TRUE}\NormalTok{)}
\NormalTok{  means     }\OtherTok{\textless{}{-}} \FunctionTok{subset}\NormalTok{(params, op }\SpecialCharTok{==} \StringTok{"\textasciitilde{}1"} \SpecialCharTok{\&}\NormalTok{ lhs }\SpecialCharTok{==}\NormalTok{ dv\_name)}
\NormalTok{  diff\_row  }\OtherTok{\textless{}{-}} \FunctionTok{subset}\NormalTok{(params, label }\SpecialCharTok{==} \StringTok{"diff"}\NormalTok{)}
\NormalTok{  d\_row     }\OtherTok{\textless{}{-}} \FunctionTok{subset}\NormalTok{(params, label }\SpecialCharTok{==} \StringTok{"d\_welch"}\NormalTok{)}
\NormalTok{  variances }\OtherTok{\textless{}{-}} \FunctionTok{subset}\NormalTok{(params, op }\SpecialCharTok{==} \StringTok{"\textasciitilde{}\textasciitilde{}"} \SpecialCharTok{\&}\NormalTok{ lhs }\SpecialCharTok{==}\NormalTok{ dv\_name }\SpecialCharTok{\&}\NormalTok{ rhs }\SpecialCharTok{==}\NormalTok{ dv\_name)}
  
  \CommentTok{\# Back{-}transform scale{-}dependent outputs; Cohen\textquotesingle{}s d is dimensionless so no transform}
  \FunctionTok{data.frame}\NormalTok{(}
    \AttributeTok{variable =} \FunctionTok{sub}\NormalTok{(}\StringTok{"\_s$"}\NormalTok{, }\StringTok{""}\NormalTok{, dv\_name),   }\CommentTok{\# strip \_s suffix for display}
    
    \AttributeTok{pre\_M   =} \FunctionTok{round}\NormalTok{(means}\SpecialCharTok{$}\NormalTok{est[means}\SpecialCharTok{$}\NormalTok{group }\SpecialCharTok{==} \DecValTok{1}\NormalTok{]         }\SpecialCharTok{*}\NormalTok{ scale\_factor, }\DecValTok{2}\NormalTok{),}
    \AttributeTok{pre\_var =} \FunctionTok{round}\NormalTok{(variances}\SpecialCharTok{$}\NormalTok{est[variances}\SpecialCharTok{$}\NormalTok{group }\SpecialCharTok{==} \DecValTok{1}\NormalTok{] }\SpecialCharTok{*}\NormalTok{ scale\_factor}\SpecialCharTok{\^{}}\DecValTok{2}\NormalTok{, }\DecValTok{2}\NormalTok{),}
    
    \AttributeTok{post\_M   =} \FunctionTok{round}\NormalTok{(means}\SpecialCharTok{$}\NormalTok{est[means}\SpecialCharTok{$}\NormalTok{group }\SpecialCharTok{==} \DecValTok{2}\NormalTok{]         }\SpecialCharTok{*}\NormalTok{ scale\_factor, }\DecValTok{2}\NormalTok{),}
    \AttributeTok{post\_var =} \FunctionTok{round}\NormalTok{(variances}\SpecialCharTok{$}\NormalTok{est[variances}\SpecialCharTok{$}\NormalTok{group }\SpecialCharTok{==} \DecValTok{2}\NormalTok{] }\SpecialCharTok{*}\NormalTok{ scale\_factor}\SpecialCharTok{\^{}}\DecValTok{2}\NormalTok{, }\DecValTok{2}\NormalTok{),}
    
    \AttributeTok{diff\_est =} \FunctionTok{round}\NormalTok{(diff\_row}\SpecialCharTok{$}\NormalTok{est      }\SpecialCharTok{*}\NormalTok{ scale\_factor, }\DecValTok{2}\NormalTok{),}
    \AttributeTok{diff\_lci =} \FunctionTok{round}\NormalTok{(diff\_row}\SpecialCharTok{$}\NormalTok{ci.lower }\SpecialCharTok{*}\NormalTok{ scale\_factor, }\DecValTok{2}\NormalTok{),}
    \AttributeTok{diff\_uci =} \FunctionTok{round}\NormalTok{(diff\_row}\SpecialCharTok{$}\NormalTok{ci.upper }\SpecialCharTok{*}\NormalTok{ scale\_factor, }\DecValTok{2}\NormalTok{),}
    
    \AttributeTok{d\_est =} \FunctionTok{round}\NormalTok{(d\_row}\SpecialCharTok{$}\NormalTok{est,      }\DecValTok{2}\NormalTok{),}
    \AttributeTok{d\_lci =} \FunctionTok{round}\NormalTok{(d\_row}\SpecialCharTok{$}\NormalTok{ci.lower, }\DecValTok{2}\NormalTok{),}
    \AttributeTok{d\_uci =} \FunctionTok{round}\NormalTok{(d\_row}\SpecialCharTok{$}\NormalTok{ci.upper, }\DecValTok{2}\NormalTok{),}
    
    \AttributeTok{chi\_sq  =} \FunctionTok{round}\NormalTok{(chi\_sq, }\DecValTok{2}\NormalTok{),}
    \AttributeTok{df      =}\NormalTok{ df,}
    \AttributeTok{p\_raw   =}\NormalTok{ p\_raw,}
    
    \AttributeTok{n\_pre   =}\NormalTok{ n1,}
    \AttributeTok{n\_post  =}\NormalTok{ n2,}
    \AttributeTok{n\_total =} \FunctionTok{sum}\NormalTok{(N\_total),}
    
    \AttributeTok{stringsAsFactors =} \ConstantTok{FALSE}
\NormalTok{  )}
\NormalTok{\}}
\end{Highlighting}
\end{Shaded}

\textbf{What this is doing:}

\begin{itemize}
\tightlist
\item
  This is a custom function, written once, then looped over every DV in
  the next step.
\end{itemize}

\begin{itemize}
\tightlist
\item
  \texttt{use\_auxiliary\ \textless{}-\ !is.null(common\_auxiliary)\ \&\&\ common\_auxiliary\ !=\ dv\_name}
  checks two things: that an auxiliary variable was actually specified,
  and that it isn't the same variable as the DV currently being tested
  (which would be nonsensical). If both are true, \texttt{aux\_syntax}
  is built (the auxiliary's mean is estimated, and its covariance with
  the DV specified) --- this is the lavaan syntax that actually
  implements the FIML auxiliary-variable approach.
\end{itemize}

\begin{itemize}
\item
  \texttt{model\_unconstrained} is the \textbf{alternative hypothesis}
  model: \texttt{c(mu1,\ mu2)*1} estimates a \textbf{separate mean} for
  each of the two groups (labeled \texttt{mu1} and \texttt{mu2}) ---
  this is the ``the groups genuinely differ'' model.

  \begin{itemize}
  \item
    \texttt{\textasciitilde{}\textasciitilde{}\ c(v1,\ v2)*dv\_name}
    estimates a \textbf{separate variance} per group too, which is what
    makes this approximate a \textbf{Welch t-test} (which doesn't assume
    equal variances) rather than a standard Student's t-test.
  \item
    The \texttt{diff}, \texttt{sd\_welch}, and \texttt{d\_welch} lines
    are lavaan's \texttt{:=} syntax for defining new quantities
    \emph{derived} from the model's parameters --- here, calculating the
    mean difference and a Welch-style Cohen's d directly as part of the
    model, rather than needing a separate calculation step afterward.
  \end{itemize}
\end{itemize}

\begin{itemize}
\tightlist
\item
  \texttt{model\_constrained} is your \textbf{null hypothesis} model:
  \texttt{c(mu\_equal,\ mu\_equal)*1} forces both groups to share the
  \textbf{same} mean --- this is ``the groups don't actually differ''
  model, everything else kept identical.
\end{itemize}

\begin{itemize}
\tightlist
\item
  \texttt{sem(...,\ missing\ =\ "FIML")} is what actually runs each
  model using full information maximum likelihood --- this single
  argument is what allows participants with partial missing data to
  still contribute to parameter estimation, rather than being dropped
  via listwise deletion, addressing the exact problem described in your
  intro paragraph.
\end{itemize}

\begin{itemize}
\item
  \texttt{group\ =\ "pre\_post"} tells lavaan which variable defines
  your two comparison groups --- the model is fit \textbf{separately per
  group}, which is what allows the means (and variances) to differ
  between them in the unconstrained model.
\item
  The \texttt{group\_labels} check is a safeguard: lavaan orders groups
  alphabetically/by factor level internally, and if that order ever
  didn't match your expectation (\texttt{"pre"} first, \texttt{"post"}
  second), any downstream code assuming \texttt{group\ ==\ 1} is ``pre''
  would silently mislabel results. \texttt{stop()} halts execution with
  a clear error message the moment this happens, rather than letting a
  silent labeling mistake flow through to your final table.
\end{itemize}

\begin{itemize}
\tightlist
\item
  \texttt{lavInspect(...,\ "nobs")} extracts the sample size used in
  each group.
\end{itemize}

\begin{itemize}
\item
  \texttt{anova(fit\_constrained,\ fit\_unconstrained)} runs a
  \textbf{likelihood ratio test}, comparing how well the two models fit
  the data. This is the ``t-test equivalent'': rather than a
  t-statistic, you get a chi-square difference test, where a significant
  result means the unconstrained model (means allowed to differ) fits
  meaningfully better than the constrained one (means forced equal) ---
  i.e.~the groups differ significantly.
\item
  The \texttt{params} block tpulls every result out of the fitted model
  --- group means, variances, the mean difference with its confidence
  interval, and Cohen's d with its confidence interval --- and assembles
  them into a single tidy row of output for this one DV.

  \begin{itemize}
  \item
    This is where \textbf{\texttt{scale\_factor}} does its work:
    \texttt{pre\_M\ *\ scale\_factor},
    \texttt{post\_M\ *\ scale\_factor}, and the difference/CI values are
    all multiplied back up, undoing the /10 rescaling, so the variables
    are reported back on their original, interpretable 0--100 scale.
  \item
    \texttt{d\_est}, \texttt{d\_lci}, \texttt{d\_uci} are \textbf{not}
    multiplied by \texttt{scale\_factor} --- Cohen's d is a standardised
    effect size (already expressed in SD units), so it's scale-invariant
    by design and doesn't need back-transforming.
  \item
    \texttt{sub("\_s\$",\ "",\ dv\_name)} also strips the \texttt{\_s}
    suffix back off the variable name for the final table, so it reads
    as \texttt{"prej\_immigrant"} rather than the working name
    \texttt{"prej\_immigrant\_s"}.
  \end{itemize}
\end{itemize}

\subsection{Step 6: Running every comparison in a
loop}\label{step-6-running-every-comparison-in-a-loop}

\begin{Shaded}
\begin{Highlighting}[]
\NormalTok{results\_list\_controlled }\OtherTok{\textless{}{-}} \FunctionTok{vector}\NormalTok{(}\StringTok{"list"}\NormalTok{, }\FunctionTok{length}\NormalTok{(dv\_config\_controlled))}
\ControlFlowTok{for}\NormalTok{ (i }\ControlFlowTok{in} \FunctionTok{seq\_along}\NormalTok{(dv\_config\_controlled)) \{}
  \FunctionTok{cat}\NormalTok{(}\StringTok{"}\SpecialCharTok{\textbackslash{}n}\StringTok{\#\#\#\# Comparison"}\NormalTok{, i, }\StringTok{"of"}\NormalTok{, }\FunctionTok{length}\NormalTok{(dv\_config\_controlled), }\StringTok{":"}\NormalTok{,}
\NormalTok{      dv\_config\_controlled[[i]]}\SpecialCharTok{$}\NormalTok{dv, }\StringTok{"\#\#\#\#}\SpecialCharTok{\textbackslash{}n}\StringTok{"}\NormalTok{)}
\NormalTok{  results\_list\_controlled[[i]] }\OtherTok{\textless{}{-}} \FunctionTok{run\_group\_comparison\_controlled}\NormalTok{(}
    \AttributeTok{dv\_name           =}\NormalTok{ dv\_config\_controlled[[i]]}\SpecialCharTok{$}\NormalTok{dv,}
    \AttributeTok{common\_auxiliary  =}\NormalTok{ COMMON\_AUXILIARY,}
    \AttributeTok{scale\_factor      =}\NormalTok{ dv\_config\_controlled[[i]]}\SpecialCharTok{$}\NormalTok{scale\_factor }\SpecialCharTok{\%||\%} \DecValTok{1}\NormalTok{,}
    \AttributeTok{data              =}\NormalTok{ group\_comparisons}
\NormalTok{  )}
\NormalTok{\}}

\NormalTok{final\_results\_controlled }\OtherTok{\textless{}{-}} \FunctionTok{do.call}\NormalTok{(rbind, results\_list\_controlled)}
\FunctionTok{rownames}\NormalTok{(final\_results\_controlled) }\OtherTok{\textless{}{-}} \ConstantTok{NULL}
\end{Highlighting}
\end{Shaded}

\textbf{What this is doing:}

\begin{itemize}
\item
  \texttt{vector("list",\ length(...))} pre-allocates an empty list of
  the right size --- more efficient than growing a list one element at a
  time in a loop.
\item
  The \texttt{for} loop then runs
  \texttt{run\_group\_comparison\_controlled()} once per DV listed in
  \texttt{dv\_config\_controlled}, printing a progress message
  (\texttt{cat(...)}) each time so you can track how far through a long
  list of comparisons you are.
\item
  The
  \texttt{scale\_factor\ =\ dv\_config\_controlled{[}{[}i{]}{]}\$scale\_factor\ \%\textbar{}\textbar{}\%\ 1}
  line is exactly where the Step 2 operator gets used in practice ---
  pulling each DV's own specified scale\_factor, or defaulting to 1 if
  it was left \texttt{NULL}.
\item
  \texttt{do.call(rbind,\ results\_list\_controlled)} stacks every DV's
  one-row result into a single combined dataframe, and
  \texttt{rownames(...)\ \textless{}-\ NULL} tidies up the row
  numbering.
\end{itemize}

\subsection{Step 7: Correcting for multiple
comparisons}\label{step-7-correcting-for-multiple-comparisons}

\begin{Shaded}
\begin{Highlighting}[]
\NormalTok{final\_results\_controlled}\SpecialCharTok{$}\NormalTok{p\_adj\_BH         }\OtherTok{\textless{}{-}} \FunctionTok{p.adjust}\NormalTok{(final\_results\_controlled}\SpecialCharTok{$}\NormalTok{p\_raw, }\AttributeTok{method =} \StringTok{"BH"}\NormalTok{)}
\NormalTok{final\_results\_controlled}\SpecialCharTok{$}\NormalTok{p\_adj\_bonferroni }\OtherTok{\textless{}{-}} \FunctionTok{p.adjust}\NormalTok{(final\_results\_controlled}\SpecialCharTok{$}\NormalTok{p\_raw, }\AttributeTok{method =} \StringTok{"bonferroni"}\NormalTok{)}
\end{Highlighting}
\end{Shaded}

\textbf{What this is doing:}

\begin{itemize}
\item
  \texttt{p.adjust()} is base R's general-purpose function for
  correcting a whole vector of p-values at once.
\item
  Two different methods have been applied: Bonferroni and BH (which is
  the false discovery rate method).
\end{itemize}

\subsection{Step 8: Formatting and exporting the
table}\label{step-8-formatting-and-exporting-the-table}

\begin{Shaded}
\begin{Highlighting}[]
\NormalTok{format\_results\_table\_controlled }\OtherTok{\textless{}{-}} \ControlFlowTok{function}\NormalTok{(results) \{}
\NormalTok{  format\_mean\_sd }\OtherTok{\textless{}{-}} \ControlFlowTok{function}\NormalTok{(mean, var) }\FunctionTok{sprintf}\NormalTok{(}\StringTok{"\%.2f (\%.2f)"}\NormalTok{, mean, }\FunctionTok{sqrt}\NormalTok{(var))}
\NormalTok{  format\_mean\_ci }\OtherTok{\textless{}{-}} \ControlFlowTok{function}\NormalTok{(mean, lower, upper) }\FunctionTok{sprintf}\NormalTok{(}\StringTok{"\%.2f [\%.2f, \%.2f]"}\NormalTok{, mean, lower, upper)}
  
  \FunctionTok{data.frame}\NormalTok{(}
    \AttributeTok{Variable          =}\NormalTok{ results}\SpecialCharTok{$}\NormalTok{variable,}
    \StringTok{\textasciigrave{}}\AttributeTok{Pre M (SD)}\StringTok{\textasciigrave{}}      \OtherTok{=} \FunctionTok{mapply}\NormalTok{(format\_mean\_sd, results}\SpecialCharTok{$}\NormalTok{pre\_M,  results}\SpecialCharTok{$}\NormalTok{pre\_var),}
    \StringTok{\textasciigrave{}}\AttributeTok{Post M (SD)}\StringTok{\textasciigrave{}}     \OtherTok{=} \FunctionTok{mapply}\NormalTok{(format\_mean\_sd, results}\SpecialCharTok{$}\NormalTok{post\_M, results}\SpecialCharTok{$}\NormalTok{post\_var),}
    \StringTok{\textasciigrave{}}\AttributeTok{Mean difference}\StringTok{\textasciigrave{}} \OtherTok{=} \FunctionTok{sprintf}\NormalTok{(}\StringTok{"\%.2f"}\NormalTok{, results}\SpecialCharTok{$}\NormalTok{diff\_est),}
    \StringTok{\textasciigrave{}}\AttributeTok{Cohen\textquotesingle{}s d [95\% CI]}\StringTok{\textasciigrave{}} \OtherTok{=} \FunctionTok{mapply}\NormalTok{(format\_mean\_ci, results}\SpecialCharTok{$}\NormalTok{d\_est, results}\SpecialCharTok{$}\NormalTok{d\_lci, results}\SpecialCharTok{$}\NormalTok{d\_uci),}
    \StringTok{\textasciigrave{}}\AttributeTok{Chi{-}square}\StringTok{\textasciigrave{}}      \OtherTok{=} \FunctionTok{sprintf}\NormalTok{(}\StringTok{"\%.2f"}\NormalTok{, results}\SpecialCharTok{$}\NormalTok{chi\_sq),}
    \AttributeTok{p                 =}\NormalTok{ results}\SpecialCharTok{$}\NormalTok{p\_raw,}
    \StringTok{\textasciigrave{}}\AttributeTok{p (FDR)}\StringTok{\textasciigrave{}}         \OtherTok{=}\NormalTok{ results}\SpecialCharTok{$}\NormalTok{p\_adj\_BH,}
    \StringTok{\textasciigrave{}}\AttributeTok{p (Bonf)}\StringTok{\textasciigrave{}}        \OtherTok{=}\NormalTok{ results}\SpecialCharTok{$}\NormalTok{p\_adj\_bonferroni,}
    \AttributeTok{check.names =} \ConstantTok{FALSE}\NormalTok{, }\AttributeTok{stringsAsFactors =} \ConstantTok{FALSE}
\NormalTok{  )}
\NormalTok{\}}
\end{Highlighting}
\end{Shaded}

\textbf{What this is doing:}

\begin{itemize}
\item
  This function turns the raw numeric results into publication-ready
  formatted text --- e.g.~combining a mean and its variance into a
  single \texttt{"3.45\ (1.02)"} style string (converting variance to SD
  first, via \texttt{sqrt(var)}, since SD rather than variance is the
  conventional format for reporting), and combining Cohen's d with its
  confidence interval into \texttt{"0.32\ {[}0.18,\ 0.46{]}"} style APA
  formatting.
\item
  \texttt{check.names\ =\ FALSE} stops R from automatically mangling
  your deliberately spaced/punctuated column names (like
  \texttt{"Pre\ M\ (SD)"}) into syntactically ``safe'' but ugly names
  (like \texttt{"Pre.M..SD."}).
\end{itemize}

\begin{Shaded}
\begin{Highlighting}[]
\NormalTok{nice\_formatted\_table\_controlled }\OtherTok{\textless{}{-}} \FunctionTok{nice\_table}\NormalTok{(}
\NormalTok{  formatted\_table\_controlled,}
  \AttributeTok{col.format.p =} \DecValTok{7}\SpecialCharTok{:}\DecValTok{9}\NormalTok{,}
  \AttributeTok{spacing      =} \DecValTok{1}\NormalTok{,}
  \AttributeTok{title        =} \FunctionTok{c}\NormalTok{(}\StringTok{"Table S3"}\NormalTok{, }\StringTok{"Pre{-}post comparisons controlling for age, education, and conservatism, with raw, BH{-}corrected, and Bonferroni{-}corrected p{-}values."}\NormalTok{)}
\NormalTok{)}

\NormalTok{sect\_properties }\OtherTok{\textless{}{-}} \FunctionTok{prop\_section}\NormalTok{(}
  \AttributeTok{page\_size    =} \FunctionTok{page\_size}\NormalTok{(}\AttributeTok{orient =} \StringTok{"landscape"}\NormalTok{),}
  \AttributeTok{page\_margins =} \FunctionTok{page\_mar}\NormalTok{(}\AttributeTok{top =} \FloatTok{0.5}\NormalTok{, }\AttributeTok{bottom =} \FloatTok{0.5}\NormalTok{, }\AttributeTok{left =} \FloatTok{0.5}\NormalTok{, }\AttributeTok{right =} \FloatTok{0.5}\NormalTok{),}
  \AttributeTok{type         =} \StringTok{"continuous"}
\NormalTok{)}

\FunctionTok{save\_as\_docx}\NormalTok{(}
\NormalTok{  nice\_formatted\_table\_controlled,}
  \AttributeTok{path       =} \StringTok{"Tables/Table 3 Pre{-}post comparisons AGE EDUC CONS STATE\_2 controlled.docx"}\NormalTok{,}
  \AttributeTok{pr\_section =}\NormalTok{ sect\_properties}
\NormalTok{)}
\end{Highlighting}
\end{Shaded}

\textbf{What this is doing:}

\begin{itemize}
\item
  Using \texttt{nice\_table()} from the \textbf{rempsyc} package, APA
  styling has been applied to the finished dataframe.
\item
  \texttt{col.format.p\ =\ 7:9} tells it specifically which columns (7
  through 9 --- the raw, FDR, and Bonferroni p-value columns) should be
  formatted using APA's p-value conventions
  (e.g.~\texttt{p\ \textless{}\ .001} rather than a long decimal).
\item
  \texttt{title} adds a table number and caption directly into the
  exported table.
\item
  \texttt{prop\_section()} (from the \textbf{officer} package) defines a
  landscape page layout with tight margins, needed here because this
  table is unusually wide (nine columns).
\item
  \texttt{save\_as\_docx()} is \textbf{rempsyc's} own direct Word export
  function.
\end{itemize}

\chapter{Regression}\label{regression}

Page under construction\ldots{} 🚧

\section{Univariate and multivariate
regression}\label{univariate-and-multivariate-regression}

TBD.

\subsection{Building the models}\label{building-the-models-1}

TBD.

\subsection{Visualising the models}\label{visualising-the-models}

TBD.

\section{Binary logistic regression}\label{binary-logistic-regression}

TBD.

\chapter{Analysis of Variance}\label{analysis-of-variance}

This page is under construction\ldots{} 🚧

\chapter{Structural Equation
Modelling}\label{structural-equation-modelling}

\section{Covariance-based SEM}\label{covariance-based-sem}

The code below is taken from Pittaway et al. (2024). This was a
complicated model that included two observed IVs, three observed DVs,
and a parallel mediation model comprising one observed and one latent
variable. In order to calculate the specific indirect and total effects,
you must specify names for each indirect effect and assign them an
alphabetical code.

\emph{Note:} I am pasting this code here so that others can access it as
needed, but I will return to tidy up this page and annotate the code
more comprehensively when I have time. If you are trying to run a
similar model and it's throwing errors, let me know and I can help
troubleshoot.

\subsection{Specify the model}\label{specify-the-model}

\begin{Shaded}
\begin{Highlighting}[]
\FunctionTok{library}\NormalTok{(lavaan)}
\FunctionTok{library}\NormalTok{(semTools)}

\DocumentationTok{\#\#\# M1 }\AlertTok{\#\#\#}
\NormalTok{m1 }\OtherTok{\textless{}{-}} \StringTok{\textquotesingle{}}
\StringTok{\# LATENT VARIABLES}
\StringTok{ANX =\textasciitilde{} anx\_aff + anx\_behav + anx\_rumin + anx\_pers}

\StringTok{\# STRUCTURAL MODEL (Direct Effects)}
\StringTok{conv\_private \textasciitilde{} e*ANX + a*ecas + l*future + o*immediate}
\StringTok{conv\_public \textasciitilde{} g*ANX + c*ecas + m*future + p*immediate}
\StringTok{rad\_public \textasciitilde{} h*ANX + d*ecas + n*future + q*immediate}
\StringTok{ecas \textasciitilde{} b*future + j*immediate}
\StringTok{ANX \textasciitilde{} f*future + k*immediate + i*ecas}

\StringTok{\# COVARIANCES}
\StringTok{conv\_private \textasciitilde{}\textasciitilde{} conv\_public + rad\_public}
\StringTok{conv\_public \textasciitilde{}\textasciitilde{} rad\_public}

\StringTok{\# INDIRECT AND TOTAL EFFECTS}

\StringTok{\#\# CFC {-}\textgreater{} PRIV}
\StringTok{\#Indirect effect 1: cfc {-}\textgreater{} aeca {-}\textgreater{} priv}
\StringTok{\#Indirect effect 2: cfc {-}\textgreater{} aeca {-}\textgreater{} anx {-}\textgreater{} priv}
\StringTok{\#Indirect effect 3: cfc {-}\textgreater{} anx {-}\textgreater{} priv}

\StringTok{de\_cfc\_1 := l}
\StringTok{ie\_cfc\_1 := b*a}
\StringTok{ie\_cfc\_2 := b*i*e}
\StringTok{ie\_cfc\_3 := f*e}
\StringTok{TI\_cfc\_1 := ie\_cfc\_1 + ie\_cfc\_2 + ie\_cfc\_3}
\StringTok{TE\_cfc\_1 := l + TI\_cfc\_1}


\StringTok{\#\# CFC {-}\textgreater{} PUB}
\StringTok{\#Indirect effect 4: cfc {-}\textgreater{} aeca {-}\textgreater{} pub}
\StringTok{\#Indirect effect 5: cfc {-}\textgreater{} aeca {-}\textgreater{} anx {-}\textgreater{} pub}
\StringTok{\#Indirect effect 6: cfc {-}\textgreater{} anx {-}\textgreater{} pub}

\StringTok{de\_cfc\_2 := m}
\StringTok{ie\_cfc\_4 := b*c}
\StringTok{ie\_cfc\_5 := b*i*g}
\StringTok{ie\_cfc\_6 := f*g}
\StringTok{TI\_cfc\_2 := ie\_cfc\_4 + ie\_cfc\_5 + ie\_cfc\_6}
\StringTok{TE\_cfc\_2 := m + TI\_cfc\_2}


\StringTok{\#\# CFC {-}\textgreater{} RAD}
\StringTok{\#Indirect effect 7: cfc {-}\textgreater{} aeca {-}\textgreater{} rad}
\StringTok{\#Indirect effect 8: cfc {-}\textgreater{} aeca {-}\textgreater{} anx {-}\textgreater{} rad}
\StringTok{\#Indirect effect 9: cfc {-}\textgreater{} anx {-}\textgreater{} rad}

\StringTok{de\_cfc\_3 := n}
\StringTok{ie\_cfc\_7 := b*d}
\StringTok{ie\_cfc\_8 := b*i*h}
\StringTok{ie\_cfc\_9 := f*h}
\StringTok{TI\_cfc\_3 := ie\_cfc\_7 + ie\_cfc\_8 + ie\_cfc\_9}
\StringTok{TE\_cfc\_3 := n + TI\_cfc\_3}



\StringTok{\#\# CIC {-}\textgreater{} PRIV}
\StringTok{\#Indirect effect 1: cic {-}\textgreater{} aeca {-}\textgreater{} priv}
\StringTok{\#Indirect effect 2: cic {-}\textgreater{} aeca {-}\textgreater{} anx {-}\textgreater{} priv}
\StringTok{\#Indirect effect 3: cic {-}\textgreater{} anx {-}\textgreater{} priv}

\StringTok{de\_cic\_1 := o}
\StringTok{ie\_cic\_1 := j*a}
\StringTok{ie\_cic\_2 := j*i*e}
\StringTok{ie\_cic\_3 := k*e}
\StringTok{TI\_cic\_1 := ie\_cic\_1 + ie\_cic\_2 + ie\_cic\_3}
\StringTok{TE\_cic\_1 := o + TI\_cic\_1}


\StringTok{\#\# CIC {-}\textgreater{} PUB}
\StringTok{\#Indirect effect 4: cic {-}\textgreater{} aeca {-}\textgreater{} pub}
\StringTok{\#Indirect effect 5: cic {-}\textgreater{} aeca {-}\textgreater{} anx {-}\textgreater{} pub}
\StringTok{\#Indirect effect 6: cic {-}\textgreater{} anx {-}\textgreater{} pub}

\StringTok{de\_cic\_2 := p}
\StringTok{ie\_cic\_4 := j*c}
\StringTok{ie\_cic\_5 := j*i*g}
\StringTok{ie\_cic\_6 := k*g}
\StringTok{TI\_cic\_2 := ie\_cic\_4 + ie\_cic\_5 + ie\_cic\_6}
\StringTok{TE\_cic\_2 := p + TI\_cic\_2}


\StringTok{\#\# CIC {-}\textgreater{} RAD}
\StringTok{\#Indirect effect 7: cic {-}\textgreater{} aeca {-}\textgreater{} rad}
\StringTok{\#Indirect effect 8: cic {-}\textgreater{} aeca {-}\textgreater{} anx {-}\textgreater{} rad}
\StringTok{\#Indirect effect 9: cic {-}\textgreater{} anx {-}\textgreater{} rad}

\StringTok{de\_cic\_3 := q}
\StringTok{ie\_cic\_7 := j*d}
\StringTok{ie\_cic\_8 := j*i*h}
\StringTok{ie\_cic\_9 := k*h}
\StringTok{TI\_cic\_3 := ie\_cic\_7 + ie\_cic\_8 + ie\_cic\_9}
\StringTok{TE\_cic\_3 := q + TI\_cic\_3}



\StringTok{\# AECA {-}\textgreater{} PRIV}
\StringTok{\#Indirect effect 1: aeca {-}\textgreater{} anx {-}\textgreater{} priv}

\StringTok{de\_aeca\_1 := a}
\StringTok{ie\_aeca\_1 := i*e}
\StringTok{TE\_aeca\_1 := a + ie\_aeca\_1}

\StringTok{\# AECA {-}\textgreater{} PUB}
\StringTok{\#Indirect effect 2: aeca {-}\textgreater{} anx {-}\textgreater{} pub}

\StringTok{de\_aeca\_2 := c}
\StringTok{ie\_aeca\_2 := i*g}
\StringTok{TE\_aeca\_2 := c + ie\_aeca\_2}

\StringTok{\# AECA {-}\textgreater{} RAD}
\StringTok{\#Indirect effect 3: aeca {-}\textgreater{} anx {-}\textgreater{} rad}

\StringTok{de\_aeca\_3 := d}
\StringTok{ie\_aeca\_3 := i*h}
\StringTok{TE\_aeca\_3 := d + ie\_aeca\_3}


\StringTok{\# CFC {-}\textgreater{} ANX}
\StringTok{\#Indirect effect 1: cfc {-}\textgreater{} aeca {-}\textgreater{} anx}

\StringTok{de\_anx\_1 := f}
\StringTok{ie\_anx\_1 := b*i}
\StringTok{TE\_anx\_1 := f + ie\_anx\_1}

\StringTok{\# CIC {-}\textgreater{} ANX}
\StringTok{\#Indirect effect 2: cic {-}\textgreater{} aeca {-}\textgreater{} anx}

\StringTok{de\_anx\_2 := k}
\StringTok{ie\_anx\_2 := j*i}
\StringTok{TE\_anx\_2 := k + ie\_anx\_2}
\StringTok{\textquotesingle{}}
\end{Highlighting}
\end{Shaded}

\subsection{Run the model}\label{run-the-model}

\begin{Shaded}
\begin{Highlighting}[]
\CommentTok{\# Run model}
\NormalTok{fit\_m1 }\OtherTok{\textless{}{-}} \FunctionTok{sem}\NormalTok{(m1, dat\_s1, }\AttributeTok{missing =} \StringTok{"fiml"}\NormalTok{, }\AttributeTok{se =} \StringTok{"bootstrap"}\NormalTok{)}
\FunctionTok{fitmeasures}\NormalTok{(fit\_m1, }\AttributeTok{fit.measures =} \FunctionTok{c}\NormalTok{(}\StringTok{"chisq"}\NormalTok{, }\StringTok{"pvalue"}\NormalTok{, }\StringTok{"cfi"}\NormalTok{, }\StringTok{"tli"}\NormalTok{, }\StringTok{"rmsea"}\NormalTok{, }\StringTok{"srmr"}\NormalTok{))}

\CommentTok{\# Get standardized path coefficients}
\FunctionTok{summary}\NormalTok{(fit\_m1, }\AttributeTok{standardized =}\NormalTok{ T) }

\CommentTok{\# Check modification indices \textgreater{} 3.85}
\FunctionTok{modificationindices}\NormalTok{(fit\_m1, }\AttributeTok{sort. =}\NormalTok{ T, }\AttributeTok{minimum.value =} \FloatTok{3.85}\NormalTok{) }
\end{Highlighting}
\end{Shaded}

\textbf{What this is doing:}

\begin{itemize}
\item
  Full information maximum likelihood estimation is turned on, which
  means that partially missing cases are not dropped from the model and
  are instead used to improve the estimation of model effects.
\item
  \texttt{se\ =\ "bootstrap"} turns on \textbf{bootstrapped standard
  errors} instead of the model's default (usually asymptotic,
  formula-based) standard errors. Bootstrapping works by repeatedly
  resampling your data (with replacement) many times, refitting the
  model on each resample, and using the spread of results across those
  resamples to estimate how much your parameter estimates would vary
  from sample to sample.

  \begin{itemize}
  \tightlist
  \item
    This is a common choice when your data doesn't meet the normality
    assumptions that standard SEM standard errors rely on. It makes
    fewer distributional assumptions, at the cost of being more
    computationally intensive.
  \end{itemize}
\item
  \texttt{fitmeasures(fit\_m1,\ fit.measures\ =\ c(...))} extracts
  specific \textbf{model fit statistics}, rather than printing every fit
  measure lavaan can calculate (which is a long list by default).

  \begin{itemize}
  \item
    \texttt{cfi} and \texttt{tli} (Comparative Fit Index and
    Tucker-Lewis Index). Both compare the model to a baseline ``null''
    model, with values closer to 1 indicating better fit, and values
    above .95 conventionally considered good (Kline, 2015).
  \item
    \texttt{rmsea} (Root Mean Square Error of Approximation). This is a
    measure of how well the model approximates the true population
    covariance structure, with values below .06 conventionally
    considered good fit (Kline, 2015).
  \item
    \texttt{srmr} (Standardized Root Mean Square Residual). This is the
    average discrepancy between observed and model-implied correlations,
    with values below .08 conventionally considered acceptable (Hu \&
    Bentler, 1999).
  \item
    Note that the chi-square statistic and corresponding p-value is very
    sensitive to sample size (Kline, 2015).
  \end{itemize}
\item
  \texttt{summary(fit\_m1,\ standardized\ =\ TRUE)} prints the full
  model summary, including all path coefficients, variances, and
  covariances --- \texttt{standardized\ =\ TRUE} adds an extra column
  showing \textbf{standardized} coefficients, which are the ones
  typically reported in a results write-up.
\item
  \texttt{modificationindices(fit\_m1,\ sort.\ =\ TRUE,\ minimum.value\ =\ 3.85)}
  checks for potential missing paths in your model --- for every
  parameter \emph{not} currently included, it estimates how much your
  model's chi-square would improve if that parameter were added.

  \begin{itemize}
  \item
    \texttt{sort.\ =\ TRUE} orders the output from largest to smallest
    improvement, so the most impactful potential additions appear first.
  \item
    \texttt{minimum.value\ =\ 3.85} filters out trivial suggestions,
    only showing modification indices above this threshold --- 3.85
    corresponds to roughly \texttt{p\ \textless{}\ .05} on a chi-square
    distribution with 1 degree of freedom, so this is filtering to
    changes that would represent a statistically meaningful improvement
    in fit if added.
  \end{itemize}
\end{itemize}

\subsection{Multigroup SEM}\label{multigroup-sem}

To meaningfully compare across groups, measurement invariance must first
be established Kline (2015). A baseline model is established by
estimating the parameters separately for each group, and configural
invariance is met if the fit is acceptable and if the pattern of factor
loadings is consistent. The baseline model is then compared to
increasingly more constrained models. First, the factor loadings are
constrained to be equal, and metric invariance is met if doing so does
not worsen model fit compared to baseline. Then, intercepts are
constrained in addition to the factor loadings, and scalar invariance is
met if this more constrained model does not decrease model fit. For all
measurement invariance models, the chi-square is reported; however, the
practical fit indices and change in CFI and RMSEA (ΔCFI \textless{} .01
and ΔRMSEA \textless{} .015 for large samples; Chen (2007)) were used to
assess changes to model fit.

\subsubsection{Build the multigroup
model}\label{build-the-multigroup-model}

If assessing indirect effects, the model will have to be built slightly
differently so that there are unique parameters for each indirect/total
effect for each group. The following code is again taken from Pittaway
et al. (2024).

\begin{Shaded}
\begin{Highlighting}[]
\DocumentationTok{\#\#\# M2 {-} Separate Defined Parameters }\AlertTok{\#\#\#}
\NormalTok{m2 }\OtherTok{\textless{}{-}} \StringTok{\textquotesingle{}}
\StringTok{\# LATENT VARIABLES}
\StringTok{ANX =\textasciitilde{} anx\_aff + anx\_behav + anx\_rumin + anx\_pers}

\StringTok{\# STRUCTURAL MODEL (Direct Effects)}
\StringTok{conv\_private \textasciitilde{} c(e1, e2)*ANX + c(a1, a2)*ecas + c(l1, l2)*future + c(o1, o2)*immediate}
\StringTok{conv\_public \textasciitilde{} c(g1, g2)*ANX + c(c1, c2)*ecas + c(m1, m2)*future + c(p1, p2)*immediate}
\StringTok{rad\_public \textasciitilde{} c(h1, h2)*ANX + c(d1, d2)*ecas + c(n1, n2)*future + c(q1, q2)*immediate}
\StringTok{ecas \textasciitilde{} c(b1, b2)*future + c(j1, j2)*immediate}
\StringTok{ANX \textasciitilde{} c(f1, f2)*future + c(k1, k2)*immediate + c(i1, i2)*ecas}

\StringTok{\# COVARIANCES}
\StringTok{conv\_private \textasciitilde{}\textasciitilde{} conv\_public + rad\_public}
\StringTok{conv\_public \textasciitilde{}\textasciitilde{} rad\_public}

\StringTok{\# INDIRECT AND TOTAL EFFECTS}

\StringTok{\# GROUP 1}

\StringTok{\#\# CFC {-}\textgreater{} PRIV}
\StringTok{\#Indirect effect 1: cfc {-}\textgreater{} aeca {-}\textgreater{} priv}
\StringTok{\#Indirect effect 2: cfc {-}\textgreater{} aeca {-}\textgreater{} anx {-}\textgreater{} priv}
\StringTok{\#Indirect effect 3: cfc {-}\textgreater{} anx {-}\textgreater{} priv}

\StringTok{de\_cfc\_1\_1 := l1}
\StringTok{ie\_cfc\_1\_1 := b1*a1}
\StringTok{ie\_cfc\_2\_1 := b1*i1*e1}
\StringTok{ie\_cfc\_3\_1 := f1*e1}
\StringTok{TI\_cfc\_1\_1 := ie\_cfc\_1\_1 + ie\_cfc\_2\_1 + ie\_cfc\_3\_1}
\StringTok{TE\_cfc\_1\_1 := l1 + TI\_cfc\_1\_1}


\StringTok{\#\# CFC {-}\textgreater{} PUB}
\StringTok{\#Indirect effect 4: cfc {-}\textgreater{} aeca {-}\textgreater{} pub}
\StringTok{\#Indirect effect 5: cfc {-}\textgreater{} aeca {-}\textgreater{} anx {-}\textgreater{} pub}
\StringTok{\#Indirect effect 6: cfc {-}\textgreater{} anx {-}\textgreater{} pub}

\StringTok{de\_cfc\_2\_1 := m1}
\StringTok{ie\_cfc\_4\_1 := b1*c1}
\StringTok{ie\_cfc\_5\_1 := b1*i1*g1}
\StringTok{ie\_cfc\_6\_1 := f1*g1}
\StringTok{TI\_cfc\_2\_1 := ie\_cfc\_4\_1 + ie\_cfc\_5\_1 + ie\_cfc\_6\_1}
\StringTok{TE\_cfc\_2\_1 := m1 + TI\_cfc\_2\_1}


\StringTok{\#\# CFC {-}\textgreater{} RAD}
\StringTok{\#Indirect effect 7: cfc {-}\textgreater{} aeca {-}\textgreater{} rad}
\StringTok{\#Indirect effect 8: cfc {-}\textgreater{} aeca {-}\textgreater{} anx {-}\textgreater{} rad}
\StringTok{\#Indirect effect 9: cfc {-}\textgreater{} anx {-}\textgreater{} rad}

\StringTok{de\_cfc\_3\_1 := n1}
\StringTok{ie\_cfc\_7\_1 := b1*d1}
\StringTok{ie\_cfc\_8\_1 := b1*i1*h1}
\StringTok{ie\_cfc\_9\_1 := f1*h1}
\StringTok{TI\_cfc\_3\_1 := ie\_cfc\_7\_1 + ie\_cfc\_8\_1 + ie\_cfc\_9\_1}
\StringTok{TE\_cfc\_3\_1 := n1 + TI\_cfc\_3\_1}



\StringTok{\#\# CIC {-}\textgreater{} PRIV}
\StringTok{\#Indirect effect 1: cic {-}\textgreater{} aeca {-}\textgreater{} priv}
\StringTok{\#Indirect effect 2: cic {-}\textgreater{} aeca {-}\textgreater{} anx {-}\textgreater{} priv}
\StringTok{\#Indirect effect 3: cic {-}\textgreater{} anx {-}\textgreater{} priv}

\StringTok{de\_cic\_1\_1 := o1}
\StringTok{ie\_cic\_1\_1 := j1*a1}
\StringTok{ie\_cic\_2\_1 := j1*i1*e1}
\StringTok{ie\_cic\_3\_1 := k1*e1}
\StringTok{TI\_cic\_1\_1 := ie\_cic\_1\_1 + ie\_cic\_2\_1 + ie\_cic\_3\_1}
\StringTok{TE\_cic\_1\_1 := o1 + TI\_cic\_1\_1}


\StringTok{\#\# CIC {-}\textgreater{} PUB}
\StringTok{\#Indirect effect 4: cic {-}\textgreater{} aeca {-}\textgreater{} pub}
\StringTok{\#Indirect effect 5: cic {-}\textgreater{} aeca {-}\textgreater{} anx {-}\textgreater{} pub}
\StringTok{\#Indirect effect 6: cic {-}\textgreater{} anx {-}\textgreater{} pub}

\StringTok{de\_cic\_2\_1 := p1}
\StringTok{ie\_cic\_4\_1 := j1*c1}
\StringTok{ie\_cic\_5\_1 := j1*i1*g1}
\StringTok{ie\_cic\_6\_1 := k1*g1}
\StringTok{TI\_cic\_2\_1 := ie\_cic\_4\_1 + ie\_cic\_5\_1 + ie\_cic\_6\_1}
\StringTok{TE\_cic\_2\_1 := p1 + TI\_cic\_2\_1}


\StringTok{\#\# CIC {-}\textgreater{} RAD}
\StringTok{\#Indirect effect 7: cic {-}\textgreater{} aeca {-}\textgreater{} rad}
\StringTok{\#Indirect effect 8: cic {-}\textgreater{} aeca {-}\textgreater{} anx {-}\textgreater{} rad}
\StringTok{\#Indirect effect 9: cic {-}\textgreater{} anx {-}\textgreater{} rad}

\StringTok{de\_cic\_3\_1 := q1}
\StringTok{ie\_cic\_7\_1 := j1*d1}
\StringTok{ie\_cic\_8\_1 := j1*i1*h1}
\StringTok{ie\_cic\_9\_1 := k1*h1}
\StringTok{TI\_cic\_3\_1 := ie\_cic\_7\_1 + ie\_cic\_8\_1 + ie\_cic\_9\_1}
\StringTok{TE\_cic\_3\_1 := q1 + TI\_cic\_3\_1}



\StringTok{\# AECA {-}\textgreater{} PRIV}
\StringTok{\#Indirect effect 1: aeca {-}\textgreater{} anx {-}\textgreater{} priv}

\StringTok{de\_aeca\_1\_1 := a1}
\StringTok{ie\_aeca\_1\_1 := i1*e1}
\StringTok{TE\_aeca\_1\_1 := a1 + ie\_aeca\_1\_1}

\StringTok{\# AECA {-}\textgreater{} PUB}
\StringTok{\#Indirect effect 2: aeca {-}\textgreater{} anx {-}\textgreater{} pub}

\StringTok{de\_aeca\_2\_1 := c1}
\StringTok{ie\_aeca\_2\_1 := i1*g1}
\StringTok{TE\_aeca\_2\_1 := c1 + ie\_aeca\_2\_1}

\StringTok{\# AECA {-}\textgreater{} RAD}
\StringTok{\#Indirect effect 3: aeca {-}\textgreater{} anx {-}\textgreater{} rad}

\StringTok{de\_aeca\_3\_1 := d1}
\StringTok{ie\_aeca\_3\_1 := i1*h1}
\StringTok{TE\_aeca\_3\_1 := d1 + ie\_aeca\_3\_1}


\StringTok{\# CFC {-}\textgreater{} ANX}
\StringTok{\#Indirect effect 1: cfc {-}\textgreater{} aeca {-}\textgreater{} anx}

\StringTok{de\_anx\_1\_1 := f1}
\StringTok{ie\_anx\_1\_1 := b1*i1}
\StringTok{TE\_anx\_1\_1 := f1 + ie\_anx\_1\_1}

\StringTok{\# CIC {-}\textgreater{} ANX}
\StringTok{\#Indirect effect 2: cic {-}\textgreater{} aeca {-}\textgreater{} anx}

\StringTok{de\_anx\_2\_1 := k1}
\StringTok{ie\_anx\_2\_1 := j1*i1}
\StringTok{TE\_anx\_2\_1 := k1 + ie\_anx\_2\_1}




\StringTok{\# GROUP 2}

\StringTok{\#\# CFC {-}\textgreater{} PRIV}
\StringTok{\#Indirect effect 1: cfc {-}\textgreater{} aeca {-}\textgreater{} priv}
\StringTok{\#Indirect effect 2: cfc {-}\textgreater{} aeca {-}\textgreater{} anx {-}\textgreater{} priv}
\StringTok{\#Indirect effect 3: cfc {-}\textgreater{} anx {-}\textgreater{} priv}

\StringTok{de\_cfc\_1\_2 := l2}
\StringTok{ie\_cfc\_1\_2 := b2*a2}
\StringTok{ie\_cfc\_2\_2 := b2*i2*e2}
\StringTok{ie\_cfc\_3\_2 := f2*e2}
\StringTok{TI\_cfc\_1\_2 := ie\_cfc\_1\_2 + ie\_cfc\_2\_2 + ie\_cfc\_3\_2}
\StringTok{TE\_cfc\_1\_2 := l2 + TI\_cfc\_1\_2}


\StringTok{\#\# CFC {-}\textgreater{} PUB}
\StringTok{\#Indirect effect 4: cfc {-}\textgreater{} aeca {-}\textgreater{} pub}
\StringTok{\#Indirect effect 5: cfc {-}\textgreater{} aeca {-}\textgreater{} anx {-}\textgreater{} pub}
\StringTok{\#Indirect effect 6: cfc {-}\textgreater{} anx {-}\textgreater{} pub}

\StringTok{de\_cfc\_2\_2 := m2}
\StringTok{ie\_cfc\_4\_2 := b2*c2}
\StringTok{ie\_cfc\_5\_2 := b2*i2*g2}
\StringTok{ie\_cfc\_6\_2 := f2*g2}
\StringTok{TI\_cfc\_2\_2 := ie\_cfc\_4\_2 + ie\_cfc\_5\_2 + ie\_cfc\_6\_2}
\StringTok{TE\_cfc\_2\_2 := m2 + TI\_cfc\_2\_2}


\StringTok{\#\# CFC {-}\textgreater{} RAD}
\StringTok{\#Indirect effect 7: cfc {-}\textgreater{} aeca {-}\textgreater{} rad}
\StringTok{\#Indirect effect 8: cfc {-}\textgreater{} aeca {-}\textgreater{} anx {-}\textgreater{} rad}
\StringTok{\#Indirect effect 9: cfc {-}\textgreater{} anx {-}\textgreater{} rad}

\StringTok{de\_cfc\_3\_2 := n2}
\StringTok{ie\_cfc\_7\_2 := b2*d2}
\StringTok{ie\_cfc\_8\_2 := b2*i2*h2}
\StringTok{ie\_cfc\_9\_2 := f2*h2}
\StringTok{TI\_cfc\_3\_2 := ie\_cfc\_7\_2 + ie\_cfc\_8\_2 + ie\_cfc\_9\_2}
\StringTok{TE\_cfc\_3\_2 := n2 + TI\_cfc\_3\_2}



\StringTok{\#\# CIC {-}\textgreater{} PRIV}
\StringTok{\#Indirect effect 1: cic {-}\textgreater{} aeca {-}\textgreater{} priv}
\StringTok{\#Indirect effect 2: cic {-}\textgreater{} aeca {-}\textgreater{} anx {-}\textgreater{} priv}
\StringTok{\#Indirect effect 3: cic {-}\textgreater{} anx {-}\textgreater{} priv}

\StringTok{de\_cic\_1\_2 := o2}
\StringTok{ie\_cic\_1\_2 := j2*a2}
\StringTok{ie\_cic\_2\_2 := j2*i2*e2}
\StringTok{ie\_cic\_3\_2 := k2*e2}
\StringTok{TI\_cic\_1\_2 := ie\_cic\_1\_2 + ie\_cic\_2\_2 + ie\_cic\_3\_2}
\StringTok{TE\_cic\_1\_2 := o2 + TI\_cic\_1\_2}


\StringTok{\#\# CIC {-}\textgreater{} PUB}
\StringTok{\#Indirect effect 4: cic {-}\textgreater{} aeca {-}\textgreater{} pub}
\StringTok{\#Indirect effect 5: cic {-}\textgreater{} aeca {-}\textgreater{} anx {-}\textgreater{} pub}
\StringTok{\#Indirect effect 6: cic {-}\textgreater{} anx {-}\textgreater{} pub}

\StringTok{de\_cic\_2\_2 := p2}
\StringTok{ie\_cic\_4\_2 := j2*c2}
\StringTok{ie\_cic\_5\_2 := j2*i2*g2}
\StringTok{ie\_cic\_6\_2 := k2*g2}
\StringTok{TI\_cic\_2\_2 := ie\_cic\_4\_2 + ie\_cic\_5\_2 + ie\_cic\_6\_2}
\StringTok{TE\_cic\_2\_2 := p2 + TI\_cic\_2\_2}


\StringTok{\#\# CIC {-}\textgreater{} RAD}
\StringTok{\#Indirect effect 7: cic {-}\textgreater{} aeca {-}\textgreater{} rad}
\StringTok{\#Indirect effect 8: cic {-}\textgreater{} aeca {-}\textgreater{} anx {-}\textgreater{} rad}
\StringTok{\#Indirect effect 9: cic {-}\textgreater{} anx {-}\textgreater{} rad}

\StringTok{de\_cic\_3\_2 := q2}
\StringTok{ie\_cic\_7\_2 := j2*d2}
\StringTok{ie\_cic\_8\_2 := j2*i2*h2}
\StringTok{ie\_cic\_9\_2 := k2*h2}
\StringTok{TI\_cic\_3\_2 := ie\_cic\_7\_2 + ie\_cic\_8\_2 + ie\_cic\_9\_2}
\StringTok{TE\_cic\_3\_2 := q2 + TI\_cic\_3\_2}



\StringTok{\# AECA {-}\textgreater{} PRIV}
\StringTok{\#Indirect effect 1: aeca {-}\textgreater{} anx {-}\textgreater{} priv}

\StringTok{de\_aeca\_1\_2 := a2}
\StringTok{ie\_aeca\_1\_2 := i2*e2}
\StringTok{TE\_aeca\_1\_2 := a2 + ie\_aeca\_1\_2}

\StringTok{\# AECA {-}\textgreater{} PUB}
\StringTok{\#Indirect effect 2: aeca {-}\textgreater{} anx {-}\textgreater{} pub}

\StringTok{de\_aeca\_2\_2 := c2}
\StringTok{ie\_aeca\_2\_2 := i2*g2}
\StringTok{TE\_aeca\_2\_2 := c2 + ie\_aeca\_2\_2}

\StringTok{\# AECA {-}\textgreater{} RAD}
\StringTok{\#Indirect effect 3: aeca {-}\textgreater{} anx {-}\textgreater{} rad}

\StringTok{de\_aeca\_3\_2 := d2}
\StringTok{ie\_aeca\_3\_2 := i2*h2}
\StringTok{TE\_aeca\_3\_2 := d2 + ie\_aeca\_3\_2}


\StringTok{\# CFC {-}\textgreater{} ANX}
\StringTok{\#Indirect effect 1: cfc {-}\textgreater{} aeca {-}\textgreater{} anx}

\StringTok{de\_anx\_1\_2 := f2}
\StringTok{ie\_anx\_1\_2 := b2*i2}
\StringTok{TE\_anx\_1\_2 := f2 + ie\_anx\_1\_2}

\StringTok{\# CIC {-}\textgreater{} ANX}
\StringTok{\#Indirect effect 2: cic {-}\textgreater{} aeca {-}\textgreater{} anx}

\StringTok{de\_anx\_2\_2 := k2}
\StringTok{ie\_anx\_2\_2 := j2*i2}
\StringTok{TE\_anx\_2\_2 := k2 + ie\_anx\_2\_2}
\StringTok{\textquotesingle{}}
\end{Highlighting}
\end{Shaded}

\subsubsection{Measurement invariance}\label{measurement-invariance}

First, establish the baseline model by estimating the parameters
separately for each group.

\begin{Shaded}
\begin{Highlighting}[]
\CommentTok{\# Establish baseline model by estimating parameters separately for each group}
\NormalTok{baseline\_1 }\OtherTok{\textless{}{-}} \FunctionTok{sem}\NormalTok{(m2, dat\_s1, }\AttributeTok{group =} \StringTok{"counterbalancing"}\NormalTok{, }\AttributeTok{missing =} \StringTok{"fiml"}\NormalTok{) }

\FunctionTok{fitmeasures}\NormalTok{(baseline\_1, }\AttributeTok{fit.measures =} \FunctionTok{c}\NormalTok{(}\StringTok{"chisq"}\NormalTok{, }\StringTok{"pvalue"}\NormalTok{, }\StringTok{"cfi"}\NormalTok{, }\StringTok{"tli"}\NormalTok{, }\StringTok{"rmsea"}\NormalTok{, }\StringTok{"srmr"}\NormalTok{))}

\FunctionTok{summary}\NormalTok{(baseline\_1)}
\end{Highlighting}
\end{Shaded}

\begin{itemize}
\tightlist
\item
  Configural invariance is met if the model fits and the pattern of
  factor loadings is the same for each group.
\end{itemize}

Next, establish metric invariance by constraining the factor loadings to
be equal.

\begin{Shaded}
\begin{Highlighting}[]
\CommentTok{\# Establish metric invariance by constraining the factor loadings to be equal}
\NormalTok{metric\_1 }\OtherTok{\textless{}{-}} \FunctionTok{sem}\NormalTok{(m2, dat\_s1, }\AttributeTok{group =} \StringTok{"counterbalancing"}\NormalTok{, }\AttributeTok{group.equal =} \StringTok{"loadings"}\NormalTok{, }\AttributeTok{missing =} \StringTok{"fiml"}\NormalTok{, }\AttributeTok{se =} \StringTok{"bootstrap"}\NormalTok{)}

\FunctionTok{fitmeasures}\NormalTok{(metric\_1, }\AttributeTok{fit.measures =} \FunctionTok{c}\NormalTok{(}\StringTok{"chisq"}\NormalTok{, }\StringTok{"pvalue"}\NormalTok{, }\StringTok{"cfi"}\NormalTok{, }\StringTok{"tli"}\NormalTok{, }\StringTok{"rmsea"}\NormalTok{, }\StringTok{"srmr"}\NormalTok{))}
\end{Highlighting}
\end{Shaded}

Next, establish scalar invariance by constraining factor loadings +
observed means to be equal.

\begin{Shaded}
\begin{Highlighting}[]
\CommentTok{\# Establish scalar invariance by constraining factor loadings + observed means to be equal}
\NormalTok{scalar\_1 }\OtherTok{\textless{}{-}} \FunctionTok{sem}\NormalTok{(m2, dat\_s1, }\AttributeTok{group =} \StringTok{"counterbalancing"}\NormalTok{, }\AttributeTok{group.equal =} \FunctionTok{c}\NormalTok{(}\StringTok{"loadings"}\NormalTok{, }\StringTok{"intercepts"}\NormalTok{), }\AttributeTok{se =} \StringTok{"bootstrap"}\NormalTok{, }\AttributeTok{missing =} \StringTok{"fiml"}\NormalTok{)}

\FunctionTok{fitmeasures}\NormalTok{(scalar\_1, }\AttributeTok{fit.measures =} \FunctionTok{c}\NormalTok{(}\StringTok{"chisq"}\NormalTok{, }\StringTok{"pvalue"}\NormalTok{, }\StringTok{"cfi"}\NormalTok{, }\StringTok{"tli"}\NormalTok{, }\StringTok{"rmsea"}\NormalTok{, }\StringTok{"srmr"}\NormalTok{)) }

\CommentTok{\# Compare models}
\FunctionTok{summary}\NormalTok{(}\FunctionTok{compareFit}\NormalTok{(baseline\_1, metric\_1, scalar\_1))}
\end{Highlighting}
\end{Shaded}

\begin{itemize}
\item
  Metric invariance is established if the change in CFI and RMSEA are
  below the recommended thresholds.
\item
  Scalar invariance is established if the change in CFI and RMSEA are
  below the recommended thresholds.
\end{itemize}

If scalar invariance is established, then the model is appropriately
measurement-invariant across groups which means you can meaningfully
compare the pathways and effect sizes. At this point, you can get the
standardised coefficients for the scalar model. If you were creating a
new scale, however, and you would feasibly expect that the groups SHOULD
be the same, you would have to establish one final level of invariance
(the strictest level ever), which I cannot currently remember. If it
becomes relevant to me, I'll update this page.

\begin{tcolorbox}[enhanced jigsaw, toprule=.15mm, coltitle=black, colframe=quarto-callout-note-color-frame, colbacktitle=quarto-callout-note-color!10!white, opacityback=0, left=2mm, breakable, rightrule=.15mm, bottomrule=.15mm, leftrule=.75mm, titlerule=0mm, arc=.35mm, opacitybacktitle=0.6, toptitle=1mm, bottomtitle=1mm, title=\textcolor{quarto-callout-note-color}{\faInfo}\hspace{0.5em}{Note}, colback=white]

If running complicated models with bootstrapping, expect the models to
take a long time to estimate. This is normal!

\end{tcolorbox}

\subsubsection{Get multigroup model
coefficients}\label{get-multigroup-model-coefficients}

\begin{Shaded}
\begin{Highlighting}[]
\CommentTok{\# Get standardised coefficients}
\FunctionTok{summary}\NormalTok{(scalar\_1, }\AttributeTok{standardized =}\NormalTok{ T)}
\end{Highlighting}
\end{Shaded}

\section{Partial-least-squares SEM}\label{partial-least-squares-sem}

Regular SEM, as outlined above, is \emph{covariance-based SEM} (CB-SEM).
CB-SEM is primarily used to confirm or reject theories and their
underlying hypotheses. It works by determining how closely a proposed
theoretical model can reproduce the covariance matrix for an observed
sample dataset.

\emph{Partial least squares SEM} (PLS-SEM) is a second method of
conducting SEM. It was introduced more as a `causal-predictive' approach
to SEM, and focuses on explaining the variance in the model's DVs. It's
also better at coping with smaller samples and complex models. Another
good thing about PLS-SEM is that it is non-parametric, meaning that
non-normally distributed data is not a problem.

It uses a different approach to analysis that is explained in detail in
the following sources:

\begin{itemize}
\item
  Hair, J. F., Ringle, C. M., \& Sarstedt, M. (2011). PLS-SEM: Indeed a
  silver bullet. \emph{Journal of Marketing Theory and Practice},
  \emph{19}(2), 139--152.
  \url{https://doi.org/10.2753/MTP1069-6679190202}
\item
  Hair, J. F., Hult, G. T. M., Ringle, C. M., Sarstedt, M., Danks, N.
  P., \& Ray, S. (2021). \emph{Partial Least Squares Structural Equation
  Modeling (PLS-SEM) Using R: A Workbook}. Springer International
  Publishing. \url{https://doi.org/10.1007/978-3-030-80519-7}
\item
  Hair, J. F., Hult, G. T. M., Ringle, C. M., Sarstedt, M., Danks, N.
  P., \& Ray, S. (2021). \emph{Partial Least Squares Structural Equation
  Modeling (PLS-SEM) Using R: A Workbook}. Springer International
  Publishing. \url{https://doi.org/10.1007/978-3-030-80519-7}
\end{itemize}

The one downside is that there are no established goodness-of-fit
indicators, so its use for theory testing and confirmation is pretty
limited. But, if you have already run confirmatory analyses showing
support for your model (e.g., in an earlier paper), then PLS-SEM can
overcome sample size/distributional issues.

\textbf{Packages required}:

\begin{itemize}
\item
  cSEM
  (\url{https://cran.r-project.org/web/packages/cSEM/vignettes/cSEM.html})
\item
  lavaan (\url{https://lavaan.ugent.be/tutorial/})
\end{itemize}

The cSEM package uses lavaan syntax to run the PLS-SEM.

\subsection{Build the model using
lavaan}\label{build-the-model-using-lavaan}

The following code is taken from Pittaway et al. (2025). This tests
practically the same model as in the example above, except with
additional longer indirect pathways to the three DVs.

\begin{Shaded}
\begin{Highlighting}[]
\DocumentationTok{\#\#\# M1 }\AlertTok{\#\#\#}
\NormalTok{m1\_csem }\OtherTok{\textless{}{-}} \StringTok{\textquotesingle{}}
\StringTok{ANX \textless{}\textasciitilde{} anx\_aff\_w1 + anx\_behav\_w1 + anx\_rumin\_w1 + anx\_pers\_w1}
\StringTok{CFC \textless{}\textasciitilde{} cfc\_w1}
\StringTok{CIC \textless{}\textasciitilde{} cic\_w1}
\StringTok{PRIV \textless{}\textasciitilde{} conv\_priv\_w1}
\StringTok{PUB \textless{}\textasciitilde{} conv\_pub\_w1}
\StringTok{RAD \textless{}\textasciitilde{} rad\_pub\_w1}
\StringTok{ECAS \textless{}\textasciitilde{} ecas\_w1}
\StringTok{aPRIV \textless{}\textasciitilde{} act\_conv\_priv\_w2}
\StringTok{aPUB \textless{}\textasciitilde{} act\_conv\_pub\_w2}
\StringTok{aRAD \textless{}\textasciitilde{} act\_rad\_pub\_w2}

\StringTok{\# STRUCTURAL MODEL }
\StringTok{PRIV \textasciitilde{} ANX + ECAS + CFC + CIC}
\StringTok{PUB \textasciitilde{} ANX + ECAS + CFC + CIC}
\StringTok{RAD \textasciitilde{} ANX + ECAS + CFC + CIC}
\StringTok{ECAS \textasciitilde{} CFC + CIC}
\StringTok{ANX \textasciitilde{} CFC + CIC + ECAS}
\StringTok{aPRIV \textasciitilde{} PRIV}
\StringTok{aPUB \textasciitilde{} PUB}
\StringTok{aRAD \textasciitilde{} RAD}
\StringTok{\textquotesingle{}}
\end{Highlighting}
\end{Shaded}

\textbf{Note:} Unlike regular SEM, the cSEM package cannot handle the
separate indirect effects.

\subsection{Implement PLS-SEM}\label{implement-pls-sem}

\begin{Shaded}
\begin{Highlighting}[]
\CommentTok{\# Basic implementation}
\NormalTok{results\_csem }\OtherTok{\textless{}{-}} \FunctionTok{csem}\NormalTok{(}
  \AttributeTok{.data =}\NormalTok{ dat, }
  \AttributeTok{.model =}\NormalTok{ m1\_csem, }
  \AttributeTok{.approach\_weights =} \StringTok{"PLS{-}PM"}\NormalTok{,}
  \AttributeTok{.approach\_paths =} \StringTok{"OLS"}
\NormalTok{  )}

\NormalTok{bootstrap\_results }\OtherTok{\textless{}{-}} \FunctionTok{resamplecSEMResults}\NormalTok{(}
  \AttributeTok{.object =}\NormalTok{ results\_csem,}
  \AttributeTok{.resample\_method =} \StringTok{"bootstrap"}\NormalTok{,}
  \AttributeTok{.R =} \DecValTok{1000}\NormalTok{,}
  \AttributeTok{.seed =} \DecValTok{123456}
\NormalTok{)}

\CommentTok{\# Assess measurement model}
\NormalTok{results\_csem}

\CommentTok{\# View results}
\FunctionTok{summarize}\NormalTok{(bootstrap\_results)}
\NormalTok{results\_csem}\SpecialCharTok{$}\NormalTok{Estimates}\SpecialCharTok{$}\NormalTok{R2}
\FunctionTok{verify}\NormalTok{(results\_csem)}

\FunctionTok{assess}\NormalTok{(bootstrap\_results)}
\FunctionTok{predict}\NormalTok{(bootstrap\_results, }\AttributeTok{.approach\_predict =} \StringTok{"direct"}\NormalTok{)}
\end{Highlighting}
\end{Shaded}

\subsection{Evaluate + interpret the
model}\label{evaluate-interpret-the-model}

First, assess the variance inflation factors (VIFs) to ensure they are
all below the recommended threshold of 3, which indicates that
collinearity that would otherwise bias the regression results is not a
problem in the present data (Hair et al., 2019).

Evaluation of the structural model is done in the following ways:

\begin{enumerate}
\def\labelenumi{\arabic{enumi}.}
\item
  Significance of the hypothesised paths
\item
  The size of the effects
\item
  The explanatory power of the model (R2): \emph{R2} values greater than
  .75, .50, and .25 denoted substantial, moderate, and weak power,
  respectively (Hair et al., 2011).
\item
  The predictive relevance of the model (Q2): \emph{Q2} values of
  greater than 0, .25, and .50 indicated small, medium, and large
  predictive relevance, while negative values indicated no predictive
  relevance (Hair et al., 2011).
\end{enumerate}

\subsection{Manually calculate indirect
effects}\label{manually-calculate-indirect-effects}

The total indirect effects and total effects are generated in the
initial output, but you have to write a function to generate the
specific indirect pathways if you have multiple. This code was adapted
from \url{https://github.com/FloSchuberth/cSEM/issues/457}.

A separate function is created for every unique indirect pathway in the
model. This is an example for one unique indirect pathway.

\begin{Shaded}
\begin{Highlighting}[]
\CommentTok{\# Define function to calculate the specific indirect effect CFC {-}\textgreater{} PRIV {-}\textgreater{} aPRIV}
\NormalTok{indirectCFC\_PRIV\_aPRIV }\OtherTok{\textless{}{-}} \ControlFlowTok{function}\NormalTok{(.object) \{}
\NormalTok{  .object}\SpecialCharTok{$}\NormalTok{Estimates}\SpecialCharTok{$}\NormalTok{Path\_estimates[}\StringTok{\textquotesingle{}PRIV\textquotesingle{}}\NormalTok{, }\StringTok{\textquotesingle{}CFC\textquotesingle{}}\NormalTok{] }\SpecialCharTok{*}\NormalTok{ .object}\SpecialCharTok{$}\NormalTok{Estimates}\SpecialCharTok{$}\NormalTok{Path\_estimates[}\StringTok{\textquotesingle{}aPRIV\textquotesingle{}}\NormalTok{, }\StringTok{\textquotesingle{}PRIV\textquotesingle{}}\NormalTok{]}
\NormalTok{\}}

\CommentTok{\# Run cSEM with the user{-}defined function}
\NormalTok{result\_indirectCFC\_PRIV\_aPRIV }\OtherTok{\textless{}{-}} \FunctionTok{csem}\NormalTok{(}
  \AttributeTok{.data =}\NormalTok{ dat,}
  \AttributeTok{.model =}\NormalTok{ m1\_csem,}
  \AttributeTok{.resample\_method =} \StringTok{"bootstrap"}\NormalTok{,}
  \AttributeTok{.R =} \DecValTok{100}\NormalTok{,}
  \AttributeTok{.seed =} \DecValTok{147150973}\NormalTok{,}
  \AttributeTok{.user\_funs =}\NormalTok{ indirectCFC\_PRIV\_aPRIV}
\NormalTok{)}

\CommentTok{\# Get the coefficient of the indirect effect}
\NormalTok{result\_indirectCFC\_PRIV\_aPRIV}\SpecialCharTok{$}\NormalTok{Estimates}\SpecialCharTok{$}\NormalTok{Estimates\_resample}\SpecialCharTok{$}\NormalTok{Estimates1}
\FunctionTok{infer}\NormalTok{(result\_indirectCFC\_PRIV\_aPRIV, }\AttributeTok{.quantity =} \StringTok{\textquotesingle{}CI\_percentile\textquotesingle{}}\NormalTok{)}
\end{Highlighting}
\end{Shaded}

\chapter{Meta-Analysis}\label{meta-analysis}

Contact me for a full three-level random-effects meta-analysis workflow
from start to finish. I can't upload it yet, as it uses data from a
study that has not yet been published. When it has been published, I
will upload it here.

But in the meantime\ldots{} let me know if you want me to send you the
code and explanations.

📧
\href{mailto:charlie.pittaway@flinders.edu.au}{\nolinkurl{charlie.pittaway@flinders.edu.au}}

\part{Data Visualisation}

\chapter{Themes and Plot Colouring}\label{themes-and-plot-colouring}

\section{APA-formatted figures}\label{apa-formatted-figures}

There are two approaches here: using a purpose-built theme from an
existing package, or building the APA look manually with
\texttt{theme()}. Both are worth knowing.

\textbf{Packages required:}

\begin{itemize}
\item
  ggplot2 (part of \textbf{tidyverse})
\item
  jtools (J. A. Long, 2024)
\end{itemize}

\subsection{Using the theme\_apa() function in
ggplot2}\label{using-the-theme_apa-function-in-ggplot2}

\begin{Shaded}
\begin{Highlighting}[]
\CommentTok{\#| eval: false}
\FunctionTok{ggplot}\NormalTok{(data, }\FunctionTok{aes}\NormalTok{(}\AttributeTok{x =}\NormalTok{ condition, }\AttributeTok{y =}\NormalTok{ outcome)) }\SpecialCharTok{+}
  \FunctionTok{geom\_boxplot}\NormalTok{() }\SpecialCharTok{+}
  \FunctionTok{labs}\NormalTok{(}\AttributeTok{x =} \StringTok{"Condition"}\NormalTok{, }\AttributeTok{y =} \StringTok{"Outcome Score"}\NormalTok{) }\SpecialCharTok{+}
  \FunctionTok{theme\_apa}\NormalTok{()}
\end{Highlighting}
\end{Shaded}

\textbf{What this is doing:}

\begin{itemize}
\tightlist
\item
  \texttt{theme\_apa()} applies a pre-built theme designed to match APA
  figure conventions in one line --- removing gridlines, using a clean
  black-and-white base, and adjusting font/line styling to match APA's
  expectations for published figures.
\end{itemize}

\begin{tcolorbox}[enhanced jigsaw, toprule=.15mm, coltitle=black, colframe=quarto-callout-tip-color-frame, colbacktitle=quarto-callout-tip-color!10!white, opacityback=0, left=2mm, breakable, rightrule=.15mm, bottomrule=.15mm, leftrule=.75mm, titlerule=0mm, arc=.35mm, opacitybacktitle=0.6, toptitle=1mm, bottomtitle=1mm, title=\textcolor{quarto-callout-tip-color}{\faLightbulb}\hspace{0.5em}{Tip}, colback=white]

\texttt{theme\_apa()} has a few optional arguments worth knowing:
\texttt{legend.pos\ =\ "none"} removes the legend entirely (useful if
your plot doesn't need one), and \texttt{facet.title.size} adjusts facet
label sizing if you're using
\texttt{facet\_wrap()}/\texttt{facet\_grid()}. Run \texttt{?theme\_apa}
for the full list, and visit this cookbook for the full ggplot2
customisation: (Chang, 2026).

\end{tcolorbox}

\subsection{Building the APA look
manually}\label{building-the-apa-look-manually}

These are the APA conventions specified manually in ggplot2:

\begin{Shaded}
\begin{Highlighting}[]
\CommentTok{\#| eval: false}
\FunctionTok{ggplot}\NormalTok{(dat, }\FunctionTok{aes}\NormalTok{(}\AttributeTok{x =}\NormalTok{ condition, }\AttributeTok{y =}\NormalTok{ outcome)) }\SpecialCharTok{+}
  \FunctionTok{geom\_boxplot}\NormalTok{() }\SpecialCharTok{+}
  \FunctionTok{labs}\NormalTok{(}\AttributeTok{x =} \StringTok{"Condition"}\NormalTok{, }\AttributeTok{y =} \StringTok{"Outcome Score"}\NormalTok{) }\SpecialCharTok{+}
  \FunctionTok{theme\_classic}\NormalTok{(}\AttributeTok{base\_size =} \DecValTok{12}\NormalTok{) }\SpecialCharTok{+}
  \FunctionTok{theme}\NormalTok{(}
    \AttributeTok{axis.line =} \FunctionTok{element\_line}\NormalTok{(}\AttributeTok{colour =} \StringTok{"black"}\NormalTok{),}
    \AttributeTok{axis.text =} \FunctionTok{element\_text}\NormalTok{(}\AttributeTok{colour =} \StringTok{"black"}\NormalTok{),}
    \AttributeTok{legend.title =} \FunctionTok{element\_blank}\NormalTok{(),}
    \AttributeTok{legend.position =} \StringTok{"bottom"}\NormalTok{,}
    \AttributeTok{panel.grid =} \FunctionTok{element\_blank}\NormalTok{(),}
    \AttributeTok{plot.title =} \FunctionTok{element\_text}\NormalTok{(}\AttributeTok{hjust =} \FloatTok{0.5}\NormalTok{)}
\NormalTok{  )}
\end{Highlighting}
\end{Shaded}

\begin{tcolorbox}[enhanced jigsaw, toprule=.15mm, coltitle=black, colframe=quarto-callout-tip-color-frame, colbacktitle=quarto-callout-tip-color!10!white, opacityback=0, left=2mm, breakable, rightrule=.15mm, bottomrule=.15mm, leftrule=.75mm, titlerule=0mm, arc=.35mm, opacitybacktitle=0.6, toptitle=1mm, bottomtitle=1mm, title=\textcolor{quarto-callout-tip-color}{\faLightbulb}\hspace{0.5em}{Tip}, colback=white]

You can apply these APA formatting conventions to any kind of ggplot2
figure: just add the \texttt{theme()} part to the end of the plot code.

\end{tcolorbox}

\section{Colour blindness}\label{colour-blindness}

Use the \textbf{viridis} colour package for \textbf{ggplot2}. This has
colours designed to be seen regardless of colour blindness.

More information available here:
\hyperref[0]{https://ggplot2.tidyverse.org/reference/scale\_viridis.html}

\part{Bayesian Analysis}

\chapter{Bayesian Meta Analysis}\label{bayesian-meta-analysis}

This page is under construction\ldots{} 🚧

In the meantime, here is a vignette for a package I found to run
Bayesian meta-analysis:
\url{https://cran.r-project.org/web//packages/bayesmeta/vignettes/bayesmeta.html}

And a link to a paper about Bayesian meta-analysis:
\url{https://link.springer.com/content/pdf/10.1007/s40273-025-01510-2.pdf}

\part{Miscellaneous}

\chapter{Index of R Packages}\label{index-of-r-packages}

Here is a list of things you might want to do in R, and the packages
that I have found to help with them. This list will be constantly
updated as I learn new methodologies and find new and better processes
to do things.

FYI, ``vignettes'' are instructions manuals and examples of how to use
the package.

\section{Data Wrangling}\label{data-wrangling}

\begin{longtable}[]{@{}
  >{\raggedright\arraybackslash}p{(\linewidth - 4\tabcolsep) * \real{0.3333}}
  >{\raggedright\arraybackslash}p{(\linewidth - 4\tabcolsep) * \real{0.3333}}
  >{\raggedright\arraybackslash}p{(\linewidth - 4\tabcolsep) * \real{0.3333}}@{}}
\toprule\noalign{}
\begin{minipage}[b]{\linewidth}\raggedright
Package (Citation)
\end{minipage} & \begin{minipage}[b]{\linewidth}\raggedright
Vignettes
\end{minipage} & \begin{minipage}[b]{\linewidth}\raggedright
Notes
\end{minipage} \\
\midrule\noalign{}
\endhead
\bottomrule\noalign{}
\endlastfoot
tidyverse (Wickham et al., 2019) &
\url{https://tidyverse.tidyverse.org/} & An essential, non-negotiable
collection of packages. \\
\end{longtable}

\section{Psychology Convenience
Packages}\label{psychology-convenience-packages}

\begin{longtable}[]{@{}
  >{\raggedright\arraybackslash}p{(\linewidth - 6\tabcolsep) * \real{0.0658}}
  >{\raggedright\arraybackslash}p{(\linewidth - 6\tabcolsep) * \real{0.1567}}
  >{\raggedright\arraybackslash}p{(\linewidth - 6\tabcolsep) * \real{0.6144}}
  >{\raggedright\arraybackslash}p{(\linewidth - 6\tabcolsep) * \real{0.1567}}@{}}
\toprule\noalign{}
\begin{minipage}[b]{\linewidth}\raggedright
Package (Citation)
\end{minipage} & \begin{minipage}[b]{\linewidth}\raggedright
Vignettes
\end{minipage} & \begin{minipage}[b]{\linewidth}\raggedright
Notes
\end{minipage} & \begin{minipage}[b]{\linewidth}\raggedright
Used in
\end{minipage} \\
\midrule\noalign{}
\endhead
\bottomrule\noalign{}
\endlastfoot
psych (William Revelle, 2026) & \url{https://personality-project.org/r/}
& A general purpose toolbox developed originally for personality,
psychometric theory and experimental psychology. &
\begin{minipage}[t]{\linewidth}\raggedright
\begin{itemize}
\tightlist
\item
  \href{data-cleaning-overview.qmd}{Data cleaning}
\end{itemize}
\end{minipage} \\
rempsyc (Thériault, 2023) &
\url{https://rempsyc.remi-theriault.com/index.html} & R package of
convenience functions often used by psychological researchers (e.g.,
plots, nice APA tables, easily run statistical tests or check
assumptions, and automatise various other tasks). &
\begin{minipage}[t]{\linewidth}\raggedright
\begin{itemize}
\tightlist
\item
  \href{t-tests.qmd}{T-tests}
\end{itemize}
\end{minipage} \\
\end{longtable}

\section{Reproducible Tables}\label{reproducible-tables-1}

\begin{longtable}[]{@{}
  >{\raggedright\arraybackslash}p{(\linewidth - 6\tabcolsep) * \real{0.0422}}
  >{\raggedright\arraybackslash}p{(\linewidth - 6\tabcolsep) * \real{0.1179}}
  >{\raggedright\arraybackslash}p{(\linewidth - 6\tabcolsep) * \real{0.7118}}
  >{\raggedright\arraybackslash}p{(\linewidth - 6\tabcolsep) * \real{0.1252}}@{}}
\toprule\noalign{}
\begin{minipage}[b]{\linewidth}\raggedright
Package (Citation)
\end{minipage} & \begin{minipage}[b]{\linewidth}\raggedright
Vignettes
\end{minipage} & \begin{minipage}[b]{\linewidth}\raggedright
Notes
\end{minipage} & \begin{minipage}[b]{\linewidth}\raggedright
Used in
\end{minipage} \\
\midrule\noalign{}
\endhead
\bottomrule\noalign{}
\endlastfoot
apaTables (Stanley, 2026) &
\url{https://dstanley4.github.io/apaTables/index.html} & ⭐A perfect
angel of a package for reporting correlation tables, regression tables,
and ANOVA tables with APA 7 specifications.

The downside of this package is that it has a very narrow focus on just
APA conventions and just these few tables types, with limited
customisation/table building options. & 📌To add: Link to personal
guides for initial manuscript tables (corr tables, etc.) \\
tableone (Yoshida \& Bartel, 2022) &
\url{https://cran.r-project.org/web/packages/tableone/vignettes/introduction.html}
& Create easy ``Table 1'' (i.e., demographic) tables. &
\begin{minipage}[t]{\linewidth}\raggedright
\begin{itemize}
\tightlist
\item
  \href{table-one.qmd}{Building a demographics table with Table One}
\end{itemize}
\end{minipage} \\
gtsummary (Sjoberg et al., 2021) &
\url{https://www.danieldsjoberg.com/gtsummary/index.html} & A more
comprehensive, highly customisable package for creating
publication-ready tables. ⭐Can present multiple regression models
side-by-side in the one table (which apaTables cannot do). But I don't
know how to use it, so if you figure it out, let me know. & \\
flextable (Gohel \& Skintzos, 2026) &
\url{https://ardata-fr.github.io/flextable-book/index.html} &
\multicolumn{2}{>{\raggedright\arraybackslash}p{(\linewidth - 6\tabcolsep) * \real{0.8370} + 2\tabcolsep}@{}}{%
Provides a framework to easily create tables for \textbar{} reporting
and publications. \textbar{} \textbar{} ⚠️It can be a bit difficult to
use for a beginner, but I often embed it into my workflow to create
custom tables (with a lot of help from the internet). \textbar{}
\textbar{} \textbar{} \textbar{} \textbar{} \textbar{}} \\
rempsyc (Thériault, 2023) &
\url{https://rempsyc.remi-theriault.com/index.html} & This convenience
package includes functions to easily create ``nice'' APA tables (i.e.,
correctly formatted). For example, lots of outputs created in R will
save p-values, lower 95\% CIs, and upper 95\% CIs in separate columns.
If these are named correctly, the nice\_table() function in
\textbf{rempsych} will automatically detect them and format them
correctly. Similarly, you can easily create APA-formatted t-test tables,
planned contrasts, moderation tables, violin plots, and scatter plots. &
\begin{minipage}[t]{\linewidth}\raggedright
\begin{itemize}
\tightlist
\item
  \href{t-tests.qmd}{T-tests}
\end{itemize}
\end{minipage} \\
\end{longtable}

\section{Data Quality}\label{data-quality}

\begin{longtable}[]{@{}
  >{\raggedright\arraybackslash}p{(\linewidth - 6\tabcolsep) * \real{0.1489}}
  >{\raggedright\arraybackslash}p{(\linewidth - 6\tabcolsep) * \real{0.3121}}
  >{\raggedright\arraybackslash}p{(\linewidth - 6\tabcolsep) * \real{0.1702}}
  >{\raggedright\arraybackslash}p{(\linewidth - 6\tabcolsep) * \real{0.3546}}@{}}
\toprule\noalign{}
\begin{minipage}[b]{\linewidth}\raggedright
Package (Citation)
\end{minipage} & \begin{minipage}[b]{\linewidth}\raggedright
Vignettes
\end{minipage} & \begin{minipage}[b]{\linewidth}\raggedright
Notes
\end{minipage} & \begin{minipage}[b]{\linewidth}\raggedright
Used in
\end{minipage} \\
\midrule\noalign{}
\endhead
\bottomrule\noalign{}
\endlastfoot
naniar (Tierney \& Cook, 2023) &
\url{https://naniar.njtierney.com/index.html} & Explore missing data. &
\begin{minipage}[t]{\linewidth}\raggedright
\begin{itemize}
\tightlist
\item
  \href{data-cleaning-overview.qmd}{Data cleaning}
\end{itemize}
\end{minipage} \\
\end{longtable}

\chapter{Official R Guides}\label{official-r-guides}

Here are a collection of R cookbooks and guides that have informed my
own workflows stored in this repository.

\begin{itemize}
\item
  Learning Statistics with R (Navarro, 2026) - unbelievably helpful!
  ⭐⭐⭐
\item
  The R Cookbook (J. (JD). Long \& Teetor, 2019) - basics ⭐⭐
\item
  The R Graphics Cookbook (Chang, 2026) - basics of figures ⭐⭐
\item
  Doing Meta-Analysis in R: A Hands-On Guide (Harrer et al., 2021) -
  supremely useful for understanding three-level random-effects
  meta-analysis ⭐⭐⭐
\item
  And \href{https://github.com/qinwf/awesome-R}{here} you can find a
  list of fantastic R packages under a number of categories. For
  example, you might be interested in using R for text mining, or you
  might want to conduct Bayesian analysis.
\end{itemize}

\part{References}

\chapter*{References}\label{references-5}
\addcontentsline{toc}{chapter}{References}

\markboth{References}{References}

\phantomsection\label{refs}
\begin{CSLReferences}{1}{0}
\bibitem[\citeproctext]{ref-BenjaminiHochberg1995}
Benjamini, Y., \& Hochberg, Y. (1995). Controlling the false discovery
rate: A practical and powerful approach to multiple testing.
\emph{Journal of the Royal Statistical Society: Series B
(Methodological)}, \emph{57}(1), 289--300.
\url{https://doi.org/10.1111/j.2517-6161.1995.tb02031.x}

\bibitem[\citeproctext]{ref-Chang}
Chang, W. (2026). \emph{R graphics cookbook, 2nd edition}.
\url{https://r-graphics.org/}

\bibitem[\citeproctext]{ref-chen2007}
Chen, F. F. (2007). Sensitivity of goodness of fit indices to lack of
measurement invariance. \emph{Structural Equation Modeling: A
Multidisciplinary Journal}, \emph{14}(3), 464--504.
\url{https://doi.org/10.1080/10705510701301834}

\bibitem[\citeproctext]{ref-cohen2013applied}
Cohen, J., Cohen, P., West, S. G., \& Aiken, L. S. (2013). \emph{Applied
multiple regression/correlation analysis for the behavioral sciences}.
Routledge.

\bibitem[\citeproctext]{ref-Cooper2024}
Cooper, G. (2024). \emph{CoopG/tableone2flextable}.
\url{https://github.com/CoopG/tableone2flextable}

\bibitem[\citeproctext]{ref-eisinga2013}
Eisinga, R., Grotenhuis, M. te, \& Pelzer, B. (2013). The reliability of
a two-item scale: Pearson, Cronbach, or Spearman-Brown?
\emph{International Journal of Public Health}, \emph{58}(4), 637--642.
\url{https://doi.org/10.1007/s00038-012-0416-3}

\bibitem[\citeproctext]{ref-EndersBandalos2001}
Enders, C. K., \& Bandalos, D. L. (2001). The relative performance of
full information maximum likelihood estimation for missing data in
structural equation models. \emph{Structural Equation Modeling: A
Multidisciplinary Journal}, \emph{8}(3), 430--457.
\url{https://doi.org/10.1207/S15328007SEM0803_5}

\bibitem[\citeproctext]{ref-ferr2025}
Ferr, H. (2025). The normal distribution is not normal in psychological
data: Moving beyond parametric dogma. \emph{PLOS Mental Health},
\emph{2}(8), e0000403.
\url{https://doi.org/10.1371/journal.pmen.0000403}

\bibitem[\citeproctext]{ref-officer}
Gohel, D., Moog, S., \& Heckmann, M. (2026). \emph{Officer: Manipulation
of microsoft word and PowerPoint documents}.
\url{https://ardata-fr.github.io/officeverse/}

\bibitem[\citeproctext]{ref-flextable}
Gohel, D., \& Skintzos, P. (2026). \emph{Flextable: Functions for
tabular reporting}. \url{https://github.com/davidgohel/flextable}

\bibitem[\citeproctext]{ref-hair2019}
Hair, J. F., Risher, J. J., Sarstedt, M., \& Ringle, C. M. (2019). When
to use and how to report the results of PLS-SEM. \emph{European Business
Review}, \emph{31}(1), 2--24.
\url{https://doi.org/10.1108/EBR-11-2018-0203}

\bibitem[\citeproctext]{ref-HarrerCuijpersFurukawa}
Harrer, M., Cuijpers, P., Furukawa, T. A., \& Ebert, D. D. (2021).
\emph{Doing Meta-Analysis in R: A Hands-On Guide}.
\url{https://doing-meta.guide/}

\bibitem[\citeproctext]{ref-hu1999}
Hu, L., \& Bentler, P. M. (1999). Cutoff criteria for fit indexes in
covariance structure analysis: Conventional criteria versus new
alternatives. \emph{Structural Equation Modeling: A Multidisciplinary
Journal}, \emph{6}(1), 1--55.
\url{https://doi.org/10.1080/10705519909540118}

\bibitem[\citeproctext]{ref-fastdummies2026}
Kaplan, J. (2026). \emph{fastDummies: Fast creation of dummy (binary)
columns and rows from categorical variables}.
\url{https://github.com/jacobkap/fastDummies}

\bibitem[\citeproctext]{ref-kline2015}
Kline, R. B. (2015). \emph{Principles and practice of structural
equation modeling}. Guilford Publications.

\bibitem[\citeproctext]{ref-kugler2018coding}
Kugler, K. C., Dziak, J. J., \& Trail, J. (2018). Coding and
interpretation of effects in analysis of data from a factorial
experiment. In \emph{Optimization of behavioral, biobehavioral, and
biomedical interventions: Advanced topics} (pp. 175--205). Springer.

\bibitem[\citeproctext]{ref-little1988a}
Little, R. J. A. (1988). A Test of Missing Completely at Random for
Multivariate Data with Missing Values. \emph{Journal of the American
Statistical Association}, \emph{83}(404), 1198--1202.
\url{https://doi.org/10.1080/01621459.1988.10478722}

\bibitem[\citeproctext]{ref-jtools}
Long, J. A. (2024). Jtools: Analysis and presentation of social
scientific data. \emph{Journal of Open Source Software}, \emph{9}(101),
6610. \url{https://doi.org/10.21105/joss.06610}

\bibitem[\citeproctext]{ref-LongTeetor}
Long, J. (JD)., \& Teetor, P. (2019). \emph{R cookbook, 2nd edition}.
\url{https://rc2e.com/}

\bibitem[\citeproctext]{ref-Navarro}
Navarro, D. (2026). \emph{Learning Statistics with R}.
\url{https://learningstatisticswithr.com/}

\bibitem[\citeproctext]{ref-Perneger1998}
Perneger, T. V. (1998). What's wrong with Bonferroni adjustments.
\emph{BMJ}, \emph{316}(7139), 1236--1238.
\url{https://doi.org/10.1136/bmj.316.7139.1236}

\bibitem[\citeproctext]{ref-PittawayFieldingLouis2024}
Pittaway, C. R., Fielding, K. S., \& Louis, W. R. (2024). Pathways to
conventional and radical climate action: The role of temporal
orientation, environmental cognitive alternatives, and eco-anxiety.
\emph{Global Environmental Change}, \emph{87}, 102886.
\url{https://doi.org/10.1016/j.gloenvcha.2024.102886}

\bibitem[\citeproctext]{ref-Pittaway_Louis_Fielding2025}
Pittaway, C. R., Louis, W., \& Fielding, K. (2025). Temporal
orientations and climate action over time: A two-wave study of the
associations between consideration of consequences and climate action
intentions and behaviour. \emph{Journal of Environmental Psychology},
\emph{108}, 102833. \url{https://doi.org/10.1016/j.jenvp.2025.102833}

\bibitem[\citeproctext]{ref-splithalfr2025}
Pronk, T. (2025). \emph{Splithalfr: Estimates split-half reliabilities
for scoring algorithms of cognitive tasks and questionnaires (version
3.0.0) {[}computer software{]}}.
\url{https://doi.org/10.5281/zenodo.7777894}

\bibitem[\citeproctext]{ref-razali2011}
Razali, N. M., \& Wah, Y. B. (2011). Power comparisons of shapiro-wilk,
kolmogorov-smirnov, lilliefors and anderson-darling tests. \emph{Journal
of Statistical Modeling and Analytics}, \emph{2}(1), 21--33.

\bibitem[\citeproctext]{ref-lavaanarticle}
Rosseel, Y. (2012). {lavaan}: An {R} package for structural equation
modeling. \emph{Journal of Statistical Software}, \emph{48}(2), 1--36.
\url{https://doi.org/10.18637/jss.v048.i02}

\bibitem[\citeproctext]{ref-lavaanpackage}
Rosseel, Y., Jorgensen, T. D., \& De Wilde, L. (2026). \emph{{lavaan}:
Latent variable analysis}.
\url{https://doi.org/10.32614/CRAN.package.lavaan}

\bibitem[\citeproctext]{ref-gtsummary}
Sjoberg, D. D., Whiting, K., Curry, M., Lavery, J. A., \& Larmarange, J.
(2021). Reproducible summary tables with the gtsummary package.
\emph{{The R Journal}}, \emph{13}, 570--580.
\url{https://doi.org/10.32614/RJ-2021-053}

\bibitem[\citeproctext]{ref-apaTables2026}
Stanley, D. (2026). \emph{apaTables: Create american psychological
association (APA) style tables}.
\url{https://github.com/dstanley4/apaTables}

\bibitem[\citeproctext]{ref-rempsyc}
Thériault, R. (2023). {rempsyc}: Convenience functions for psychology.
\emph{Journal of Open Source Software}, \emph{8}(87), 5466.
\url{https://doi.org/10.21105/joss.05466}

\bibitem[\citeproctext]{ref-naniar}
Tierney, N., \& Cook, D. (2023). Expanding tidy data principles to
facilitate missing data exploration, visualization and assessment of
imputations. \emph{Journal of Statistical Software}, \emph{105}(7),
1--31. \url{https://doi.org/10.18637/jss.v105.i07}

\bibitem[\citeproctext]{ref-tidyverse2019}
Wickham, H., Averick, M., Bryan, J., Chang, W., McGowan, L. D.,
François, R., Grolemund, G., Hayes, A., Henry, L., Hester, J., Kuhn, M.,
Pedersen, T. L., Miller, E., Bache, S. M., Müller, K., Ooms, J.,
Robinson, D., Seidel, D. P., Spinu, V., \ldots{} Yutani, H. (2019).
Welcome to the {tidyverse}. \emph{Journal of Open Source Software},
\emph{4}(43), 1686. \url{https://doi.org/10.21105/joss.01686}

\bibitem[\citeproctext]{ref-readr2026}
Wickham, H., Hester, J., \& Bryan, J. (2026). \emph{Readr: Read
rectangular text data}. \url{https://readr.tidyverse.org}

\bibitem[\citeproctext]{ref-haven2025}
Wickham, H., Miller, E., \& Smith, D. (2025). \emph{Haven: Import and
export 'SPSS', 'stata' and 'SAS' files}.
\url{https://haven.tidyverse.org}

\bibitem[\citeproctext]{ref-psych2026}
William Revelle. (2026). \emph{Psych: Procedures for psychological,
psychometric, and personality research}. Northwestern University.
\url{https://CRAN.R-project.org/package=psych}

\bibitem[\citeproctext]{ref-yazici2007comparison}
Yazici, B., \& Yolacan, S. (2007). A comparison of various tests of
normality. \emph{Journal of Statistical Computation and Simulation},
\emph{77}(2), 175--183.

\bibitem[\citeproctext]{ref-tableone}
Yoshida, K., \& Bartel, A. (2022). \emph{Tableone: Create 'table 1' to
describe baseline characteristics with or without propensity score
weights}. \url{https://CRAN.R-project.org/package=tableone}

\end{CSLReferences}




\end{document}
